\documentclass[11pt]{article}

% ============================================================
% Packages
% ============================================================
\usepackage[a4paper,margin=1in]{geometry}
\usepackage{amsmath,amssymb,amsthm}
\usepackage{mathtools}
\usepackage{bm}
\usepackage{physics}   % optional; remove if you prefer explicit notation
\usepackage{hyperref}
\usepackage{graphicx}
\usepackage{booktabs}
\usepackage{enumitem}
\usepackage{tabularx}

\hypersetup{
  colorlinks=true,
  linkcolor=blue,
  citecolor=blue,
  urlcolor=blue
}

% ============================================================
% Macros / Notation
% ============================================================
\newcommand{\ii}{\mathrm{i}}
\newcommand{\ee}{\mathrm{e}}
%\newcommand{\dd}{\mathrm{d}}
\newcommand{\R}{\mathbb{R}}

% Coordinates and index sets
\newcommand{\X}{\mathbf{X}}     % bulk spatial (x,y,z,w)
\newcommand{\x}{\mathbf{x}}     % brane spatial (x,y,z)
\newcommand{\gradfour}{\nabla_4}
\newcommand{\gradthree}{\nabla_3}
\newcommand{\lapfour}{\nabla_4^2}
\newcommand{\lapthree}{\nabla_3^2}

% Bulk index conventions (use in text):
%   M,N,... in {0,1,2,3,4}, with 4 ≡ w
%   μ,ν,... in {0,1,2,3}
%   i,j,k,... in {1,2,3}

% Fields
\newcommand{\A}{A}              % gauge potential A_M
\newcommand{\F}{F}              % field strength F_{MN}
\newcommand{\J}{J}              % total current J^M
\newcommand{\Js}{J_s}           % species current J_s^M
\newcommand{\Z}{Z}              % EM localization weight Z(w)
\newcommand{\W}{W}              % brane projection weight W(w)

% Extra-dimension coordinate and operators
\newcommand{\w}{w}
\newcommand{\delw}{\partial_w}

% Brane-projected observable operator (define operationally in the text)
\newcommand{\Proj}[1]{\overline{#1}} % e.g. \Proj{Q} := \int dw W(w) Q

% Plasma species labels
\newcommand{\ion}{\mathrm{i}}
\newcommand{\ele}{\mathrm{e}}

% EM fields (3D brane components)
\newcommand{\E}{\mathbf{E}}
\newcommand{\B}{\mathbf{B}}
%\newcommand{\v}{\mathbf{v}}
\newcommand{\Jb}{\mathbf{J}}
\newcommand{\rhob}{\rho}

% Extra / mixed fields
\newcommand{\Ew}{E_w}           % F_{0w}
\newcommand{\Bw}{\mathbf{B}_w}  % shorthand for mixed components (define precisely)

% Leakage source in projected continuity
\newcommand{\Sleak}{S_{\rm leak}}

% Effective coupling after reduction
\newcommand{\muzero}{\mu_0}
\newcommand{\muzeff}{\mu_0^{\rm eff}}
\newcommand{\Zint}{Z_{\rm int}}

% KK / mode labels
\newcommand{\mkk}{m_n}
\newcommand{\phin}{\phi_n}

% Theorem-like environments (optional)
\newtheorem{claim}{Claim}
\newtheorem{definition}{Definition}

% ============================================================
% Title
% ============================================================
\title{Plasma Dynamics from the Unified 4D Model:\\
A 4+1D Alternative to MHD with a Controlled MHD Limit and 4D Interaction Corrections}
\author{Trevor Norris}
\date{\today}

\begin{document}
\maketitle

\begin{abstract}
We present a 4+1D plasma interaction model derived from a unified bulk action with localized electromagnetism and an operational brane-observer map. The model provides exact bulk conservation laws, exact projected (brane) balance laws with explicit leakage terms, and controlled reductions that recover standard 3+1 Maxwell theory, two-fluid plasma equations, and the ideal-MHD induction equation. We then identify and organize the leading beyond-MHD corrections that arise when 4D interactions are retained: a dynamical scalar-photon sector \(A_w\), brane--bulk charge/current exchange \(J^w\), and finite-localization (KK/mode) corrections that modify propagation and topology transport. Finally, we outline a simulation architecture for massive plasma computations that evolves either the full 4+1D system or a mode-reduced Hermite basis in \(w\), with diagnostics designed to quantify reconnection and topology change as genuine 4D transport rather than ad hoc resistivity.
\end{abstract}

% ============================================================
\section{Introduction and Motivation}
\label{sec:intro}

Magnetohydrodynamics (MHD) is the dominant continuum framework for macroscopic plasma dynamics across laboratory, space, and astrophysical settings. In its standard regime of validity---quasi-neutrality, low-frequency dynamics compared to light and plasma frequencies, and lengthscales large compared to kinetic scales---MHD provides a compact and remarkably predictive description. However, the most important (and often most computationally expensive) plasma phenomena frequently occur precisely where ideal MHD is structurally stressed: strong gradients, topological reorganization of field structure, and rapid energy conversion. The canonical example is magnetic reconnection, where ideal flux-freezing would prohibit topology change, yet observed plasmas change connectivity on dynamical timescales.

A large portion of ``extended MHD'' and kinetic plasma theory can be understood as principled ways to weaken ideal constraints by introducing additional degrees of freedom (Hall drift, electron pressure tensor, electron inertia, collisions, anomalous transport, etc.). In practice, many large-scale simulations still rely on closure choices whose physical meaning is indirect: reconnection is triggered by an imposed resistivity, by numerical diffusion, or by semi-empirical terms that are difficult to relate to a deeper ontology.

This paper develops an alternative starting point. Instead of taking 3D MHD as the fundamental description, we treat the observable 3D world as an \emph{operational brane} that measures a projection of a conservative \emph{bulk} dynamics living in four spatial dimensions \((x,y,z,w)\). The key claim is not that MHD is incorrect, but that it is a controlled \emph{limit} of a more complete interaction model. In that parent theory, topological evolution that appears non-ideal from a brane perspective can arise from conservative transport into the extra spatial direction \(w\) and into transverse field structure (mode towers), rather than from phenomenological dissipation imposed by hand.

Concretely, we study a coupled 4+1D system consisting of:
(i) a localized electromagnetic gauge field \(A_M(t,\x,w)\) whose kinetic term is weighted by a localization function \(Z(w)\), and
(ii) a multi-species charged medium that may be represented either as a multi-field nonrelativistic matter sector (Schr\"odinger/GNLS-type) or as a two-fluid/kinetic plasma sector.
The observable brane fields are defined by a separate projection kernel \(W(w)\), emphasizing that ``the brane'' is not a boundary condition but a measurement map. This separation (bulk dynamics vs.\ brane observables) is essential: the bulk can remain conservative while the projected brane system becomes an open subsystem with explicit exchange terms.

\subsection{What This Paper Claims (and What It Does Not)}
\label{sec:claims_scope}

\paragraph{Exact identities vs.\ controlled limits.}
We distinguish three categories of statements throughout:
\begin{enumerate}[leftmargin=1.5em]
\item \textbf{Exact identities} that follow from the bulk field equations and definitions (e.g.\ Bianchi identities, bulk charge conservation, exact projected balance laws with leakage sources).
\item \textbf{Controlled limits} in which the bulk system reduces to familiar 3+1D theories under explicit assumptions (e.g.\ zero-mode dominance in \(w\), negligible transverse current \(J^w\), quasi-neutrality, and small electron inertia).
\item \textbf{Regime statements} that specify when particular correction terms are expected to matter (e.g.\ when \(w\)-gradients are excited, when the scalar-photon sector \(A_w\) becomes dynamical, or when a finite number of transverse modes fails to capture observed dynamics).
\end{enumerate}

\paragraph{Deriving MHD is a requirement, not the endpoint.}
A central goal of the paper is to show that standard MHD emerges as a \emph{recoverable} description. We provide (i) a controlled reduction from localized 4+1D electromagnetism to effective 3+1D Maxwell theory on the brane, and (ii) a controlled reduction from a two-fluid brane plasma model to the ideal-MHD induction equation. These reductions are algebraic and are additionally verified by a symbolic reduction script (not included in the paper as source code).

\paragraph{The new physics is organized as explicit correction channels.}
The second central goal is to identify the leading beyond-MHD terms that the parent 4D model predicts when the controlled assumptions are relaxed. These corrections are not introduced phenomenologically; they arise from retained degrees of freedom and from the non-commutativity of projection with nonlinear dynamics. The paper organizes the most important extensions into:
\begin{enumerate}[leftmargin=1.5em]
\item a dynamical \emph{scalar-photon} sector \(A_w\) and associated mixed field components \(F_{\mu w}\), which drive bulk drift and new force channels;
\item explicit \emph{brane--bulk exchange} via nonzero transverse current \(J^w\), yielding open-system source terms in projected brane balance laws;
\item \emph{finite-localization / transverse-mode} corrections (a KK/Hermite tower for Gaussian \(Z(w)\)) that provide computable short-range modifications and additional energy/topology storage channels;
\item a projected unresolved-stress contribution (a ``topology stress tensor'') that quantifies how sub-\(w\) structure acts as an effective non-ideal term in the brane equations without requiring dissipation into heat.
\end{enumerate}

\paragraph{Hybrid relativity (scope of validity).}
The electromagnetic sector is evolved as a relativistic hyperbolic system (wave propagation at speed \(c\)), while the baseline matter/plasma sector is nonrelativistic (as is standard in Schr\"odinger--Maxwell, two-fluid--Maxwell, and Vlasov--Maxwell simulation frameworks). The intent of this paper is to target MHD-like and extended-MHD-like regimes; fully relativistic plasma matter sectors (e.g.\ Dirac/Klein--Gordon descriptions or explicitly relativistic kinetic closures) are deferred to future work.

\paragraph{What we do not claim.}
We do not claim that the present paper completes a full kinetic microphysics program for all plasma regimes, nor that any single truncation of the transverse-mode tower is universally adequate. Instead, we provide a systematic parent model, a controlled reduction to known limits, and a hierarchy of correction terms with explicit simulation knobs that can be increased until convergence is achieved.

\paragraph{Charge ontology and scope.}
This plasma paper treats species charges as fixed input parameters and does not attempt to derive particle identity. For elementary defect branches, the electric sign is a fixed topological orientation label \(\eta_Q=\pm 1\) and the observed brane magnitude is inherited from localization thickness; in the plasma equations we therefore retain \(q_s\) simply as the fixed species charge label. Circulation and canonical vorticity remain part of the magnetic--vortical sector and are not used here to define electric charge.

\subsection{Roadmap}
This paper is organized as follows. Section~\ref{sec:conventions} defines the bulk/brane geometry, index conventions, and the operational projection map used to define brane observables. Section~\ref{sec:em_sector} presents the localized 4+1D electromagnetic sector, emphasizing the physical role of the \(A_w\) degree of freedom and the transverse-mode (KK/Hermite) structure induced by a localized kinetic weight \(Z(w)\). Section~\ref{sec:matter_sector} introduces the plasma matter sector (multi-species field, two-fluid, or kinetic) and the associated currents and conservation laws. Section~\ref{sec:full_system} collects the coupled bulk system to be simulated and summarizes bulk invariants.

Section~\ref{sec:mhd_reduction} presents the controlled reductions that recover effective 3+1 Maxwell theory on the brane and then recover the standard MHD limit from brane two-fluid dynamics. Section~\ref{sec:beyond_mhd} is the main novelty section: it derives and organizes the beyond-MHD correction channels associated with \(A_w\), \(J^w\), finite localization (mode towers), and projected unresolved stresses. Section~\ref{sec:computation} describes how the model is intended to be simulated at scale, contrasting direct 4D-grid evolution with a mode-reduced 3D\(\times\)modes scheme (Hermite in \(w\)) designed for massive computations. Section~\ref{sec:validation} outlines validation benchmarks in the controlled MHD limit and proposed tests that isolate genuinely 4D interaction effects. We conclude in Section~\ref{sec:conclusion} with a summary of the model hierarchy and near-term simulation targets.

% ============================================================
\section{Conventions, Geometry, and Operational Brane Observables}
\label{sec:conventions}

This section fixes notation and defines what is meant by ``bulk'' versus ``brane''
quantities in a way that is suitable for both analysis and simulation. The guiding
principle is that the bulk theory is posed on four spatial dimensions
\(\X=(x,y,z,w)\), while brane-facing observables are defined \emph{operationally}
by a projection (measurement) map rather than by imposing boundary matching
conditions.

\subsection{Bulk Geometry and Index Conventions}
\label{sec:bulk_geometry}

\paragraph{Coordinates.}
We work on a \((4+1)\)-dimensional spacetime with coordinates
\begin{equation}
x^M = (t, x, y, z, w) \equiv (t,\x,w),
\qquad
\X \equiv (x,y,z,w)\in\R^4,
\qquad
\x \equiv (x,y,z)\in\R^3.
\label{eq:coords_def}
\end{equation}
The operational brane is centered at \(w=0\). In the full theory \(w\) is a bona
fide spatial coordinate; ``on the brane'' will always mean ``as measured by the
projection defined in Sec.~\ref{sec:projection_defs}.''

\paragraph{Indices.}
We use the following index sets:
\begin{align}
&\text{Bulk spacetime indices:} && M,N,\ldots \in \{0,1,2,3,4\},\quad 0\equiv t,\;4\equiv w,
\nonumber\\
&\text{Brane spacetime indices:} && \mu,\nu,\ldots \in \{0,1,2,3\},
\nonumber\\
&\text{Bulk spatial indices:} && A,B,\ldots \in \{1,2,3,4\}\equiv\{x,y,z,w\},
\nonumber\\
&\text{Brane spatial indices:} && a,b,\ldots \in \{1,2,3\}\equiv\{x,y,z\}.
\label{eq:index_sets}
\end{align}
We frequently decompose bulk fields into brane and transverse components, e.g.
\begin{equation}
A_M=(A_0,A_a,A_w),
\qquad
\partial_A=(\partial_a,\partial_w).
\label{eq:field_decomposition}
\end{equation}

\paragraph{Metric and derivatives.}
For the electromagnetic sector we use a flat \((4+1)\) Minkowski metric with
mostly-plus signature,
\begin{equation}
\eta_{MN}=\mathrm{diag}(-1,+1,+1,+1,+1),
\label{eq:metric_def}
\end{equation}
used for index bookkeeping and for defining the Maxwell kinetic operator. We set
the brane light-cone speed to unity in this paper (restore \(c\) by
\(t\mapsto ct\) in \(\eta_{MN}\) and \(\partial_t\mapsto c^{-1}\partial_t\)).
Derivatives are \(\partial_M\equiv \partial/\partial x^M\). We define
\begin{equation}
\gradthree = (\partial_x,\partial_y,\partial_z),\quad
\lapthree = \partial_x^2+\partial_y^2+\partial_z^2,
\qquad
\gradfour = (\gradthree,\partial_w),\quad
\lapfour = \lapthree + \partial_w^2.
\label{eq:diff_ops}
\end{equation}
The \((4+1)\) d'Alembertian is
\begin{equation}
\square_5 \equiv \eta^{MN}\partial_M\partial_N
= -\partial_t^2 + \lapthree + \partial_w^2.
\label{eq:box5_def}
\end{equation}

\paragraph{Integration measures and boundary terms.}
We write
\begin{equation}
\int \dd^4X \equiv \int_{\R^4}\dd x\,\dd y\,\dd z\,\dd w,
\qquad
\int \dd^3x \equiv \int_{\R^3}\dd x\,\dd y\,\dd z.
\label{eq:measures}
\end{equation}
When integrating by parts in \(w\), we will either (i) retain the boundary flux
terms explicitly, or (ii) state sufficient decay/localization assumptions for
those terms to vanish as \(|w|\to\infty\).

\paragraph{Hybrid relativity (model symmetry).}
The electromagnetic sector is evolved as a relativistic hyperbolic field theory
in \((4+1)\) dimensions. The plasma matter sector used in this paper is
nonrelativistic (two-fluid or Schr\"odinger/GNLS type). Thus the combined system is
not assumed to be fully \((4+1)\)-Lorentz invariant; ``covariant'' below refers to
\emph{gauge covariance} of matter couplings and to the covariant form of the Maxwell
sector in flat \((4+1)\) spacetime.

\subsection{Localization and Projection Operators}
\label{sec:projection_defs}

A central distinction in this framework is between:
\begin{itemize}[leftmargin=1.7em]
\item \textbf{Localization:} a weight \(\Z(w)\) that appears in the \emph{bulk action}
and localizes gauge dynamics near \(w=0\).
\item \textbf{Projection:} a (generally independent) measurement kernel \(\W(w)\)
that defines \emph{brane observables} as weighted averages over \(w\).
\end{itemize}

\paragraph{Localization profile \texorpdfstring{$\Z(w)$}{Z(w)}.}
The electromagnetic kinetic term is multiplied by a nonnegative localization profile
\(\Z(w)\) with finite integral
\begin{equation}
\Zint \equiv \int_{-\infty}^{+\infty}\dd w\;\Z(w) < \infty.
\label{eq:Zint_def}
\end{equation}
A canonical analytic choice is a Gaussian profile,
\begin{equation}
\Z(w)=\exp\!\left(-\frac{w^2}{\lambda^2}\right),
\qquad
\Zint=\lambda\sqrt{\pi},
\label{eq:Z_gaussian_def}
\end{equation}
but none of the definitions in this section require a specific form.

\paragraph{Operational brane projection \texorpdfstring{$\W(w)$}{W(w)}.}
We define brane observables by a normalized kernel \(\W(w)\),
\begin{equation}
\int_{-\infty}^{+\infty}\dd w\;\W(w)=1,
\qquad
\W(w)\ge 0,
\qquad
\W(w)\ \text{peaked near}\ w=0.
\label{eq:W_normalized}
\end{equation}
For any bulk field \(Q(t,\x,w)\), the brane-measured quantity is
\begin{equation}
\Proj{Q}(t,\x) \equiv \int_{-\infty}^{+\infty}\dd w\;\W(w)\,Q(t,\x,w).
\label{eq:brane_projection}
\end{equation}
In this normalized convention, \(\Proj{Q}\) has the same units as \(Q\). If one
instead wants a strictly dimension-reduced density per brane volume (integrated
through the transverse direction), one may use an unnormalized kernel; we adopt
\eqref{eq:W_normalized} because it corresponds to an operational ``finite
resolution'' measurement across \(w\) and makes leakage bookkeeping explicit.

\paragraph{Matched kernels (optional).}
A particularly natural operational choice is the matched kernel
\begin{equation}
\W(w)=\frac{\Z(w)}{\Zint},
\label{eq:W_matched}
\end{equation}
which weights measurements in the same region where the gauge dynamics is
localized. We keep \(\W\) independent in the formalism to separate ``how the bulk
evolves'' (set by \(\Z\)) from ``what a brane observer measures'' (set by \(\W\)).

\subsection{Fields and Brane-Observable Decompositions}
\label{sec:field_decompositions}

\paragraph{Gauge field and field strength.}
The bulk gauge potential is \(A_M(t,\x,w)\), with field strength
\begin{equation}
F_{MN} \equiv \partial_M A_N - \partial_N A_M.
\label{eq:F_def_conv}
\end{equation}
We frequently separate brane and transverse components:
\begin{equation}
F_{\mu\nu},\qquad F_{\mu w}.
\label{eq:F_split}
\end{equation}
Gauge transformations act as \(A_M\mapsto A_M+\partial_M\chi\), leaving \(F_{MN}\)
invariant.

\paragraph{Brane-facing electric and magnetic fields.}
For later comparison with standard plasma notation, we define brane-facing fields
from the \(\W\)-projected field strength:
\begin{align}
E_a^{\rm (brane)}(t,\x)
&\equiv \Proj{F_{a0}}(t,\x),
\label{eq:E_brane_def}
\\
B_a^{\rm (brane)}(t,\x)
&\equiv \frac12\,\epsilon_{abc}\,\Proj{F_{bc}}(t,\x),
\label{eq:B_brane_def}
\end{align}
where \(\epsilon_{abc}\) is the Levi-Civita symbol on the brane with
\(\epsilon_{xyz}=+1\). The bulk theory also contains physical mixed components
that have no analogue in strictly \(3+1\) Maxwell theory:
\begin{equation}
E_w(t,\x,w)\equiv F_{w0},
\qquad
C_a(t,\x,w)\equiv F_{aw}.
\label{eq:mixed_fields_def}
\end{equation}
Their brane-observed counterparts are \(\Proj{E_w}\) and \(\Proj{C_a}\). In the
controlled MHD/Maxwell limit these mixed components are assumed negligible; in
the full 4D interaction theory they provide additional force and topology-transport
channels.

\paragraph{Plasma observables (multi-species).}
For each plasma species \(s\) (e.g.\ electrons and ions), let \(\rho_s(t,\x,w)\)
denote the bulk number density and \(j_s^A(t,\x,w)\) the bulk number current in
the four spatial directions. Brane-observed densities and brane spatial currents
are defined by
\begin{equation}
\rho_{s,\rm brane}(t,\x)\equiv \Proj{\rho_s}(t,\x),
\qquad
j_{s,\rm brane}^a(t,\x)\equiv \Proj{j_s^a}(t,\x).
\label{eq:brane_species_defs}
\end{equation}
When \(\rho_{s,\rm brane}>0\), we define a brane-observed species velocity by
\begin{equation}
v_{s,\rm brane}^a \equiv \frac{j_{s,\rm brane}^a}{\rho_{s,\rm brane}}.
\label{eq:brane_velocity_def}
\end{equation}
This is a diagnostic/observable definition; it does \emph{not} imply that the bulk
flow is strictly confined to the brane.

\subsection{Exact Projected Balance Laws and Leakage Terms}
\label{sec:leakage_identity}

The projection \eqref{eq:brane_projection} does not, in general, commute with bulk
conservation laws. This is the mathematical origin of brane ``non-ideality'' in a
conservative bulk theory.

\paragraph{Exact leakage identity (species continuity).}
Assume each species satisfies a bulk continuity equation
\begin{equation}
\partial_t\rho_s + \partial_a j_s^a + \partial_w j_s^w = 0.
\label{eq:bulk_continuity_species_leakage}
\end{equation}
Multiply by \(\W(w)\) and integrate over \(w\). Using integration by parts in the
\(\partial_w j_s^w\) term yields the exact projected balance law
\begin{equation}
\partial_t \rho_{s,\rm brane} + \partial_a j_{s,\rm brane}^a
=
-\Big[\W(w)\,j_s^w\Big]_{w=-\infty}^{w=+\infty}
+\int_{-\infty}^{+\infty}\dd w\;\W'(w)\,j_s^w(t,\x,w).
\label{eq:projected_continuity_leakage}
\end{equation}
We define the right-hand side as the leakage/source term
\begin{equation}
\Sleak^{(s)}(t,\x)
\equiv
-\Big[\W\,j_s^w\Big]_{-\infty}^{+\infty}
+\int_{-\infty}^{+\infty}\dd w\;\W'(w)\,j_s^w.
\label{eq:Sleak_def}
\end{equation}
If \(\W(w)\) is compactly supported or decays sufficiently fast and \(j_s^w\) does
not grow at infinity, the boundary flux term vanishes and the leakage is carried
entirely by the overlap with \(\W'(w)\). In either case, \eqref{eq:projected_continuity_leakage}
is an identity: brane densities can change due to transverse transport even when
the bulk evolution is exactly conservative.

\paragraph{Interpretation.}
In the controlled MHD limit, one assumes \(\Sleak^{(s)}\approx 0\) (no net exchange
through \(w\) at the measurement scale), and \eqref{eq:projected_continuity_leakage}
reduces to ordinary 3D continuity on the brane. In the full 4D interaction model,
\(\Sleak^{(s)}\) is retained and becomes a physically meaningful diagnostic:
reconnection-like and topology-changing behavior in the brane-facing fields can
be reinterpreted as conservative transport into transverse degrees of freedom
(\(j^w\), mixed EM components \(F_{\mu w}\), and/or mode towers).

\paragraph{Projection versus reduction (summary).}
We close this section by emphasizing the separation:
\begin{itemize}[leftmargin=1.7em]
\item \textbf{Reduction} uses \(\Zint=\int \Z(w)\dd w\) to obtain effective brane
couplings (e.g.\ \(\mu_0^{\rm eff}\)) under a stated ansatz (e.g.\ a brane-dominant
zero mode).
\item \textbf{Projection} uses \(\W(w)\) to define what is observed on the brane
and yields exact open-system identities such as \eqref{eq:projected_continuity_leakage}.
\end{itemize}
This distinction is essential for organizing ``MHD as a controlled limit'' and
``beyond-MHD physics as controlled relaxations'' within a single bulk theory.

% ============================================================
\section{4+1D Electromagnetic Sector with Localization}
\label{sec:em_sector}

This section defines the electromagnetic (gauge) sector of the model in \(4+1\)
dimensions with a transverse localization profile \(\Z(w)\). The key features are:
(i) the bulk equations are those of Maxwell theory with a weighted kinetic term,
(ii) brane-facing Maxwell theory is recovered as a \emph{controlled} reduction with
an effective coupling \(\mu_{0,\mathrm{eff}}\), and (iii) for Gaussian localization
the transverse operator admits a closed-form Hermite (KK) tower whose exchange
produces fixed-coefficient Yukawa corrections. These structures provide both a
clean derivation of the MHD limit (later) and a computable hierarchy of beyond-MHD
degrees of freedom (the \(A_w\) sector and transverse-mode excitations).

\subsection{Localized Maxwell action and field equations}
\label{sec:localized_maxwell}

\paragraph{Action.}
The bulk gauge potential is \(\A_M(t,\x,w)\) with field strength
\begin{equation}
  \F_{MN} \equiv \partial_M \A_N - \partial_N \A_M,
  \qquad
  \F^{MN} \equiv \eta^{MP}\eta^{NQ}\F_{PQ},
\end{equation}
where \(\eta_{MN}=\mathrm{diag}(-1,+1,+1,+1,+1)\) is the flat bulk metric from
Sec.~\ref{sec:conventions}. The localized Maxwell action is
\begin{equation}
S_{\rm EM}[\A;\Z]
= \int \dd t\,\dd^3x\,\dd w\;
\left[
-\frac{\Z(w)}{4\muzero}\,\F_{MN}\F^{MN}
-\frac{1}{2\xi\muzero}\,(\partial\!\cdot\!\A)^2
- \A_M \J^M
\right].
\label{eq:action_em}
\end{equation}
The gauge-fixing term (parameter \(\xi\)) is included only to make the kinetic
operator invertible; all gauge-invariant observables are \(\xi\)-independent for
conserved currents. The interaction is written \(\mathcal{L}_{\rm int}=-\A_M\J^M\),
so that under a gauge variation \(\delta \A_M=\partial_M\chi\),
\(\delta\mathcal{L}_{\rm int}=-\partial_M(\chi\J^M)+\chi\,\partial_M\J^M\); thus
gauge invariance of the coupling is equivalent to \(\partial_M\J^M=0\) up to
boundary terms.

\paragraph{Euler--Lagrange equations.}
Varying \eqref{eq:action_em} yields the localized Maxwell equations
\begin{equation}
\partial_M\!\left(\Z(w)\F^{MN}\right)
+\frac{1}{\xi}\,\partial^N(\partial\!\cdot\!\A)
=\muzero\,\J^N.
\label{eq:maxwell_localized}
\end{equation}
These reduce to the usual Maxwell equations when \(\Z\equiv 1\). Here \(\Z(w)\)
breaks full \(4+1\) translation/Lorentz symmetry in the \(w\)-direction while
preserving brane Lorentz covariance in the \((t,\x)\) directions.

\paragraph{Bianchi identities (homogeneous Maxwell).}
By definition of \(\F_{MN}\),
\begin{equation}
\partial_{[L}\F_{MN]}=0,
\label{eq:bianchi}
\end{equation}
an exact off-shell identity. The brane-facing homogeneous equations follow upon
projection to brane observables (Sec.~\ref{sec:projection_defs}).

\paragraph{Divergence identity and current conservation.}
Taking \(\partial_N\) of \eqref{eq:maxwell_localized} and using antisymmetry of
\(\F^{MN}\) gives
\begin{equation}
\frac{1}{\xi}\,\square_5(\partial\!\cdot\!\A)
=\muzero\,\partial_N\J^N,
\qquad
\square_5 \equiv \eta^{MN}\partial_M\partial_N
= -\partial_t^2+\lapthree+\partial_w^2.
\label{eq:div_identity}
\end{equation}
In Lorenz gauge \((\partial\!\cdot\!\A=0)\), the field equations enforce bulk
current conservation \(\partial_N\J^N=0\). Conversely, for conserved sources the
gauge-fixing term simply selects a representative in the gauge orbit without
affecting gauge-invariant quantities.

\subsection{3+1 decomposition and the \texorpdfstring{$A_w$}{Aw} degree of freedom}
\label{sec:Aw_scalar_photon}

We decompose the gauge potential and current into brane and transverse components,
\begin{equation}
\A_M=(A_0,A_a,A_w),
\qquad
\J^M=(J^0,J^a,J^w),
\qquad a\in\{x,y,z\}.
\end{equation}
In addition to the brane-facing electric/magnetic fields, the \(4+1\) theory
contains mixed field components with no analogue in strictly \(3+1\) Maxwell
theory:
\begin{equation}
E_w \equiv \F_{w0} = -\partial_t A_w - \partial_w A_0,
\qquad
C_a \equiv \F_{aw} = \partial_a A_w - \partial_w A_a.
\label{eq:mixed_fields_em}
\end{equation}
These enter forces on charged matter through the full Lorentz tensor structure
\(q\,\F_{iN}u^N\) (later sections). In particular, a dynamical \(A_w\) generically
drives transverse drift and induces \(J^w\neq 0\), providing a natural,
conservative channel for brane-observed non-ideality (leakage and topology export)
even when the bulk evolution is otherwise ``ideal.''

\paragraph{Axial gauge as a convenience (not a physical assumption).}
A common simplification is axial gauge \(A_w=0\). Since under a gauge
transformation \(\A_w\mapsto \A_w+\partial_w\chi\), axial gauge can be reached
locally by choosing
\begin{equation}
\chi_{\rm Ax}(t,\x,w)\equiv -\int_{0}^{w}A_w(t,\x,w')\,\dd w',
\end{equation}
provided \(A_w\) is sufficiently regular in \(w\). The residual gauge freedom then
satisfies \(\partial_w\chi=0\), i.e. \(\chi=\chi_b(t,\x)\), which is the usual
\(3+1\) gauge freedom. In this plasma paper we keep the full \(A_w\) sector
available; axial gauge will be used only when deriving controlled brane limits or
when defining the transverse spectral basis.

\subsection{Controlled brane Maxwell limit and effective coupling}
\label{sec:controlled_brane_maxwell}

The localization profile is assumed nonnegative and integrable,
\begin{equation}
\Zint \equiv \int_{-\infty}^{+\infty}\Z(w)\,\dd w < \infty.
\label{eq:Zint_def_em}
\end{equation}
A canonical analytic choice (used for explicit spectra below) is Gaussian
localization,
\begin{equation}
\Z(w)=\exp\!\left(-\frac{w^2}{\lambda^2}\right),
\qquad
\Zint=\lambda\sqrt{\pi}.
\label{eq:Z_gaussian_em}
\end{equation}

\paragraph{Exact \(w\)-integrated identity.}
Consider the \(\nu\)-components (\(\nu\in\{0,1,2,3\}\)) of
\eqref{eq:maxwell_localized} and integrate over \(w\):
\begin{align}
\partial_\mu\!\left(\int_{-\infty}^{+\infty}\Z\,\F^{\mu\nu}\,\dd w\right)
\;+\;
\Big[\Z\,\F^{w\nu}\Big]_{w=-\infty}^{w=+\infty}
\;+\;
\frac{1}{\xi}\,\partial^\nu\!\left(\int_{-\infty}^{+\infty}(\partial\!\cdot\!\A)\,\dd w\right)
&=\muzero\,\J_{\rm eff}^\nu,
\label{eq:integrated_exact_em}
\\
\J_{\rm eff}^\nu(t,\x)
&\equiv \int_{-\infty}^{+\infty} \J^\nu(t,\x,w)\,\dd w.
\end{align}
For localized field configurations the boundary flux term vanishes. It is then
natural to define the \(\Z\)-weighted brane field strength
\begin{equation}
\langle \F^{\mu\nu}\rangle_{\Z}
\;\equiv\;
\frac{1}{\Zint}\int_{-\infty}^{+\infty}\Z(w)\,\F^{\mu\nu}(t,\x,w)\,\dd w.
\label{eq:F_Z_weighted_def}
\end{equation}
Dividing \eqref{eq:integrated_exact_em} by \(\Zint\) gives the exact integrated
brane equation
\begin{equation}
\partial_\mu\langle \F^{\mu\nu}\rangle_{\Z}
+\frac{1}{\xi\,\Zint}\,\partial^\nu\!\left(\int(\partial\!\cdot\!\A)\,\dd w\right)
=\frac{\muzero}{\Zint}\,\J_{\rm eff}^\nu.
\label{eq:brane_maxwell_exact_weighted_em}
\end{equation}
This is \emph{not yet} standard Maxwell, because it still contains the explicit
\(w\)-integral of the gauge-fixing term and the \(\Z\)-weighted average of the bulk
field strength.

\paragraph{Controlled zero-mode reduction.}
To recover ordinary \(3+1\) Maxwell theory we now take a controlled limit:
\begin{enumerate}[leftmargin=1.6em]
\item \textbf{Zero-mode dominance:} the brane components are approximately
\(w\)-independent in the localization region,
\(\A_\mu(t,\x,w)\simeq a_\mu(t,\x)\), \(\partial_w\A_\mu\simeq 0\).
\item \textbf{Negligible mixed flux:} \(\F^{w\nu}\) is negligible and/or sufficiently
localized so that \(\big[\Z\F^{w\nu}\big]_{-\infty}^{+\infty}=0\) and
\(\langle \F^{\mu\nu}\rangle_{\Z}\simeq f^{\mu\nu}\) with
\(f_{\mu\nu}\equiv \partial_\mu a_\nu-\partial_\nu a_\mu\).
\item \textbf{No net transverse current in this limit:} the \(N=w\) equation is
consistent with \(J^w\simeq 0\) (see below).
\end{enumerate}
Under these assumptions, \eqref{eq:brane_maxwell_exact_weighted_em} reduces to
\begin{equation}
\partial_\mu f^{\mu\nu}
+\frac{1}{\xi}\,\partial^\nu(\partial_\mu a^\mu)
=\muzeff\,\J_{\rm eff}^\nu,
\qquad
\muzeff \equiv \frac{\muzero}{\Zint}.
\label{eq:brane_maxwell_gf_em}
\end{equation}
If the effective brane current is brane-localized (e.g.\ \(J^\mu=J_b^\mu\delta(w)\))
then \(\J_{\rm eff}^\nu=\J_b^\nu\). For Gaussian \(\Z\),
\begin{equation}
\muzeff=\frac{\muzero}{\lambda\sqrt{\pi}}.
\label{eq:mu0_eff_em}
\end{equation}
In Lorenz gauge on the brane \((\partial_\mu a^\mu=0)\) this becomes the standard
inhomogeneous Maxwell equation \(\partial_\mu f^{\mu\nu}=\muzeff\,\J_b^\nu\).

\paragraph{Consistency of the transverse component.}
The \(N=w\) component of \eqref{eq:maxwell_localized} is
\begin{equation}
\partial_M\!\left(\Z\,\F^{Mw}\right)+\frac{1}{\xi}\,\partial^w(\partial\!\cdot\!\A)=\muzero\,\J^w.
\label{eq:eom_w_em}
\end{equation}
In axial gauge \(A_w=0\) with the zero-mode ansatz \(\partial_w A_\mu\simeq 0\),
one has \(\F^{\mu w}=-\partial^w A^\mu\simeq 0\) and \(\partial^w(\partial\!\cdot\!A)\simeq 0\),
so \eqref{eq:eom_w_em} reduces to \(\J^w\simeq 0\). Thus the controlled brane Maxwell
limit is internally consistent precisely when transverse current is negligible at
the reduction scale. In the full plasma system we \emph{do not} impose \(\J^w=0\);
instead, \(\J^w\) becomes a dynamical diagnostic and a source of beyond-MHD terms.

\subsection{Mode expansion in the extra dimension}
\label{sec:kk_modes_em}

When \(w\)-dependence is retained, the localized Maxwell sector contains a
normalizable massless mode plus a discrete tower of massive excitations. For
Gaussian \(\Z(w)\) these modes can be computed in closed form and provide a
numerically natural basis (Hermite spectral) for large-scale simulations.

\paragraph{Sturm--Liouville eigenproblem.}
In axial gauge \(A_w=0\) and (for convenience) Lorenz gauge for the brane
components, a separated expansion
\begin{equation}
\A_\mu(t,\x,w)=\sum_{n\ge 0} a_\mu^{(n)}(t,\x)\,f_n(w)
\label{eq:kk_expansion_em}
\end{equation}
leads to the transverse Sturm--Liouville problem
\begin{equation}
-\frac{\dd}{\dd w}\!\left(\Z(w)\,\frac{\dd f_n}{\dd w}\right)
= m_n^2\,\Z(w)\,f_n(w),
\label{eq:sl_problem_em}
\end{equation}
with inner product \(\langle f,g\rangle_{\Z}=\int \Z(w)f(w)g(w)\,\dd w\).

\paragraph{Gaussian localization: Hermite spectrum.}
For \(\Z(w)=\exp(-w^2/\lambda^2)\), defining \(y=w/\lambda\) turns
\eqref{eq:sl_problem_em} into the Hermite equation. A complete polynomial basis is
\begin{equation}
f_n(w)=H_n\!\left(\frac{w}{\lambda}\right),
\qquad
m_n^2=\frac{2n}{\lambda^2},
\qquad n=0,1,2,\ldots,
\label{eq:hermite_spectrum_em}
\end{equation}
with \(n=0\) the massless (Maxwell) mode and \(n\ge 1\) massive modes.

\paragraph{Norms and brane coupling weights.}
The natural norm induced by the action is
\begin{equation}
\|f_n\|_{\Z}^2 \equiv \int_{-\infty}^{+\infty}\Z(w)\,f_n(w)^2\,\dd w
=\lambda\sqrt{\pi}\,2^n n!.
\label{eq:hermite_norm_em}
\end{equation}
For a brane-localized source, each mode is excited in proportion to its value at
the brane. A convenient dimensionful coupling weight is
\begin{equation}
c_n \equiv \frac{f_n(0)^2}{\|f_n\|_{\Z}^2}.
\label{eq:cn_def_em}
\end{equation}
Since \(H_n(0)=0\) for odd \(n\), odd modes do not couple to brane sources:
\begin{equation}
c_{2m+1}=0.
\label{eq:odd_vanish_em}
\end{equation}
For even \(n=2m\), using \(H_{2m}(0)=(-1)^m(2m)!/m!\),
\begin{equation}
c_{2m}
=\frac{1}{\lambda\sqrt{\pi}}\frac{1}{4^m}\binom{2m}{m},
\qquad
c_0=\frac{1}{\lambda\sqrt{\pi}}=\frac{1}{\Zint}.
\label{eq:cn_even_em}
\end{equation}

\paragraph{Mode-by-mode brane wave operators.}
For conserved brane sources, each mode propagates on the brane with a covariant
massive operator. A convenient summary (in Lorenz gauge) is
\begin{equation}
\left(\square_4 + m_n^2\right)a^{(n)}_\mu(t,\x)
=\muzero\,c_n\,\J_{b\,\mu}(t,\x),
\qquad
\square_4 \equiv -\partial_t^2+\lapthree,
\label{eq:mode_wave_eq_em}
\end{equation}
with \(m_n^2\) and \(c_n\) given by \eqref{eq:hermite_spectrum_em} and
\eqref{eq:cn_def_em}--\eqref{eq:cn_even_em}. The zero mode reproduces Maxwell on the
brane with effective coupling \(\muzeff=\muzero c_0\), while the massive modes are
Proca-like degrees of freedom with dispersion \(\omega^2=\mathbf{k}^2+m_n^2\).

\paragraph{Static brane potential: Coulomb plus Yukawa tower.}
In the static limit, the brane-to-brane response is a sum of Yukawa propagators.
For a fixed point source \(q_{\rm src}\) on the brane, the scalar potential on the brane takes the
form
\begin{equation}
A_0(r)
=\frac{\muzero q_{\rm src}}{4\pi r}\sum_{n\ge 0} c_n\,\ee^{-m_n r},
\qquad r\equiv|\x|.
\label{eq:A0_KK_sum_em}
\end{equation}
Separating the zero mode gives
\begin{equation}
A_0(r)
=\frac{\muzeff q_{\rm src}}{4\pi r}
\left[1+\sum_{m\ge 1}\frac{c_{2m}}{c_0}\,\ee^{-m_{2m}r}\right],
\qquad
m_{2m}=\frac{2\sqrt{m}}{\lambda},
\label{eq:A0_mu_eff_form_em}
\end{equation}
with coefficient ratios
\begin{equation}
\frac{c_{2m}}{c_0}=\frac{1}{4^m}\binom{2m}{m}
=1,\;\frac12,\;\frac{3}{8},\;\frac{5}{16},\ldots\quad (m=0,1,2,3,\ldots).
\label{eq:ratio_cn_em}
\end{equation}
Thus for \(r\gg \lambda\) the leading deviation from Coulomb has fixed coefficient
\(\tfrac12\) and range set by \(\lambda\):
\begin{equation}
A_0(r)
=\frac{\muzeff q_{\rm src}}{4\pi r}
\left[1+\frac12\,\ee^{-2r/\lambda}
+\mathcal{O}\!\left(\ee^{-2\sqrt{2}\,r/\lambda}\right)\right].
\label{eq:leading_yukawa_em}
\end{equation}
Here \(q_{\rm src}\) is a fixed source label; for an elementary defect branch in the
canonically normalized brane packaging one may identify
\(q_{\rm src}=q_{\rm eff}=\eta_Q e_{\rm eff}\).
For plasma modeling, \(\lambda\) therefore sets a concrete crossover scale: at
separations \(r\gg\lambda\) the interaction is effectively Maxwellian, while at
\(r\sim\lambda\) additional transverse modes and mixed components can contribute.
In later sections we show how these same degrees of freedom appear as explicit
beyond-MHD correction channels in the brane-facing fluid and induction equations.

% ============================================================
\section{4+1D Plasma Matter Sector}
\label{sec:matter_sector}

This section specifies the matter sector used to represent a plasma in the bulk
\((t,\x,w)\). The overarching requirement is that the matter variables provide a
conserved bulk current \(\J^M\) suitable for sourcing the localized Maxwell system
(Section~\ref{sec:em_sector}), while also admitting controlled limits that recover
standard brane two-fluid plasma theory and MHD (Section~\ref{sec:mhd_reduction}).

We present two complementary matter realizations:
\begin{enumerate}[leftmargin=1.6em]
\item a multi-species nonrelativistic gauge-covariant field theory (Schr\"odinger/GNLS--Maxwell),
which is particularly well-suited for topology change because the fundamental
field \(\psi_s\) remains regular at \(\rho_s\to 0\) (phase slips are possible);
\item an equivalent kinetic/plasma form (Vlasov--Maxwell in 4 spatial dimensions),
which is well-suited for PIC-style massive simulations.
\end{enumerate}
Both realizations lead to the same bulk charge/current bookkeeping and the same
brane-projected leakage identities (Section~\ref{sec:leakage_identity}).

\subsection{Choice of Plasma Model: Two-Fluid GNLS or Vlasov}
\label{sec:plasma_choice}

\paragraph{Hybrid relativity (scope).}
The electromagnetic sector is a relativistic hyperbolic field theory in \(4+1\)
dimensions (wave propagation at speed \(c\)), while the baseline matter sector is
taken to be nonrelativistic. This ``Schr\"odinger--Maxwell'' or ``two-fluid--Maxwell''
hybrid is standard in plasma modeling in the MHD-like regime. If fully relativistic
particle dynamics is required, the matter sector can be upgraded (e.g.\ relativistic
kinetic theory), but that extension is not needed for the controlled MHD limit
demonstrated here.

\paragraph{Species.}
We consider \(S\) charged species indexed by \(s\in\{1,\ldots,S\}\), each with
mass \(m_s\) and fixed species charge parameter \(q_s\). This paper treats \(q_s\)
as a plasma-model input rather than as a derivation of particle identity. For
elementary defect branches, \(q_s\) may be viewed as the reduced brane charge
inherited from the corrected localized-charge ontology; for composite ions,
\(q_s\) may simply be treated as the usual effective charge state. In
applications, \(s=\ele\) (electrons) and \(s=\ion\) (ions) are the minimal choice.

\subsubsection*{Option A (field-based): multi-species Schr\"odinger/GNLS--Maxwell}

\paragraph{Gauge-covariant derivatives.}
For each species,
\begin{equation}
D_t^{(s)} \equiv \partial_t + \frac{\ii q_s}{\hbar}A_0,
\qquad
D_A^{(s)} \equiv \partial_A - \frac{\ii q_s}{\hbar}A_A,
\qquad A\in\{x,y,z,w\}.
\label{eq:cov_derivatives_species}
\end{equation}

\paragraph{Matter action (nonrelativistic, gauge-covariant).}
A minimal multi-species matter action consistent with the unified model style is
\begin{equation}
\begin{split}
S_{\rm matt}[\{\psi_s\},A]
&=
\sum_{s=1}^S
\int \dd t\,\dd^3x\,\dd w\;
\Biggl[
\frac{\ii\hbar}{2}\Bigl(\psi_s^\ast D_t^{(s)}\psi_s - \psi_s (D_t^{(s)}\psi_s)^\ast\Bigr)
\\
&\qquad
-\frac{\hbar^2}{2m_s}(D_A^{(s)}\psi_s)^\ast(D_A^{(s)}\psi_s)
- V_{{\rm conf},s}(\x,w)\,|\psi_s|^2
- U_s(|\psi_s|^2)
\Biggr],
\end{split}
\label{eq:matter_action}
\end{equation}
where \(V_{{\rm conf},s}\) is an optional confinement/geometry potential and
\(U_s(\rho_s)\) is an equation-of-state (EOS) potential encoding pressure/enthalpy
effects in the hydrodynamic limit. The matter model is \emph{gauge covariant} by
construction, and gauge invariance of the interaction term implies current
conservation (see below).

\paragraph{Field equations (GNLS).}
Variation of \eqref{eq:matter_action} yields, for each species,
\begin{equation}
\ii\hbar D_t^{(s)}\psi_s
=
\left[
-\frac{\hbar^2}{2m_s}D_A^{(s)}D_A^{(s)}
+V_{{\rm conf},s}
+h_s(\rho_s)
\right]\psi_s,
\qquad
\rho_s\equiv|\psi_s|^2,
\qquad
h_s(\rho_s)\equiv \frac{\dd U_s}{\dd \rho_s}.
\label{eq:gnls_species}
\end{equation}
A frequently used stiff-polytrope choice compatible with earlier unified-model
fixes is \(U_s(\rho_s)=\frac{K_s}{4}\rho_s^5\), giving \(h_s(\rho_s)=\frac{5K_s}{4}\rho_s^4\);
however we keep \(U_s\) general in this section because plasma closures may vary by
regime.

\subsubsection*{Option B (kinetic): 4D Vlasov--Maxwell}

As an alternative (or complement) to the field-based realization, each species can
be described by a phase-space distribution \(f_s(t,\x,w;\mathbf{v})\) where
\(\mathbf{v}\in\R^4\) is the bulk spatial velocity \((v_x,v_y,v_z,v_w)\). The
collisionless Vlasov equation in the bulk is
\begin{equation}
\partial_t f_s + v^A\partial_A f_s
+ \frac{q_s}{m_s}\left(E_A + v^B F_{AB}\right)\partial_{v_A} f_s
= 0,
\qquad A,B\in\{x,y,z,w\}.
\label{eq:vlasov_4d}
\end{equation}
Collisions (when desired) can be included through an additional operator
\(C_s[f]\) on the right-hand side (e.g.\ BGK, Landau, or effective friction models),
provided the chosen operator preserves overall charge conservation.

Bulk moments define density and current:
\begin{equation}
\rho_s(t,\x,w)=\int_{\R^4} f_s\,\dd^4 v,
\qquad
j_s^A(t,\x,w)=\int_{\R^4} v^A f_s\,\dd^4 v.
\label{eq:vlasov_moments}
\end{equation}

\subsection{Species Variables, Currents, and Conservation}
\label{sec:species_currents}

\paragraph{Number density and number current.}
In both realizations (GNLS or Vlasov), the species number density is denoted
\(\rho_s(t,\x,w)\). The associated bulk number current is denoted \(j_s^A(t,\x,w)\).
For GNLS,
\begin{equation}
j_s^A
=
\frac{\hbar}{m_s}\,\Im\!\left(\psi_s^\ast D_A^{(s)}\psi_s\right),
\qquad
A\in\{x,y,z,w\}.
\label{eq:gnls_number_current}
\end{equation}
For Vlasov, \(j_s^A\) is the velocity moment in \eqref{eq:vlasov_moments}.

\paragraph{Exact bulk continuity.}
Each species satisfies an exact bulk continuity equation,
\begin{equation}
\partial_t \rho_s + \partial_A j_s^A = 0.
\label{eq:bulk_continuity_species}
\end{equation}
This identity guarantees gauge invariance of the coupling to the Maxwell sector.

\paragraph{Charge density and charge current (Maxwell sources).}
We define the species 5-current \(\Js^M\) by
\begin{equation}
\Js^0 \equiv q_s\rho_s,
\qquad
\Js^A \equiv q_s j_s^A,
\label{eq:species_charge_current}
\end{equation}
and the total conserved source is
\begin{equation}
\J^M \equiv \J_{\rm ext}^M + \sum_{s=1}^S \Js^M,
\qquad
\partial_M \J^M = 0.
\label{eq:total_current_def}
\end{equation}
Here \(\J_{\rm ext}^M\) is reserved for additional sources not represented by the
chosen matter closure (e.g.\ externally imposed drive currents or explicit defect
worldlines in hybrid simulations).

\paragraph{Projected (brane) continuity with leakage.}
Applying the brane projection map \(\Proj{\cdot}\) from
Sec.~\ref{sec:projection_defs} to \eqref{eq:bulk_continuity_species} yields the exact
open-system identity
\begin{equation}
\partial_t \rho_{s,\rm brane} + \partial_a j_{s,\rm brane}^a = \Sleak^{(s)}(t,\x),
\label{eq:projected_continuity_species_repeat}
\end{equation}
with leakage source \(\Sleak^{(s)}\) given explicitly in
\eqref{eq:Sleak_def}. This term vanishes only under additional controlled
assumptions (e.g.\ negligible transverse flux \(j_s^w\) at the measurement scale).

\subsection{Forces and Canonical Momentum}
\label{sec:canonical_momentum}

\subsubsection*{Hydrodynamic (Madelung) representation and canonical momentum}

For the field-based model, write the species wavefunction in amplitude--phase form
(Madelung transform)
\begin{equation}
\psi_s(t,\x,w)=\sqrt{\rho_s(t,\x,w)}\,\exp\!\big(\ii\theta_s(t,\x,w)\big).
\label{eq:madelung}
\end{equation}
Define the bulk species velocity \(\mathbf{v}_s\in\R^4\) by \(j_s^A=\rho_s v_s^A\).
Then
\begin{equation}
v_s^A
=
\frac{\hbar}{m_s}\left(\partial^A\theta_s - \frac{q_s}{\hbar}A^A\right).
\label{eq:velocity_def}
\end{equation}
The corresponding canonical momentum is
\begin{equation}
p^{\rm can}_{s,A} \equiv \hbar\,\partial_A\theta_s,
\qquad
p^{\rm kin}_{s,A}\equiv m_s v_{s,A},
\qquad
p^{\rm can}_{s,A} = p^{\rm kin}_{s,A} + q_s A_A.
\label{eq:canonical_momentum}
\end{equation}
This identity is important computationally: evolving \(\psi_s\) (or equivalently
\(\theta_s\) and \(\rho_s\)) evolves canonical momentum in a way that remains
well-defined even when \(\rho_s\to 0\), whereas raw velocity \(\mathbf{v}_s\) can
become singular in such regions. This feature is one reason the field-based system
is attractive for reconnection/topology-change studies.

\subsubsection*{Euler-like momentum equation (bulk 4D space)}

Inserting \eqref{eq:madelung} into the GNLS equation \eqref{eq:gnls_species} yields
the exact continuity equation \eqref{eq:bulk_continuity_species} and an Euler-like
momentum equation for each species:
\begin{equation}
m_s\left(\partial_t + v_s^B\partial_B\right)v_{s,A}
=
q_s\left(E_A + v_s^B F_{AB}\right)
-\partial_A\!\left(
V_{{\rm conf},s}
+h_s(\rho_s)
+Q_s(\rho_s)
\right),
\label{eq:euler_species_bulk}
\end{equation}
where
\begin{equation}
E_A \equiv F_{A0} = -\partial_t A_A - \partial_A A_0,
\qquad
F_{AB}=\partial_A A_B-\partial_B A_A,
\label{eq:EA_FAB_defs}
\end{equation}
and the quantum (dispersive) potential is
\begin{equation}
Q_s(\rho_s)
\equiv
-\frac{\hbar^2}{2m_s}\frac{\lapfour\sqrt{\rho_s}}{\sqrt{\rho_s}}.
\label{eq:quantum_potential}
\end{equation}
The classical two-fluid limit is obtained by dropping \(Q_s\) (or taking a controlled
\(\hbar\to 0\) limit) while retaining the barotropic closure encoded in \(h_s(\rho_s)\).

\paragraph{Transverse forces and the role of \texorpdfstring{$A_w$}{Aw}.}
Equation \eqref{eq:euler_species_bulk} holds for all bulk spatial components
\(A\in\{x,y,z,w\}\). In particular, the \(A=w\) component reads
\begin{equation}
m_s\left(\partial_t + v_s^B\partial_B\right)v_{s,w}
=
q_s\left(E_w + v_s^a F_{wa}\right)
-\partial_w\!\left(
V_{{\rm conf},s}
+h_s
+Q_s
\right),
\label{eq:euler_species_w}
\end{equation}
with \(E_w=F_{w0}=-\partial_t A_w-\partial_w A_0\). Thus a dynamical \(A_w\) (and/or
mixed components \(F_{wa}\)) generically drives transverse drift \(v_{s,w}\neq 0\),
which in turn produces \(j_s^w\neq 0\) and therefore nonzero projected leakage
\(\Sleak^{(s)}\) on the brane. This is one of the primary 4D interaction channels
that will appear as beyond-MHD corrections in Section~\ref{sec:beyond_mhd}.

\paragraph{Collisions and effective resistivity (optional).}
When a resistive or collisional regime is desired, one may include momentum-exchange
terms in the species momentum equations (or collision operators in Vlasov) that
preserve overall charge conservation. For example, an interspecies drag term of the
form \(-m_s\nu_{ss'}(v_s^A-v_{s'}^A)\) can be added to the right-hand side of
\eqref{eq:euler_species_bulk}; in the controlled brane reduction such terms generate
standard resistive contributions to Ohm's law. In the present framework, however,
non-ideal behavior may also arise \emph{without} such dissipative closures via
brane--bulk exchange and transverse-mode excitation.

\paragraph{Useful kinematic identity (canonical vorticity).}
From \eqref{eq:velocity_def} one obtains, away from phase singularities,
\begin{equation}
\partial_A v_{s,B}-\partial_B v_{s,A}
= -\frac{q_s}{m_s}F_{AB}.
\label{eq:london_identity_bulk}
\end{equation}
This ``London-type'' identity states that canonical vorticity is locked to gauge
flux as a bulk geometric constraint; its brane-facing implications for topology,
reconnection, and projected non-ideality are developed in
Section~\ref{sec:vorticity_identity}.
It is a vortical/canonical-momentum constraint and should not be read as a rule
that defines electric charge from circulation.

% ============================================================
\section{Full Coupled 4+1D Plasma--EM System}
\label{sec:full_system}

We now collect the coupled bulk equations that define the simulation target: a
localized \(4+1\) Maxwell sector (Section~\ref{sec:em_sector}) sourced by a
multi-species plasma matter sector (Section~\ref{sec:matter_sector}). The resulting
system is conservative in the bulk (up to stated boundary/localization conditions)
and admits controlled reductions to standard \(3+1\) Maxwell theory and MHD on the
brane (Section~\ref{sec:mhd_reduction}). Throughout, the brane is not a boundary of
the PDE domain; brane-facing quantities are defined by the projection operator
\(\Proj{\cdot}\) (Section~\ref{sec:projection_defs}) and may exhibit explicit
exchange (leakage) terms even when bulk dynamics is conservative.

\subsection{Coupled action and equations of motion}
\label{sec:coupled_action_eom}

\paragraph{Total action (bulk).}
A convenient parent formulation is the sum of the localized Maxwell action
\eqref{eq:action_em} and a multi-species matter action (field-based option shown
here; kinetic option below):
\begin{equation}
S_{\rm tot}
=
S_{\rm EM}[\A;\Z]
+
\sum_{s=1}^S S_{{\rm matt},s}[\psi_s,\A]
\;+\;
S_{\rm ext}[\A],
\label{eq:action_total}
\end{equation}
where \(S_{\rm ext}\) allows additional prescribed sources (drive currents,
background fields, etc.). The localization profile \(\Z(w)\) is fixed (time
independent) and integrable.

\paragraph{Maxwell equation with localization.}
Variation with respect to \(\A_N\) yields
\begin{equation}
\partial_M\!\left(\Z(w)\F^{MN}\right)
+\frac{1}{\xi}\,\partial^N(\partial\!\cdot\!\A)
=\muzero\,\J_{\rm tot}^N,
\qquad
\J_{\rm tot}^N \equiv \J_{\rm ext}^N + \sum_{s=1}^S \Js^N,
\label{eq:maxwell_full}
\end{equation}
with the exact Bianchi identity \(\partial_{[L}\F_{MN]}=0\).
For conserved sources \(\partial_N\J_{\rm tot}^N=0\), the gauge-fixing term does not
affect gauge-invariant observables.

\paragraph{Matter evolution: field-based (GNLS) option.}
For each species \(s\), the gauge-covariant derivatives are
\begin{equation}
D_t^{(s)} \equiv \partial_t + \frac{\ii q_s}{\hbar}A_0,
\qquad
D_A^{(s)} \equiv \partial_A - \frac{\ii q_s}{\hbar}A_A,
\qquad A\in\{x,y,z,w\},
\label{eq:covD_repeat_full}
\end{equation}
and the species field equation is
\begin{equation}
\ii\hbar D_t^{(s)}\psi_s
=
\left[
-\frac{\hbar^2}{2m_s}D_A^{(s)}D_A^{(s)}
+V_{{\rm conf},s}(\x,w)
+h_s(\rho_s)
\right]\psi_s,
\qquad
\rho_s\equiv|\psi_s|^2,
\qquad
h_s(\rho_s)\equiv \frac{\dd U_s}{\dd\rho_s}.
\label{eq:gnls_full}
\end{equation}
The confinement potential \(V_{{\rm conf},s}\) may be used to localize plasma
species near the operational brane or near defect/throat geometries; for an
unconstrained bulk plasma one sets \(V_{{\rm conf},s}\equiv 0\).

\paragraph{Matter evolution: kinetic (Vlasov) option.}
Alternatively, each species may be represented by a bulk phase-space distribution
\(f_s(t,\x,w;\mathbf{v})\) with \(\mathbf{v}\in\R^4\), evolving via
\begin{equation}
\partial_t f_s + v^A\partial_A f_s
+ \frac{q_s}{m_s}\left(E_A + v^B F_{AB}\right)\partial_{v_A} f_s
= C_s[f],
\label{eq:vlasov_full}
\end{equation}
where \(C_s[f]\) is an optional collision operator chosen to preserve overall
charge conservation. The electromagnetic fields are defined from \(\F_{MN}\) by
\(E_A\equiv F_{A0}\) and \(F_{AB}\) as in \eqref{eq:EA_FAB_defs}.

\paragraph{Species sources (currents).}
In the GNLS representation, the bulk number current is
\begin{equation}
j_s^A
=
\frac{\hbar}{m_s}\Im\!\left(\psi_s^\ast D_A^{(s)}\psi_s\right),
\label{eq:number_current_full}
\end{equation}
and the species 5-current sourcing Maxwell is
\begin{equation}
\Js^0 \equiv q_s\rho_s,
\qquad
\Js^A \equiv q_s j_s^A.
\label{eq:charge_current_full}
\end{equation}
In the kinetic representation, \(\rho_s\) and \(j_s^A\) are the moments
\eqref{eq:vlasov_moments}. In either case, each species satisfies exact bulk
continuity
\begin{equation}
\partial_t\rho_s+\partial_A j_s^A=0
\qquad\Rightarrow\qquad
\partial_M \Js^M=0,
\label{eq:species_conservation_full}
\end{equation}
and therefore \(\partial_M\J_{\rm tot}^M=0\) when \(\partial_M\J_{\rm ext}^M=0\).

\paragraph{Simulation target (compact statement).}
For later reference, the coupled bulk system to be evolved can be written as:
\begin{equation}
\boxed{
\begin{aligned}
&\partial_M\!\left(\Z(w)\F^{MN}\right)
+\frac{1}{\xi}\,\partial^N(\partial\!\cdot\!\A)
=\muzero\!\left(\J_{\rm ext}^N+\sum_s \Js^N\right),
\\
&\text{either}\quad
\ii\hbar D_t^{(s)}\psi_s
=
\left[
-\frac{\hbar^2}{2m_s}D_A^{(s)}D_A^{(s)}
+V_{{\rm conf},s}+h_s(\rho_s)
\right]\psi_s
\quad (\text{field/GNLS}),
\\
&\text{or}\quad
\partial_t f_s + v^A\partial_A f_s
+ \frac{q_s}{m_s}\left(E_A + v^B F_{AB}\right)\partial_{v_A} f_s
= C_s[f]
\quad (\text{kinetic/Vlasov}),
\\
&\Js^0=q_s\rho_s,\quad \Js^A=q_s j_s^A,
\qquad
\partial_M \Js^M=0,
\qquad
\partial_{[L}\F_{MN]}=0.
\end{aligned}
}
\label{eq:full_system_box}
\end{equation}

\subsection{Bulk constraints, gauge choices, and well-posedness}
\label{sec:constraints_gauge_wellposed}

\paragraph{Constraints.}
When evolving potentials \(\A_M\), the homogeneous Maxwell equations
\eqref{eq:bianchi} are automatically satisfied. The inhomogeneous equations
\eqref{eq:maxwell_full} contain constraint content in the \(N=0\) component (a
Gauss-type law) and evolution content in the \(N=A\) components (Amp\`ere-type
laws). The gauge-fixing term in \eqref{eq:maxwell_full} promotes the system to a
manifestly hyperbolic form suitable for time-domain simulation.

\paragraph{Gauge handling.}
We do not assume \(A_w=0\) in the full theory; the \(A_w\) sector and the mixed
components \(\F_{\mu w}\) are retained as physical degrees of freedom and are a key
source of beyond-MHD corrections on the brane (Section~\ref{sec:beyond_mhd}). For
analysis and for transverse mode decompositions it can be convenient to impose
axial gauge \(A_w=0\) with residual brane gauge freedom \(\chi_b(t,\x)\), but this
is treated as a \emph{controlled} simplification rather than a defining assumption.

\paragraph{Localization and boundary conditions.}
Because \(\Z(w)\) breaks translation invariance in the \(w\)-direction, conserved
Noether momenta along \(w\) are not expected in general. Time translation symmetry
is preserved (\(\Z\) is time independent), so the system admits a conserved total
energy under suitable boundary/localization conditions. In practice, one assumes
either (i) sufficiently fast decay of fields/currents as \(|w|\to\infty\), or (ii)
effective compact support in \(w\) as realized by spectral/Hermite truncation or
numerical domain truncation (Section~\ref{sec:computation}).

\subsection{Energy ledger and conservation (bulk)}
\label{sec:bulk_energy}

A major advantage of the action-based formulation is that it provides a clean
energy ledger for verification and diagnostics in simulations. Here we record a
gauge-invariant bulk energy functional and its exchange structure between matter
and fields.

\paragraph{Electromagnetic energy density (localized).}
Define the bulk electric-type and magnetic-type components by
\begin{equation}
E_A \equiv \F_{A0},
\qquad
\F_{AB}=\partial_A A_B-\partial_B A_A,
\qquad A,B\in\{x,y,z,w\}.
\end{equation}
A convenient gauge-invariant electromagnetic energy density is
\begin{equation}
u_{\rm EM}
\equiv
\frac{\Z(w)}{2\muzero}
\left(
E_A E_A
+
\frac12\,\F_{AB}\F_{AB}
\right).
\label{eq:u_em_def}
\end{equation}
In terms of brane and mixed components, this contains contributions from the usual
brane fields \(\F_{ab}\) as well as from the mixed sector \(\F_{aw}\) and the
transverse electric component \(E_w=\F_{w0}\).

\paragraph{Matter energy density (field-based option).}
For the GNLS realization, a gauge-invariant matter energy density for each species
is
\begin{equation}
u_s
\equiv
\frac{\hbar^2}{2m_s}\,(D_A^{(s)}\psi_s)^\ast(D_A^{(s)}\psi_s)
+V_{{\rm conf},s}\,\rho_s
+U_s(\rho_s),
\label{eq:u_matt_def}
\end{equation}
so that the total matter energy density is \(u_{\rm matt}\equiv\sum_s u_s\).
(When desired, \eqref{eq:u_matt_def} can be rewritten in hydrodynamic variables as
kinetic \(+\) internal \(+\) dispersive/quantum contributions via the Madelung map.)

\paragraph{Total energy and exchange structure.}
Define the total bulk energy
\begin{equation}
\mathcal{E}_{\rm tot}(t)
\equiv
\int_{\R^3}\dd^3x\int_{\R}\dd w\;
\left(
u_{\rm EM} + u_{\rm matt}
\right),
\label{eq:E_tot_def}
\end{equation}
with \(u_{\rm EM}\) from \eqref{eq:u_em_def} and \(u_{\rm matt}\) from
\eqref{eq:u_matt_def} (or from kinetic moments in the Vlasov realization).
Under standard localization/boundary assumptions and for conservative collision
models, \(\mathcal{E}_{\rm tot}\) is conserved. Locally, the EM--matter coupling
exchanges energy through the work density \(J^A E_A\), which includes the transverse
channel \(J^w E_w\) whenever the scalar-photon sector is active:
\begin{equation}
\partial_t u_{\rm EM} + \partial_A S_{\rm EM}^A
=
-\,J_{\rm tot}^A E_A,
\qquad
\partial_t u_{\rm matt} + \partial_A S_{\rm matt}^A
=
+\,J_{\rm tot}^A E_A,
\label{eq:energy_exchange}
\end{equation}
for suitable (model-dependent) fluxes \(S_{\rm EM}^A\) and \(S_{\rm matt}^A\).
A convenient covariant generalization of the Poynting flux is
\begin{equation}
S_{\rm EM}^A \equiv \frac{\Z(w)}{\muzero}\,E_B\,\F_{BA},
\label{eq:poynting_4d}
\end{equation}
which reduces to the usual \(\E\times\B\) form when restricted to brane components.

\paragraph{Why the energy ledger matters for reconnection.}
In standard resistive MHD, reconnection is often diagnosed by conversion of magnetic
energy into internal energy through \(\eta J^2\). In the present model, reconnection-like
behavior in brane observables can occur through conservative transfer of energy into
transverse degrees of freedom: \(A_w\), \(E_w\), mixed components \(\F_{aw}\), and/or
higher transverse modes (Section~\ref{sec:kk_modes}). The energy ledger
\eqref{eq:E_tot_def}--\eqref{eq:energy_exchange} provides an unambiguous way to
distinguish genuine dissipation (collisions) from topology/energy export into bulk
channels.

% ============================================================
\section{Controlled Reduction: 4D \texorpdfstring{$\to$}{->} 3+1 Maxwell \texorpdfstring{$\to$}{->} Two-Fluid \texorpdfstring{$\to$}{->} MHD}
\label{sec:mhd_reduction}

This section establishes that standard \(3+1\) Maxwell theory and classical MHD arise
as \emph{controlled limits} of the full \(4+1\) bulk plasma--EM system
(Section~\ref{sec:full_system}). The reductions proceed in three stages:
\begin{equation}
\begin{aligned}
&\text{(bulk localized Maxwell + bulk plasma)}
\\
&\Longrightarrow \text{(effective brane Maxwell + brane two-fluid)}
\\
&\Longrightarrow \text{(single-fluid MHD)}.
\end{aligned}
\end{equation}
The purpose of these reductions is twofold:
(i) to show that MHD is recoverable and therefore is not contradicted by the parent
theory, and (ii) to expose precisely which assumptions suppress the new \(4\)D
interaction channels (leakage \(J^w\), scalar-photon sector \(A_w\), and transverse
mode excitation) that will be reinstated in Section~\ref{sec:beyond_mhd}.

\subsection{Controlled assumptions and ordering}
\label{sec:controlled_assumptions}

We group assumptions by reduction step. Each group defines an approximation regime
in which the subsequent reduced equations are valid.

\subsubsection*{Step 1: localized 4+1 Maxwell \(\to\) effective 3+1 Maxwell}
We adopt the brane-dominant zero-mode approximation on the brane components:
\begin{enumerate}[leftmargin=1.6em]
\item \textbf{Zero-mode dominance in the localization region:}
\begin{equation}
A_\mu(t,\x,w)\simeq a_\mu(t,\x),
\qquad
\partial_w A_\mu \simeq 0,
\label{eq:zero_mode_assumption}
\end{equation}
for \(|w|\) in the support of \(\Z(w)\).

\item \textbf{Negligible mixed flux in the controlled limit:}
\begin{equation}
F_{\mu w}\simeq 0,
\qquad\text{equivalently}\qquad
E_w=F_{w0}\simeq 0,\;\; C_a=F_{aw}\simeq 0,
\label{eq:mixed_suppression_assumption}
\end{equation}
so that the brane-facing dynamics is dominated by \(F_{\mu\nu}\).

\item \textbf{No net transverse current at the reduction scale:}
\begin{equation}
J^w \simeq 0.
\label{eq:Jw_suppression_assumption}
\end{equation}
(When this fails, brane observables become open-system variables with explicit
exchange; see Section~\ref{sec:beyond_mhd}.)

\item \textbf{Gauge choice (optional):} axial gauge \(A_w=0\) may be imposed for
convenience, but is not required for the full theory.
\end{enumerate}

Under these assumptions, the brane components reduce to standard \(3+1\) Maxwell
equations with an effective coupling
\begin{equation}
\muzeff \equiv \frac{\muzero}{\Zint},
\qquad
\Zint \equiv \int_{-\infty}^{+\infty}\Z(w)\,\dd w.
\label{eq:mu0eff_repeat}
\end{equation}
If desired, define the corresponding effective permittivity by
\(\epsilon_0^{\rm eff}\mu_0^{\rm eff}=c^{-2}\).
For an elementary defect branch with \(q_*=\eta_Q e_*\), \(\eta_Q=\pm1\), and
\(e_*>0\), the same localization dependence may equivalently be packaged into the
fixed reduced brane charge \(q_{\rm eff}=q_*/\sqrt{\Zint}\); in this plasma paper
we keep \(\mu_0^{\rm eff}\) as the primary Maxwell/MHD notation.
Suppressing \((A_w,F_{\mu w},J^w)\) here is a controlled far-field brane reduction
used to recover standard plasma theory, not a microscopic statement that the
underlying charged core lacks mixed-sector structure.

\subsubsection*{Step 2: bulk plasma \(\to\) brane two-fluid}
To obtain standard brane plasma equations, we assume the species are effectively
brane-supported at the measurement scale:
\begin{equation}
j_s^w \simeq 0
\qquad\Rightarrow\qquad
\Sleak^{(s)} \simeq 0,
\label{eq:twofluid_leakage_suppression}
\end{equation}
so that projected continuity reduces to ordinary \(3\)D continuity. We further assume
a classical fluid closure for each species, i.e. drop dispersive/quantum terms
(\(Q_s\)) when starting from the GNLS realization.

\subsubsection*{Step 3: brane two-fluid \(\to\) MHD}
The standard MHD reduction is obtained by the familiar controlled assumptions:
\begin{enumerate}[leftmargin=1.6em]
\item \textbf{Quasi-neutrality:} \(n_\ion \simeq n_\ele \equiv n\).
\item \textbf{Electron inertia suppression:} \(m_\ele/m_\ion \ll 1\) and electron
inertia terms are negligible in Ohm's law.
\item \textbf{Low-frequency / long-scale:} optionally neglect displacement current in
Amp\`ere's law (Darwin/MHD ordering).
\end{enumerate}
Relaxing any of these yields standard extended-MHD corrections (Hall, electron pressure,
electron inertia, resistivity). In the present framework, additional beyond-MHD
corrections also arise when the Step-1 and Step-2 assumptions are relaxed (Section~\ref{sec:beyond_mhd}).

\subsection{Two-fluid algebra on the brane}
\label{sec:two_fluid_algebra}

Under the Step-1 and Step-2 assumptions, we work with brane \(3\)-vectors
\(\E(t,\x)\), \(\B(t,\x)\), and brane currents. The effective Maxwell system is
\begin{align}
\gradthree\cdot \B &= 0,
&
\gradthree\times \E &= -\partial_t \B,
\label{eq:maxwell_hom_brane}
\\
\gradthree\cdot \E &= \frac{\rho_q}{\epsilon_0^{\rm eff}},
&
\gradthree\times \B - \frac{1}{c^2}\partial_t \E &= \muzeff\,\Jb,
\label{eq:maxwell_inhom_brane}
\end{align}
with charge density \(\rho_q\) and current density \(\Jb\) defined on the brane.

\paragraph{Brane two-fluid equations.}
For species \(s\in\{\ion,\ele\}\) with charge \(q_\ion=+e\), \(q_\ele=-e\), we take
\begin{align}
\partial_t n_s + \gradthree\cdot(n_s \mathbf{v}_s) &= 0,
\label{eq:twofluid_cont}
\\
m_s n_s\left(\partial_t + \mathbf{v}_s\cdot\gradthree\right)\mathbf{v}_s
&=
q_s n_s\left(\E + \mathbf{v}_s\times\B\right) - \gradthree p_s + \mathbf{R}_s,
\label{eq:twofluid_mom}
\end{align}
where \(p_s\) is species pressure and \(\mathbf{R}_s\) represents optional collisional
momentum exchange (e.g.\ electron--ion drag), with \(\sum_s \mathbf{R}_s=0\).
Here \(e>0\) denotes the fixed brane-observed elementary charge magnitude
(equivalently \(e_{\rm eff}\) in the localized-defect reading).

\paragraph{Single-fluid variables.}
Define the (approximate) mass density, center-of-mass velocity, and current:
\begin{equation}
\rhob \equiv m_\ion n_\ion + m_\ele n_\ele \simeq (m_\ion+m_\ele)n,
\qquad
\mathbf{v} \equiv \frac{m_\ion n_\ion \mathbf{v}_\ion + m_\ele n_\ele \mathbf{v}_\ele}{\rhob},
\qquad
\Jb \equiv e n(\mathbf{v}_\ion-\mathbf{v}_\ele).
\label{eq:mhd_defs}
\end{equation}
Under quasi-neutrality and singly-ionized ions, these relations can be inverted to
express species velocities in terms of $\mathbf{v}$ and $\Jb$:
\begin{equation}
\mathbf{v}_\ion
=
\mathbf{v} + \frac{m_\ele}{m_\ion+m_\ele}\frac{\Jb}{e n},
\qquad
\mathbf{v}_\ele
=
\mathbf{v} - \frac{m_\ion}{m_\ion+m_\ele}\frac{\Jb}{e n}.
\label{eq:vi_ve_exact}
\end{equation}
In the electron-inertia-suppressed limit \(m_\ele/m_\ion\to 0\),
\begin{equation}
\mathbf{v}_\ion \to \mathbf{v},
\qquad
\mathbf{v}_\ele \to \mathbf{v} - \frac{\Jb}{e n}.
\label{eq:vi_ve_me0}
\end{equation}

\paragraph{Generalized Ohm's law (electron momentum).}
Start from the electron momentum equation \eqref{eq:twofluid_mom}. Neglecting electron inertia
and (optionally) retaining collisions via \(\mathbf{R}_\ele\), one obtains
\begin{equation}
0
=
-e n\left(\E + \mathbf{v}_\ele\times\B\right) - \gradthree p_\ele + \mathbf{R}_\ele.
\label{eq:electron_momentum_inertialess}
\end{equation}
Solving for \(\E\) and substituting \(\mathbf{v}_\ele\) from \eqref{eq:vi_ve_me0} yields
\begin{equation}
\E + \mathbf{v}\times\B
=
\underbrace{\frac{\Jb\times\B}{e n}}_{\text{Hall}}
\;-\;
\underbrace{\frac{1}{e n}\gradthree p_\ele}_{\text{electron pressure}}
\;+\;
\underbrace{\frac{\mathbf{R}_\ele}{e n}}_{\text{collisions/resistivity}}.
\label{eq:ohm_general}
\end{equation}
The standard resistive closure corresponds to \(\mathbf{R}_\ele \approx e n\,\eta\,\Jb\),
giving \(\E + \mathbf{v}\times\B = (\Jb\times\B)/(en) - (\gradthree p_\ele)/(en) + \eta \Jb\).
The \emph{ideal MHD} limit is the controlled sub-limit in which the right-hand side of
\eqref{eq:ohm_general} is negligible:
\begin{equation}
\E = -\mathbf{v}\times\B.
\label{eq:ideal_ohm}
\end{equation}

\paragraph{Single-fluid momentum equation.}
Adding the ion and electron momentum equations \eqref{eq:twofluid_mom} and using
\(\sum_s \mathbf{R}_s=0\) yields the single-fluid momentum balance
\begin{equation}
\rhob\left(\partial_t + \mathbf{v}\cdot\gradthree\right)\mathbf{v}
=
\Jb\times\B - \gradthree p,
\qquad
p \equiv p_\ion + p_\ele,
\label{eq:mhd_momentum}
\end{equation}
up to the usual higher-order corrections (electron inertia, viscosity, anisotropic pressure)
that are included in extended-MHD closures if required.

\subsection{Induction equation and ideal MHD limit}
\label{sec:ideal_mhd}

Faraday's law \eqref{eq:maxwell_hom_brane} gives
\begin{equation}
\partial_t \B = -\gradthree\times \E.
\label{eq:faraday_again}
\end{equation}
Inserting the generalized Ohm law \eqref{eq:ohm_general} yields the generalized induction equation
\begin{equation}
\partial_t \B
=
\gradthree\times(\mathbf{v}\times\B)
-\gradthree\times\!\left(\frac{\Jb\times\B}{e n}\right)
+\gradthree\times\!\left(\frac{1}{e n}\gradthree p_\ele\right)
-\gradthree\times\!\left(\frac{\mathbf{R}_\ele}{e n}\right).
\label{eq:induction_general}
\end{equation}
In the ideal MHD limit \eqref{eq:ideal_ohm}, this reduces to the standard ideal induction equation
\begin{equation}
\partial_t \B = \gradthree\times(\mathbf{v}\times\B).
\label{eq:induction_ideal}
\end{equation}

\paragraph{Low-frequency Amp\`ere closure (optional).}
In many MHD regimes one further neglects displacement current, reducing
\eqref{eq:maxwell_inhom_brane} to
\begin{equation}
\gradthree\times\B \simeq \muzeff\,\Jb,
\label{eq:ampere_mhd}
\end{equation}
so that \(\Jb\) may be eliminated in favor of \(\B\) in \eqref{eq:mhd_momentum} and
\eqref{eq:induction_general}. This closure is controlled by \(\omega/(ck)\) and is
not essential to the logical derivation of ideal MHD from \eqref{eq:ideal_ohm}.

\paragraph{Summary: standard MHD system.}
A common ideal-MHD closure in this notation consists of
\begin{align}
\partial_t \rhob + \gradthree\cdot(\rhob \mathbf{v}) &= 0,
\label{eq:mhd_continuity}
\\
\rhob\left(\partial_t + \mathbf{v}\cdot\gradthree\right)\mathbf{v} &= \Jb\times\B - \gradthree p,
\label{eq:mhd_mom_again}
\\
\partial_t \B &= \gradthree\times(\mathbf{v}\times\B),
\qquad
\gradthree\cdot\B=0,
\label{eq:mhd_induction_again}
\end{align}
supplemented by an equation of state and (optionally) \eqref{eq:ampere_mhd} to relate
\(\Jb\) and \(\B\). The only modification relative to standard SI MHD is that the
effective coupling \(\mu_0\) is replaced by \(\muzeff=\mu_0/\Zint\), inherited from the
localized \(4+1\) Maxwell sector.

\subsection{Symbolic verification (WL referee suite)}
\label{sec:wl_verification}

The two-fluid algebra leading to \eqref{eq:vi_ve_exact}--\eqref{eq:ohm_general} and the
ideal induction equation \eqref{eq:induction_ideal} was verified symbolically using a
Wolfram Language referee script (\texttt{4d\_to\_mhd\_reduction.wl}) maintained alongside
the project sources. The script:
\begin{enumerate}[leftmargin=1.6em]
\item encodes the controlled brane/zero-mode assumptions for the EM sector and records
\(\mu_0^{\rm eff}=\mu_0/\Zint\);
\item solves the linear system relating \((\mathbf{v}_\ion,\mathbf{v}_\ele)\) to
\((\mathbf{v},\Jb)\);
\item derives the inertialess-electron generalized Ohm law and expresses it in the
standard form \(\E+\mathbf{v}\times\B=\cdots\);
\item combines Faraday's law with the ideal Ohm sub-limit to yield
\(\partial_t\B=\gradthree\times(\mathbf{v}\times\B)\);
\item verifies the Gaussian soft-wall identities used later in the paper:
\(Z_{\rm int}=\lambda\sqrt{\pi}\), \(m_n^2=2n/\lambda^2\), the weighted Hermite norms,
and the odd/even brane-coupling selection rules;
\item verifies the exact bookkeeping behind the projected generalized Ohm law,
including \(\boldsymbol{\mathcal{E}}_{\rm cov}\), \(\boldsymbol{\mathcal{E}}_w\),
and the assembled topology EMF \(\boldsymbol{\mathcal{E}}_{\rm topo}\);
\item checks representative open-system and geometry-ledger identities used in the
appendices, including the projected \(w\)-leakage integration-by-parts identity and
the ratio-penalty / tube-geometry force derivatives.
\end{enumerate}
We do not include the script source code in this paper; it is intended as a
machine-checkable supplement to the human-readable derivation above.

\paragraph{Where the suppressed 4D terms re-enter.}
The script therefore does more than record reminders: it also machine-checks the
exact appendix-level identities that expose the three primary ways in which the
controlled MHD reduction can fail in the parent theory:
(i) brane leakage in continuity/momentum via \(j^w\) (and thus \(\Sleak\neq 0\)),
(ii) additional force channels mediated by the \(A_w\) sector and mixed components
\(F_{\mu w}\), and (iii) finite-localization transverse-mode (KK/Hermite) corrections that
modify effective propagation and storage of field energy. These are developed as beyond-MHD
extensions in Section~\ref{sec:beyond_mhd}.

% ============================================================
\section{Beyond-MHD Extensions from 4D Interactions}
\label{sec:beyond_mhd}

Section~\ref{sec:mhd_reduction} established that standard MHD arises as a \emph{controlled}
brane/zero-mode limit of the full \(4+1\)D plasma--EM system. The point of that exercise
was not merely to recover known equations, but to identify precisely which assumptions
suppress the degrees of freedom that are absent in strictly \(3+1\) electromagnetism.
In this section we relax those assumptions in a controlled way and classify the resulting
\emph{beyond-MHD} correction channels.

The organizing principle is simple:
\begin{quote}
\emph{The bulk evolution is conservative and local in \(4+1\)D; the brane-facing system is an
open subsystem defined by an operational projection map \(\overline{(\cdot)}\).}
\end{quote}
Consequently, ``non-ideal'' brane behavior can arise from
(i) explicit exchange with the transverse direction \(w\) and
(ii) projection covariances (unresolved transverse structure),
even when no ad hoc dissipation is introduced in the bulk.

\subsection{What changes when we move beyond the controlled limit}
\label{sec:beyond_mhd_channels}

The controlled MHD reduction relied on suppressing transverse structure and mixed EM components:
\begin{equation}
\partial_w(\cdot)\approx 0,
\qquad
A_w\approx 0,
\qquad
F_{\mu w}\approx 0,
\qquad
J^w\approx 0,
\qquad
\text{(zero-mode dominance)}.
\label{eq:controlled_limit_suppression_recall}
\end{equation}
Beyond MHD, we explicitly \emph{allow} these quantities to be non-negligible. The resulting
new physics falls into three coupled classes:

\begin{enumerate}[leftmargin=1.6em]
\item \textbf{Mixed-sector electromagnetism (scalar-photon channel).}
The component \(A_w\) is a genuine dynamical degree of freedom in \(4+1\)D gauge theory.
It generates the transverse electric field
\begin{equation}
E_w \equiv F_{w0} = -\partial_t A_w - \partial_w A_0,
\label{eq:Ew_def_beyond}
\end{equation}
and together with the mixed components
\begin{equation}
C_a \equiv F_{aw} = \partial_a A_w - \partial_w A_a,
\qquad a\in\{x,y,z\},
\label{eq:Ca_def_beyond}
\end{equation}
it introduces new force channels that are invisible in strictly \(3+1\) Maxwell theory.

\item \textbf{Explicit brane--bulk exchange (leakage).}
Species can drift in \(w\), producing \(j_s^w\neq 0\) and therefore a projected
(open-system) continuity law on the brane. Exchange with \(w\) is not a numerical artifact;
it is a physical degree of freedom of the parent theory.

\item \textbf{Finite-localization / transverse-mode corrections.}
When the localization profile \(Z(w)\) is integrable (Gaussian ``soft wall'' in particular),
the brane-facing EM sector contains a tower of transverse modes with masses
\(m_n^2\propto n/\lambda^2\) (Appendix~\ref{sec:kk_modes}). These modes provide additional
storage/transport channels for energy and topology. In the static limit, they generate
Yukawa corrections to Coulomb interactions at scales \(r\sim\lambda\)
(Appendix~\ref{sec:kk_modes}).
\end{enumerate}

In the remainder of this section we write the brane-facing balances in a form that separates
\emph{standard extended-MHD terms} (Hall/pressure/inertia/resistivity) from the new
\emph{4D interaction terms} (mixed-sector forces, leakage, and projection covariances).

\subsection{Open-system brane balances: leakage and projection covariances}
\label{sec:open_system_brane_balances}

The brane-observed quantity \(\overline{Q}(t,\mathbf{x})\) is defined by the operational
projection kernel \(W(w)\) as
\begin{equation}
\overline{Q}(t,\mathbf{x}) \equiv \int_{-\infty}^{+\infty}\dd w\,W(w)\,Q(t,\mathbf{x},w),
\qquad
\int \dd w\,W(w)=1.
\label{eq:proj_def_beyond}
\end{equation}
Because products do not commute with projection, it is convenient to define the covariance
operator
\begin{equation}
\mathrm{Cov}_W[Q,R] \equiv \overline{QR}-\overline{Q}\,\overline{R}.
\label{eq:cov_def_beyond}
\end{equation}
These covariances are \emph{exact} and represent unresolved transverse structure (higher
\(w\)-modes in the Hermite/KK language, or fine structure on a direct \(w\)-grid).

\subsubsection*{Species continuity with leakage}

Let \(\rho_s(t,\mathbf{x},w)\) and \(j_s^A(t,\mathbf{x},w)\) be the bulk number density and
number current for species \(s\), with \(A\in\{x,y,z,w\}\). The bulk continuity equation is
\begin{equation}
\partial_t \rho_s + \partial_a j_s^a + \partial_w j_s^w = 0.
\label{eq:bulk_species_cont_beyond}
\end{equation}
Projecting with \(W(w)\) yields a brane-facing open-system balance:
\begin{equation}
\partial_t \overline{\rho_s} + \partial_a \overline{j_s^a}
=
S_{\mathrm{leak}}^{(s)},
\label{eq:proj_species_cont_beyond}
\end{equation}
where the leakage/source term is
\begin{equation}
S_{\mathrm{leak}}^{(s)}
\equiv
- \int \dd w\,W(w)\,\partial_w j_s^w
=
-\Big[W j_s^w\Big]_{-\infty}^{+\infty}
+\int \dd w\,W'(w)\,j_s^w.
\label{eq:Sleak_def_beyond}
\end{equation}
When the boundary term vanishes, \(S_{\mathrm{leak}}^{(s)}=\int \dd w\,W'(w)\,j_s^w\).
This term is measured directly in both simulation architectures:
direct \(w\)-grid codes compute it from \(j_s^w\); 3D\(\times\)modes codes compute it
from reconstructed \(j_s^w(t,\mathbf{x},w_q)\) at quadrature nodes.

\subsubsection*{Brane-facing momentum and EMF covariances}

Whenever a brane-facing equation contains a projected product, covariances appear.
The most important example for magnetic topology is the projected electromotive term
\(\overline{\mathbf{v}\times\mathbf{B}}\), which differs from
\(\overline{\mathbf{v}}\times\overline{\mathbf{B}}\) by an \emph{exact covariance}:
\begin{equation}
\overline{\mathbf{v}\times\mathbf{B}}
=
\overline{\mathbf{v}}\times\overline{\mathbf{B}}
+
\underbrace{\mathrm{Cov}_W[\mathbf{v},\mathbf{B}]_{\times}}_{\equiv\;\overline{\mathbf{v}'\times\mathbf{B}'}}.
\label{eq:EMF_cov_identity_beyond}
\end{equation}
We will interpret \(\overline{\mathbf{v}'\times\mathbf{B}'}\) as a component of a
\emph{topology EMF} because it acts as an additional non-ideal source in the brane induction
equation (see Section~\ref{sec:generalized_ohm_4d}).

\subsection{Mixed-sector forces and the physical role of \texorpdfstring{\(A_w\)}{Aw}}
\label{sec:mixed_sector_forces}

The bulk Lorentz force for species \(s\) involves the full \(4\)-spatial tensor
\(F_{AB}\) with \(A,B\in\{x,y,z,w\}\). Splitting the brane and transverse pieces gives
(for brane components \(a\)):
\begin{align}
q_s\Big(E_a + F_{ab}v_s^b + F_{aw}v_s^w\Big)
&=
q_s\Big(E_a + (\mathbf{v}_s\times\mathbf{B})_a + v_s^w C_a\Big),
\label{eq:lorentz_split_brane_beyond}
\\
q_s\Big(E_w + F_{wb}v_s^b\Big)
&=
q_s\Big(E_w - v_s^a C_a\Big),
\label{eq:lorentz_split_w_beyond}
\end{align}
using \(C_a\equiv F_{aw}\) and antisymmetry of \(F_{AB}\).

Equations \eqref{eq:lorentz_split_brane_beyond}--\eqref{eq:lorentz_split_w_beyond} identify
the new physical couplings beyond MHD:
\begin{itemize}[leftmargin=1.7em]
\item a \textbf{brane force channel} \(q_s v_s^w C_a\) that vanishes identically when either
\(v_s^w=0\) (no transverse drift) or \(C_a=0\) (no mixed structure);
\item a \textbf{transverse acceleration channel} driven by \(E_w\), which generically produces
\(v_s^w\) and hence \(J^w\).
\end{itemize}

Thus, suppressing \(A_w\) is not merely a gauge convenience: it removes a dynamical channel
that can exchange energy and momentum between the brane-observed plasma and transverse degrees
of freedom. In our simulation program, \(A_w\) is retained as dynamical except in controlled
limit tests designed to reproduce textbook MHD benchmarks.

\subsection{Finite localization and the transverse-mode tower}
\label{sec:finite_localization_beyond}

When \(Z(w)\) is a Gaussian soft wall, the brane components \(A_\mu\) admit a Hermite mode
expansion in \(w\) (Appendix~\ref{sec:kk_modes}), producing a tower of effective \(3+1\)
fields with masses \(m_n^2=2n/\lambda^2\). The zero mode (\(n=0\)) reproduces standard
Maxwell. The higher modes (\(n\ge 1\)) yield:
\begin{enumerate}[leftmargin=1.6em]
\item \textbf{Yukawa corrections} to static interactions on scales \(r\sim\lambda\);
\item \textbf{new storage/transport reservoirs} for EM energy and helicity, tracked mode-by-mode
in the 3D\(\times\)modes solver;
\item \textbf{a controlled truncation knob} \(N_w\) for massive computation: increasing \(N_w\)
systematically increases transverse fidelity and should converge brane observables and leakage
diagnostics.
\end{enumerate}

A useful practical criterion is therefore the dimensionless ratio
\begin{equation}
\chi_\lambda \equiv \lambda\,|\nabla_3 \ln(\cdot)|,
\label{eq:chi_lambda_def}
\end{equation}
evaluated on the relevant brane structures (current sheets, density gradients, etc.).
When \(\chi_\lambda\ll 1\), the zero-mode approximation is reliable and MHD-like behavior is
expected. When \(\chi_\lambda\gtrsim 1\), higher transverse modes and mixed-sector fields
become dynamically relevant and beyond-MHD corrections should be expected.

\subsection{Generalized Ohm's law and the topology EMF}
\label{sec:generalized_ohm_preview}

The brane induction equation is always obtained by projecting Faraday's law:
\begin{equation}
\partial_t \overline{\mathbf{B}} = -\nabla_3\times \overline{\mathbf{E}}.
\label{eq:proj_faraday_beyond}
\end{equation}
Therefore, ``non-ideality'' on the brane is entirely controlled by how \(\overline{\mathbf{E}}\)
relates to the brane-facing flow variables and currents.

In standard two-fluid \(\to\) MHD reductions, one obtains (schematically)
\(\mathbf{E}+\mathbf{v}\times\mathbf{B}=\eta\mathbf{J}+\cdots\). In the present framework,
two additional structures appear generically when transverse physics is active:

\paragraph{(i) Projection covariance EMF.}
Define the exact covariance EMF
\begin{equation}
\boldsymbol{\mathcal{E}}_{\rm cov}
\;\equiv\;
\overline{\mathbf{v}\times\mathbf{B}}-\overline{\mathbf{v}}\times\overline{\mathbf{B}}
\;=\;
\overline{\mathbf{v}'\times\mathbf{B}'}.
\label{eq:Ecov_def_beyond}
\end{equation}
This term vanishes in the controlled MHD limit (no \(w\)-structure) but need not vanish
even in laminar flows: it measures transverse variation, not turbulence.

\paragraph{(ii) Mixed-sector transverse EMF.}
From the Lorentz-force split \eqref{eq:lorentz_split_brane_beyond}, the electron equation
contains the new term \(v_\ele^w \mathbf{C}\). Its brane-facing contribution is the projected
quantity
\begin{equation}
\boldsymbol{\mathcal{E}}_{w}
\;\equiv\;
\overline{v_\ele^w\,\mathbf{C}},
\qquad \mathbf{C}=(C_x,C_y,C_z),\qquad C_a\equiv F_{aw}.
\label{eq:Ew_def_beyond2}
\end{equation}
This term vanishes if either \(A_w\) is suppressed or transverse drift is absent; otherwise it
provides a direct, computable mechanism for field-line slippage in brane observables.

\paragraph{Topology EMF.}
We group these contributions into a single topology-induced electromotive term
\begin{equation}
\boxed{
\boldsymbol{\mathcal{E}}_{\rm topo}
\;\equiv\;
\boldsymbol{\mathcal{E}}_{\rm cov} + \boldsymbol{\mathcal{E}}_{w}
\;+\;(\text{optional: additional two-fluid/closure pieces}).
}
\label{eq:Etopo_def_beyond}
\end{equation}
Then a brane-facing induction equation can be written in the suggestive form
\begin{equation}
\boxed{
\partial_t \overline{\mathbf{B}}
=
\nabla_3\times\big(\overline{\mathbf{v}}\times\overline{\mathbf{B}}\big)
-\nabla_3\times \boldsymbol{\mathcal{E}}_{\rm topo}
\;+\;\text{(standard Hall/pressure/inertia/resistive terms if retained)}.
}
\label{eq:induction_with_topo_emf_beyond}
\end{equation}

Equation \eqref{eq:induction_with_topo_emf_beyond} is the operational statement of how the
\(4\)D interaction model extends MHD: topology change in brane observables is controlled by a
diagnosable EMF sourced by transverse structure and mixed-sector dynamics. The associated
energy and helicity ledgers, including explicit \(w\)-exchange and covariance reservoirs,
are derived in Appendix~\ref{app:proj_energy_helicity} and form the basis for the validation
tests in Section~\ref{sec:validation}.

% ============================================================
\section{Generalized Ohm's Law with 4D/Projection Corrections}
\label{sec:generalized_ohm_4d}

In standard two-fluid plasma theory, ``non-ideality'' on the brane is captured by a
generalized Ohm law of the schematic form
\(
\mathbf{E}+\mathbf{v}\times\mathbf{B}
=\eta\mathbf{J} + \frac{\mathbf{J}\times\mathbf{B}}{en} - \frac{1}{en}\gradthree p_e + \cdots
\),
and reconnection is associated with a nonzero
\(\mathbf{E}\cdot\mathbf{B}\) in localized regions.
In the \(4+1\)D interaction model, the same logic applies, but two new structures appear
\emph{even in the absence of ad hoc resistivity}:
\begin{enumerate}[leftmargin=1.6em]
\item \textbf{Mixed-sector (transverse) couplings} involving \(A_w\), \(E_w\), and
\(C_a\equiv F_{aw}\), which are absent in strictly \(3+1\) Maxwell theory.
\item \textbf{Projection-induced covariances} (``unresolved transverse structure'') because
the brane observer measures \(\overline{(\cdot)}\) rather than the pointwise bulk field.
\end{enumerate}
The purpose of this section is to write these corrections explicitly in a form that can be
computed in both the direct \(4\)D solver and the 3D\(\times\)modes solver.

\subsection{Bulk electron momentum equation and mixed-sector decomposition}
\label{subsec:ohm_bulk}

Let the electron species have charge \(q_e=-e\) (with \(e>0\)) and bulk number density
\(n_e(t,\mathbf{x},w)\). Denote the bulk spatial velocity by
\begin{equation}
\mathbf{v}_e(t,\mathbf{x},w)\equiv (v_e^x,v_e^y,v_e^z),
\qquad
v_e^w(t,\mathbf{x},w)\equiv \text{transverse component}.
\end{equation}
Define the full \(4\)D convective derivative
\begin{equation}
D_t^{(4)} \equiv \partial_t + v_e^a\partial_a + v_e^w\partial_w,
\qquad a\in\{x,y,z\}.
\label{eq:Dt4_def}
\end{equation}

Using the bulk electromagnetic fields
\(
E_a\equiv F_{a0}
\),
\(
B_a\equiv \tfrac12\epsilon_{abc}F_{bc}
\),
and the mixed field
\(
C_a\equiv F_{aw}=\partial_a A_w-\partial_w A_a
\),
the brane components of the bulk electron momentum equation may be written schematically as
\begin{equation}
m_e n_e\,D_t^{(4)} v_{e,a}
=
-e n_e\Big(E_a + (\mathbf{v}_e\times\mathbf{B})_a + v_e^w C_a\Big)
-\partial_a p_e
+R_{e,a}
+\cdots,
\label{eq:bulk_electron_momentum_brane}
\end{equation}
where \(p_e\) is the electron pressure (or the corresponding enthalpy-gradient term in the
GNLS realization), \(R_{e,a}\) is the electron--ion momentum exchange term (optional), and
\(\cdots\) denotes any additional modeled forces (e.g.\ dispersive/quantum terms if retained).

Solving \eqref{eq:bulk_electron_momentum_brane} pointwise in \(w\) for \(\mathbf{E}\) yields the
\emph{bulk} generalized Ohm law:
\begin{equation}
\mathbf{E}
=
-\mathbf{v}_e\times\mathbf{B}
-\;v_e^w\,\mathbf{C}
-\frac{1}{e n_e}\gradthree p_e
+\frac{1}{e n_e}\mathbf{R}_e
-\frac{m_e}{e}\,D_t^{(4)}\mathbf{v}_e
+\cdots,
\label{eq:ohm_bulk_pointwise}
\end{equation}
where \(\mathbf{C}\equiv (C_x,C_y,C_z)\) and \(\mathbf{R}_e\equiv (R_{e,x},R_{e,y},R_{e,z})\).

\paragraph{New force channel (absent in \(3+1\)).}
The term \(-v_e^w\,\mathbf{C}\) is the direct imprint of the mixed field \(F_{aw}\) on the
brane electric field. It vanishes in the controlled MHD limit
(Section~\ref{sec:mhd_reduction}) but is generically present when \(A_w\) is dynamical and
electrons drift in \(w\).

\subsection{Brane projection and exact projected Ohm form}
\label{subsec:ohm_projected_exact}

Define the operational brane projection (Section~\ref{sec:conventions}) by
\begin{equation}
\overline{Q}(t,\mathbf{x}) \equiv \int_{-\infty}^{+\infty}\dd w\;W(w)\,Q(t,\mathbf{x},w),
\qquad \int\dd w\,W(w)=1.
\label{eq:ohm_proj_def}
\end{equation}
Projecting \eqref{eq:ohm_bulk_pointwise} yields the \emph{exact brane-facing identity}
\begin{equation}
\boxed{
\overline{\mathbf{E}}
=
-\overline{\mathbf{v}_e\times\mathbf{B}}
-\overline{v_e^w\,\mathbf{C}}
-\overline{\frac{1}{e n_e}\gradthree p_e}
+\overline{\frac{1}{e n_e}\mathbf{R}_e}
-\frac{m_e}{e}\,\overline{D_t^{(4)}\mathbf{v}_e}
+\cdots.
}
\label{eq:ohm_projected_exact}
\end{equation}
Equation \eqref{eq:ohm_projected_exact} is not yet in the conventional MHD form because it
contains projected products and projected ratios. This is not a defect: it is precisely the
mathematical statement that the brane observer is not evolving a closed \(3\)D system unless
additional assumptions are imposed.

\subsection{Quasi-neutral two-fluid rewrite and standard MHD terms}
\label{subsec:ohm_twofluid_rewrite_4d}

To connect \eqref{eq:ohm_projected_exact} to the familiar two-fluid/MHD hierarchy, introduce
brane-observed quasi-neutral variables
\begin{equation}
\overline{n}\;\equiv\;\overline{n_e}\simeq \overline{n_i},
\qquad
\overline{\mathbf{J}}\;\equiv\;e\Big(\overline{n_i\mathbf{v}_i}-\overline{n_e\mathbf{v}_e}\Big),
\qquad
\overline{\mathbf{v}}\;\equiv\;\text{(single-fluid brane velocity)}.
\label{eq:ohm_twofluid_vars}
\end{equation}
Under the usual ordering \(m_e\ll m_i\), one has the standard approximation
\begin{equation}
\overline{\mathbf{v}}-\overline{\mathbf{v}_e}
\;\approx\;
\frac{\overline{\mathbf{J}}}{e\,\overline{n}},
\label{eq:v_minus_ve_hall}
\end{equation}
which produces the Hall term when multiplied by \(\overline{\mathbf{B}}\).

Add and subtract \(\overline{\mathbf{v}}\times\overline{\mathbf{B}}\) to obtain
\begin{equation}
\overline{\mathbf{E}}+\overline{\mathbf{v}}\times\overline{\mathbf{B}}
=
\frac{\overline{\mathbf{J}}\times\overline{\mathbf{B}}}{e\,\overline{n}}
-\frac{1}{e\,\overline{n}}\gradthree\overline{p_e}
+\eta\,\overline{\mathbf{J}}
-\frac{m_e}{e}\,\mathcal{I}_e
+\boldsymbol{\mathcal{E}}_{\rm topo}
+\cdots,
\label{eq:ohm_brane_mhd_form}
\end{equation}
where:
\begin{itemize}[leftmargin=1.7em]
\item \(\eta\,\overline{\mathbf{J}}\) is the usual resistive term if a drag model is used
(Appendix~\ref{app:twofluid_to_mhd}); otherwise set \(\eta=0\).
\item \(\mathcal{I}_e\) denotes the \emph{chosen} electron-inertia closure in \(3\)D variables
(e.g.\ \(\mathcal{I}_e\approx \overline{D_t^{(3)}\mathbf{v}_e}\) or an equivalent expression in
\(\overline{\mathbf{J}}\)); its explicit form is not unique and depends on the closure target
(ideal MHD, Hall MHD, extended MHD).
\item \(\boldsymbol{\mathcal{E}}_{\rm topo}\) collects the genuinely new \(4\)D/projection terms,
defined explicitly below.
\end{itemize}

Equation \eqref{eq:ohm_brane_mhd_form} is the primary ``paper equation'': it shows that the brane
system obeys a generalized Ohm law with the standard extended-MHD terms \emph{plus} a new
diagnosable electromotive contribution \(\boldsymbol{\mathcal{E}}_{\rm topo}\).

\subsection{Topology EMF and ``topology stress'' decomposition}
\label{subsec:topology_emf_def}

Define the exact covariance operator (with respect to \(W\)) by
\begin{equation}
\mathrm{Cov}_W[Q,R]\equiv \overline{QR}-\overline{Q}\,\overline{R}.
\label{eq:covW_def_ohm}
\end{equation}
Then we define the topology EMF as
\begin{equation}
\boxed{
\boldsymbol{\mathcal{E}}_{\rm topo}
\;\equiv\;
\boldsymbol{\mathcal{E}}_{\rm cov}
+\boldsymbol{\mathcal{E}}_{w}
+\boldsymbol{\mathcal{E}}_{\rm proj},
}
\label{eq:Etopo_def}
\end{equation}
with the three components:
\begin{align}
\boldsymbol{\mathcal{E}}_{\rm cov}
&\equiv
\overline{\mathbf{v}_e}\times\overline{\mathbf{B}}
-\overline{\mathbf{v}_e\times\mathbf{B}}
=
-\mathrm{Cov}_W[\mathbf{v}_e,\mathbf{B}]_{\times},
\label{eq:Ecov_def_ohm}
\\
\boldsymbol{\mathcal{E}}_{w}
&\equiv
-\overline{v_e^w\,\mathbf{C}},
\qquad \mathbf{C}\equiv(F_{xw},F_{yw},F_{zw}),
\label{eq:Ew_def_ohm}
\\
\boldsymbol{\mathcal{E}}_{\rm proj}
&\equiv
\left(
-\overline{\frac{1}{e n_e}\gradthree p_e}
+\frac{1}{e\,\overline{n}}\gradthree\overline{p_e}
\right)
+
\left(
\overline{\frac{1}{e n_e}\mathbf{R}_e}
-\eta\,\overline{\mathbf{J}}
\right)
-\frac{m_e}{e}\left(\overline{D_t^{(4)}\mathbf{v}_e}-\mathcal{I}_e\right)
+\cdots.
\label{eq:Eproj_def_ohm}
\end{align}

\paragraph{Interpretation.}
\(\boldsymbol{\mathcal{E}}_{\rm topo}\) is the part of \(\overline{\mathbf{E}}+\overline{\mathbf{v}}\times\overline{\mathbf{B}}\)
that is \emph{forced} by transverse physics (mixed fields and \(w\)-structure) and by the fact that
projection does not commute with nonlinear operations. It vanishes in the controlled MHD limit
(Section~\ref{sec:mhd_reduction}), but it need not vanish even in laminar flows if there is transverse
structure.

\paragraph{Why this acts like an ``effective resistivity'' without dissipation.}
In resistive MHD, non-ideal behavior is often modeled as \(\mathbf{E}+\mathbf{v}\times\mathbf{B}\approx\eta\mathbf{J}\).
Here, the term \(\boldsymbol{\mathcal{E}}_{w}=-\overline{v_e^w\mathbf{C}}\) has the same structural role in the induction
equation (it contributes through a curl), but it corresponds to \emph{momentum/energy export into transverse degrees of freedom}
rather than collisional heating. This distinction is made quantitative by the energy and helicity ledgers with explicit leakage
(Appendix~\ref{app:proj_energy_helicity}).

\subsection{Induction equation and reconnection diagnostics}
\label{subsec:induction_topo}

Projecting Faraday's law gives
\begin{equation}
\partial_t\overline{\mathbf{B}} = -\gradthree\times \overline{\mathbf{E}}.
\end{equation}
Using \eqref{eq:ohm_brane_mhd_form}, the brane induction equation becomes
\begin{equation}
\boxed{
\partial_t\overline{\mathbf{B}}
=
\gradthree\times(\overline{\mathbf{v}}\times\overline{\mathbf{B}})
-\gradthree\times\left[
\frac{\overline{\mathbf{J}}\times\overline{\mathbf{B}}}{e\,\overline{n}}
-\frac{1}{e\,\overline{n}}\gradthree\overline{p_e}
+\eta\,\overline{\mathbf{J}}
-\frac{m_e}{e}\mathcal{I}_e
+\boldsymbol{\mathcal{E}}_{\rm topo}
\right]
+\cdots,
}
\label{eq:induction_with_topo}
\end{equation}
with \(\gradthree\cdot\overline{\mathbf{B}}=0\) maintained by the Maxwell solver (or by the chosen divergence-control scheme).

A key reconnection/topology diagnostic is the component of the non-ideal electric field parallel to the magnetic field,
\begin{equation}
\overline{E}_\parallel
\equiv
\frac{\overline{\mathbf{E}}\cdot\overline{\mathbf{B}}}{|\overline{\mathbf{B}}|},
\qquad
\overline{E}_\parallel
=
\left[
\eta\,\overline{\mathbf{J}}
-\frac{m_e}{e}\mathcal{I}_e
+\boldsymbol{\mathcal{E}}_{\rm topo}
+\cdots
\right]\cdot\frac{\overline{\mathbf{B}}}{|\overline{\mathbf{B}}|},
\label{eq:Eparallel_diag}
\end{equation}
showing explicitly that a nonzero \(\overline{E}_\parallel\) can arise from the topology EMF even when \(\eta=0\).
The associated energy/helicity accounting is handled by the projected budgets in Appendix~\ref{app:proj_energy_helicity}.

\subsection{How \texorpdfstring{$\boldsymbol{\mathcal{E}}_{\rm topo}$}{E topo} is computed in simulations}
\label{subsec:etopocompute}

The decomposition \eqref{eq:Etopo_def} is designed to be \emph{computationally explicit}:

\begin{itemize}[leftmargin=1.7em]
\item \textbf{Direct \(4\)D-grid solver.}
All quantities in \eqref{eq:Ecov_def_ohm}--\eqref{eq:Eproj_def_ohm} are evaluated on the
\((x,y,z,w)\) grid and then integrated against \(W(w)\) to form \(\overline{(\cdot)}\).

\item \textbf{3D\(\times\)modes solver (Hermite / soft-wall).}
Reconstruct the needed fields at transverse quadrature nodes \(w_q\) from the retained mode amplitudes
(Appendix~\ref{sec:kk_modes}), form the nonlinear products \(\mathbf{v}_e\times\mathbf{B}\) and \(v_e^w\mathbf{C}\) pointwise in \(w_q\),
and then apply quadrature for \(\overline{(\cdot)}\) and for the covariance differences. Mode-truncation convergence is assessed by increasing
\(N_w\) until \(\boldsymbol{\mathcal{E}}_{\rm topo}\) and the reconnection/helicity diagnostics stabilize (Section~\ref{sec:validation}).
\end{itemize}

In both architectures, \(\boldsymbol{\mathcal{E}}_{\rm topo}\) is therefore not an \emph{assumed closure} but a directly measured
quantity arising from \(4\)D interaction structure and operational projection.

% ============================================================
\section{Computation: How We Will Do Massive Plasma Simulations}
\label{sec:computation}

This section translates the model hierarchy into a concrete numerical program. The
primary objective is to enable \emph{massive} plasma simulations that retain the
essential \(4\)D interaction channels (transverse transport in \(w\), scalar-photon
dynamics \(A_w\), and transverse-mode excitation) while remaining computationally
tractable on modern HPC architectures. We therefore present two solver families:

\begin{enumerate}[leftmargin=1.6em]
\item \textbf{Direct 4D-grid solvers} that discretize \((x,y,z,w)\) explicitly. These
are the most direct representation of the bulk PDEs and serve as ``ground truth''
and validation for reduced models.

\item \textbf{3D\(\times\)modes solvers} that discretize \((x,y,z)\) on a large grid
and represent \(w\)-dependence in a truncated spectral basis (Hermite modes for
Gaussian localization). This is the recommended approach for truly massive runs.
\end{enumerate}

Throughout, we avoid evolving only brane variables: the computational strategy is
to evolve bulk degrees of freedom and then compute brane observables by the
projection map \(\Proj{\cdot}\), including explicit leakage diagnostics.

\subsection{State variables and discretization choices}
\label{sec:state_vars}

\subsubsection*{Electromagnetic variables}

We evolve the gauge potential \(\A_M(t,\x,w)\) with a hyperbolic gauge choice
(e.g.\ generalized Lorenz gauge as in \eqref{eq:maxwell_full}) so that the field
equations form a well-posed time-domain system. Evolving \(\A_M\) has two practical
advantages:
\begin{enumerate}[leftmargin=1.6em]
\item the homogeneous Maxwell equations (Bianchi identity) are satisfied by
construction;
\item the mixed sector \((A_w,\,F_{\mu w})\) is handled on the same footing as the
brane sector without additional constraint equations.
\end{enumerate}
An alternative (useful for some plasma codes) is to evolve field strengths
\((E_A,\,F_{AB})\) with explicit constraint cleaning. In this paper we prioritize
the potential formulation because the transverse-mode decomposition (Sec.~\ref{sec:kk_modes})
is naturally expressed in terms of \(\A_M\).

\subsubsection*{Plasma variables}

We support three plasma representations, ordered by increasing kinetic fidelity:

\paragraph{(i) Two-fluid in 4D space.}
Evolve \((n_s,\mathbf{v}_s)\) on the bulk spatial domain \((\x,w)\), with
\(\mathbf{v}_s\in\R^4\), and close pressure by a specified EOS. This is the
closest analogue to extended MHD, but in \(4\) spatial dimensions.

\paragraph{(ii) Field-based (multi-species GNLS).}
Evolve \(\psi_s(t,\x,w)\) for each species and compute \(\rho_s=|\psi_s|^2\) and
\(\Js^M\) from \eqref{eq:number_current_full}--\eqref{eq:charge_current_full}.
This approach evolves canonical momentum naturally (Sec.~\ref{sec:canonical_momentum})
and remains well-defined at low density, which is attractive for topology change.

\paragraph{(iii) Kinetic (4D Vlasov / PIC).}
Evolve particle distributions in \(4\)D space with \(4\)D velocity. This is the
most expensive option but provides the clearest path to collisionless reconnection
physics in regimes where fluid closures are insufficient.

\subsubsection*{Spatial discretization in \texorpdfstring{$(x,y,z)$}{(x,y,z)}}

For brane spatial directions we assume a large uniform grid for simplicity and
HPC efficiency, though adaptive mesh refinement is compatible with the mode-based
approach. We recommend:
\begin{itemize}[leftmargin=1.7em]
\item \textbf{finite differences} (2nd--4th order) for Maxwell in potentials;
\item \textbf{finite volume} (Godunov-type) for two-fluid bulk equations when shocks
or steep gradients appear;
\item \textbf{spectral/FFT Poisson solves} only in regimes where a reduced Maxwell
closure is used (e.g.\ Darwin/MHD ordering), not in the full wave evolution.
\end{itemize}

\subsubsection*{Representation of the extra dimension \texorpdfstring{$w$}{w}}

Two choices:

\paragraph{Direct \(w\)-grid (baseline).}
Use a finite interval \(w\in[-W_{\max},W_{\max}]\) with boundary conditions chosen
to respect localization (e.g.\ outgoing/radiative boundaries). This is simple but
can be expensive because stable timesteps are constrained by the finest \(\Delta w\).

\paragraph{Hermite spectral in \(w\) (recommended).}
When \(\Z(w)\) is Gaussian, the transverse operator admits a Hermite basis with a
diagonal mass spectrum \(m_n^2\propto n\) (Sec.~\ref{sec:kk_modes}). A truncated
mode expansion converts the problem into a set of coupled \(3+1\) PDEs for mode
amplitudes. This is the primary pathway to massive simulations.

\subsection{Hermite spectral decomposition in \texorpdfstring{$w$}{w}}
\label{sec:hermite_scheme}

\subsubsection*{Field expansion and truncation}

For each component of the gauge potential (and, when desired, for matter fields),
we use a truncated expansion
\begin{equation}
\A_M(t,\x,w) \;\approx\; \sum_{n=0}^{N_w-1} a_M^{(n)}(t,\x)\,\phin(w),
\label{eq:Aw_mode_expand}
\end{equation}
where \(\{\phin(w)\}\) are orthonormal with respect to the \(\Z\)-weighted inner
product\\
\(\langle f,g\rangle_{\Z}=\int \dd w\,\Z(w)f(w)g(w)\). For Gaussian \(\Z\),
a Hermite basis yields a simple spectrum
\begin{equation}
m_n^2 = \frac{2n}{\lambda^2},
\qquad n=0,1,2,\ldots,
\label{eq:mn_spectrum_repeat}
\end{equation}
and brane coupling selection rules (odd modes decouple from strictly brane-localized
sources).

The truncation \(N_w\) is a \emph{physical} knob: increasing \(N_w\) increases the
available transverse degrees of freedom for topology and energy to flow into/out
of the brane-observed sector. Convergence is assessed by monitoring the mode-energy
spectrum and leakage diagnostics (Sec.~\ref{sec:diagnostics}).

\subsubsection*{Linear operators are diagonal in mode index}

Under standard gauge conditions (e.g.\ Lorenz gauge), the linear Maxwell operator
reduces mode-by-mode to brane wave operators of the form
\begin{equation}
\left(\square_4 + m_n^2\right)a_\mu^{(n)}(t,\x)
= \text{(mode-projected sources)},
\qquad
\square_4\equiv -\partial_t^2+\lapthree.
\label{eq:mode_linear_operator}
\end{equation}
Thus the expensive part of the solver (the linear field propagation) becomes a
collection of independent, stencil-based \(3\)D wave updates indexed by \(n\),
which is ideal for vectorization, GPUs, and mode-parallel MPI.

\subsubsection*{Nonlinear couplings handled pseudo-spectrally}

Nonlinearities enter through sources \(\J^M(\psi_s,\A)\) and/or through fluid
nonlinear terms. Rather than deriving dense mode-coupling tensors in closed form,
we use a pseudo-spectral strategy in \(w\):
\begin{enumerate}[leftmargin=1.6em]
\item choose quadrature points \(\{w_q\}_{q=1}^{Q_w}\) and weights \(\{\omega_q\}\)
appropriate to the \(\Z\)-weighted inner product;
\item reconstruct fields at quadrature nodes:
\(\A_M(t,\x,w_q)=\sum_{n}a_M^{(n)}(t,\x)\phin(w_q)\) and similarly for matter variables;
\item compute nonlinear source terms pointwise at each \(w_q\);
\item project back to modes by quadrature:
\begin{equation}
J^{M,(n)}(t,\x)\;\approx\;\sum_{q=1}^{Q_w}\omega_q\,\Z(w_q)\,\phin(w_q)\,J^M(t,\x,w_q).
\end{equation}
\end{enumerate}
This keeps the mode truncation controllable and avoids introducing spurious
``hard-wall'' artifacts in \(w\).

\subsection{Coupling algorithm (field solve + plasma push)}
\label{sec:coupling_algorithm}

We use a standard split-step / operator-splitting approach tailored to the chosen
plasma model. A typical timestep from \(t^n\to t^{n+1}=t^n+\Delta t\) proceeds as:

\begin{enumerate}[leftmargin=1.6em]
\item \textbf{Reconstruct fields in \(w\) (mode solver only):}
build \(\A_M(t^n,\x,w_q)\) at quadrature nodes \(w_q\) if using the 3D\(\times\)modes scheme.

\item \textbf{Advance plasma (matter step):}
\begin{itemize}[leftmargin=1.7em]
\item \emph{Two-fluid:} update \((n_s,\mathbf{v}_s)\) with a conservative finite-volume
scheme in \((x,y,z)\) for each \(w_q\) or mode coefficient, including Lorentz force
terms using \(\F_{MN}\).
\item \emph{GNLS:} apply a Strang split: kinetic covariant step
\(-(\hbar^2/2m_s)D_A D_A\), then local nonlinear/potential step
\((V_{{\rm conf},s}+h_s)\), with gauge fields held fixed within substeps.
\item \emph{PIC/Vlasov:} push particles in \(4\)D space using the Lorentz tensor,
deposit charge/current with a charge-conserving scheme to guarantee discrete
\(\partial_M J^M=0\).
\end{itemize}

\item \textbf{Compute sources:}
evaluate \(\J^M\) from updated matter variables. In mode form this yields
\(\{J^{M,(n)}(t^{n+\frac12},\x)\}\) by projection (pseudo-spectral step).

\item \textbf{Advance Maxwell (field step):}
update the localized Maxwell system in potentials. In 3D\(\times\)modes form, each
mode is updated with a wave-like stencil including the mass term \(m_n^2\). In the
direct 4D-grid form, update \(\A_M(t,\x,w)\) with the variable-coefficient operator
\(\partial_M(\Z F^{MN})\).

\item \textbf{Constraint/gauge control:}
apply gauge-driver or constraint-damping steps (e.g.\ generalized Lorenz gauge
damping) to suppress numerical drift in \(\partial\!\cdot\!\A\). If evolving fields
directly, apply divergence cleaning to maintain Gauss constraints.

\item \textbf{Diagnostics and brane observables:}
compute brane-projected fields \(\Proj{F_{\mu\nu}}\), brane currents, and leakage
terms \(\Sleak^{(s)}\). Record energy budgets and mode spectra (Sec.~\ref{sec:diagnostics}).
\end{enumerate}

\paragraph{Time stepping and stiffness.}
Full Maxwell evolution carries the light-wave CFL condition. For MHD-like regimes
where fast light waves are not of primary interest, one may employ IMEX time
stepping (implicit Maxwell, explicit plasma) or reduced-field closures (Darwin /
low-frequency approximations) while retaining the transverse-mode structure in \(w\).
The present paper focuses on the full hyperbolic system and treats reduced closures
as optional accelerations once correctness is established.

\subsection{Parallelization and scaling considerations}
\label{sec:hpc_scaling}

\subsubsection*{Cost model}

Let \(N_xN_yN_z\) be the brane grid size, and let \(N_w\) be either the number of
\(w\)-grid points (direct 4D) or the number of transverse modes (3D\(\times\)modes).
Then, schematically:
\begin{itemize}[leftmargin=1.7em]
\item \textbf{Direct 4D grid:} memory and compute scale as
\(\mathcal{O}(N_xN_yN_zN_w)\), with Maxwell CFL constrained by \(\min(\Delta x,\Delta w)\).
\item \textbf{3D\(\times\)modes:} memory and compute scale as
\(\mathcal{O}(N_xN_yN_zN_w)\) as well, but with substantially smaller prefactors because:
(i) the \(w\)-operator is diagonal in mode index for the linear step, and
(ii) the reconstruction/projection in \(w\) is a dense operation only over a small
\(N_w\) (tens), not over a large \(w\)-grid (hundreds or thousands).
\end{itemize}

In practice, massive runs target \(N_w\sim 8\) to \(64\) modes, chosen by convergence
of the mode-energy spectrum and leakage diagnostics.

\subsubsection*{Domain decomposition and mode parallelism}

We use two independent decomposition axes:
\begin{enumerate}[leftmargin=1.6em]
\item \textbf{Spatial decomposition in \((x,y,z)\):} standard MPI halo exchange for
stencil updates of Maxwell and fluid variables.
\item \textbf{Mode decomposition in \(n\):} distribute subsets of mode indices across
ranks (or GPU streams). The linear Maxwell update is embarrassingly parallel in \(n\).
Nonlinear reconstruction/projection requires only local (per-cell) operations and
therefore minimal communication.
\end{enumerate}
A hybrid layout (space decomposition outer, mode decomposition inner) enables both
strong scaling (more ranks for fixed problem) and weak scaling (larger grids with
fixed per-rank memory).

\subsubsection*{GPU suitability}

The 3D\(\times\)modes representation is particularly GPU-friendly because each cell
carries short vectors of mode coefficients \(\{a_M^{(n)}\}_{n=0}^{N_w-1}\) and the
dominant kernels are fused per-cell updates:
\begin{itemize}[leftmargin=1.7em]
\item mode-local Maxwell stencils (vector length \(N_w\)),
\item reconstruction/projection in \(w\) (small dense matrix multiply per cell),
\item current construction for each species.
\end{itemize}

\subsection{Diagnostics for reconnection and topology transport}
\label{sec:diagnostics}

A central claim of this framework is that brane-observed ``non-ideality'' can arise
from conservative transport into transverse degrees of freedom rather than from
ad hoc resistivity. Simulations must therefore record diagnostics that separate:
(i) true dissipation (if collisions are included) from (ii) topology/energy export
into \(w\) and into higher transverse modes.

\subsubsection*{Bulk and mode-resolved energy budgets}

We track:
\begin{itemize}[leftmargin=1.7em]
\item \textbf{Total bulk energy} \(\mathcal{E}_{\rm tot}(t)\) from
\eqref{eq:E_tot_def} as a global correctness check.
\item \textbf{Mode energies} \(E_{\rm EM}^{(n)}(t)\) in the 3D\(\times\)modes solver,
computed from the quadratic forms induced by the action (including mass terms
\(m_n^2\)). A growth of \(n\ge 1\) energy quantifies transverse-mode excitation.
\item \textbf{Scalar-photon channel energy} associated with \(A_w\) and \(E_w=F_{w0}\),
including the work term \(J^w E_w\) in the exchange identity \eqref{eq:energy_exchange}.
\end{itemize}

\subsubsection*{Leakage (open-system) diagnostics}

For each species \(s\), we compute the projected leakage source \(\Sleak^{(s)}\)
from \eqref{eq:Sleak_def}. In addition we record bulk transverse fluxes:
\begin{equation}
\Phi_{s}^{(w)}(t;\x) \equiv j_s^w(t,\x,w),
\qquad
\Proj{\Phi_{s}^{(w)}}(t,\x)=\int \dd w\,\W(w)\,j_s^w(t,\x,w),
\end{equation}
as well as the integrated net exchange over a region \(\Omega\subset\R^3\):
\begin{equation}
\mathcal{L}_s(t;\Omega) \equiv \int_{\Omega}\Sleak^{(s)}(t,\x)\,\dd^3x.
\end{equation}
Nonzero \(\mathcal{L}_s\) directly measures brane--bulk exchange at the operational
measurement scale.

\subsubsection*{Reconnection-rate proxies (brane-facing)}

On the brane we compute standard reconnection indicators from brane-projected fields
(e.g.\ \(\mathbf{E}^{\rm (brane)}\), \(\mathbf{B}^{\rm (brane)}\)) while interpreting
their non-ideal parts through bulk channels:
\begin{itemize}[leftmargin=1.7em]
\item \textbf{Parallel electric field proxy:}
\(
E_{\parallel}^{\rm (brane)} \equiv
\mathbf{E}^{\rm (brane)}\cdot \mathbf{B}^{\rm (brane)} / |\mathbf{B}^{\rm (brane)}|
\),
and its line integral through current sheets.
\item \textbf{Brane helicity change:} compute magnetic helicity
\(H_{\rm brane}=\int \dd^3x\,\mathbf{A}^{\rm (brane)}\cdot\mathbf{B}^{\rm (brane)}\)
in a chosen brane gauge and track \(\dd H_{\rm brane}/\dd t\). Changes in
\(H_{\rm brane}\) are compared against energy transfer into \(A_w\), \(E_w\), and
higher modes to distinguish topology export from dissipation.
\item \textbf{Topology-stress signature:} evaluate projected unresolved-stress terms
(e.g.\ the ``topology stress tensor'' defined later in Section~\ref{sec:generalized_ohm_4d})
to quantify how sub-\(w\) structure drives effective non-ideal terms in brane balances.
\end{itemize}

\paragraph{Convergence criteria specific to this framework.}
A simulation is regarded as ``4D-converged'' not only when brane-facing fields are
grid-converged in \((x,y,z)\), but also when:
(i) leakage diagnostics stabilize under increases in \(N_w\),
(ii) mode-energy spectra show a resolved cascade or decay with \(n\),
and (iii) the bulk energy ledger closes to within the expected numerical tolerance.
These criteria are essential because reconnection-like behavior may manifest as
transport into higher transverse modes rather than as local heating.

% ============================================================
\section{Validation Plan and Benchmark Problems}
\label{sec:validation}

This section outlines a validation program designed to establish correctness at
three distinct levels:
\begin{enumerate}[leftmargin=1.6em]
\item \textbf{Bulk PDE correctness:} the \(4+1\) Maxwell--matter system is solved
accurately, respects discrete charge conservation, satisfies the homogeneous Maxwell
identities, and closes the bulk energy ledger.
\item \textbf{Controlled-limit correctness:} when the controlled assumptions in
Section~\ref{sec:mhd_reduction} are enforced (zero-mode dominance, negligible
\(F_{\mu w}\) and \(J^w\), quasi-neutral two-fluid ordering), the simulation reduces
to standard \(3\)D Maxwell/two-fluid/MHD benchmark behavior.
\item \textbf{Beyond-MHD signal isolation:} when \(4\)D interaction channels are
activated (\(A_w\), \(F_{\mu w}\), \(J^w\), transverse-mode tower), their effects appear
in the predicted observables (leakage sources, mode-energy transfer, topology-stress
signatures) while the bulk system remains conservative.
\end{enumerate}

We explicitly validate both computational architectures described in
Section~\ref{sec:computation}:
\textbf{direct \(4\)D-grid} evolution in \((x,y,z,w)\) and the \textbf{3D\(\times\)modes}
Hermite spectral reduction in \(w\). A major goal is to show that the mode solver
converges to the direct solver as \(N_w\) is increased.

\subsection{MHD benchmarks (controlled limit)}
\label{sec:validation_mhd}

All tests in this subsection are run in a controlled-limit configuration:
\begin{equation}
\partial_w(\cdot)\approx 0,\qquad A_w\approx 0,\qquad F_{\mu w}\approx 0,\qquad J^w\approx 0,
\qquad N_w=1 \ \text{(mode solver)}.
\label{eq:controlled_validation_config}
\end{equation}
In this configuration, brane projection becomes trivial (\(\Proj{Q}\approx Q\)) and the
effective coupling \(\mu_0\to\mu_0^{\rm eff}\) is fixed by \(\Zint\) as in
\eqref{eq:mu0eff_repeat}.

\subsubsection*{(A) Maxwell-sector verification in \(3+1\)}
\begin{enumerate}[leftmargin=1.6em]
\item \textbf{Vacuum plane waves.}
Initialize a vacuum wave packet in Lorenz gauge and verify:
(i) wave speed \(c\) (or \(1\) in units),
(ii) preservation of the gauge condition,
(iii) convergence rates under grid refinement, and
(iv) closure of electromagnetic energy to truncation error.

\item \textbf{Static Coulomb field.}
Place a stationary brane charge distribution \(\rho_q(\x)\) and verify the static
solution \(\E=-\gradthree \phi\) reproduces the Poisson limit of
\eqref{eq:maxwell_inhom_brane} (with \(\epsilon_0^{\rm eff}\)). This test is
useful for verifying charge deposition, Gauss constraints, and solver correctness.

\item \textbf{Charge conservation under a moving source.}
Evolve a prescribed moving brane charge/current that is analytically conserved and
verify discrete \(\partial_t\rho_q+\gradthree\cdot\mathbf{J}=0\) and stability of
Gauss constraints.
\end{enumerate}

\subsubsection*{(B) Two-fluid and ideal-MHD wave benchmarks}
\begin{enumerate}[leftmargin=1.6em]
\item \textbf{Circularly polarized Alfv\'en wave.}
Initialize a nonlinear Alfv\'en wave (a standard code verification target) and
verify amplitude preservation, phase speed, and convergence.

\item \textbf{Fast/slow magnetosonic waves.}
Initialize linear perturbations about a uniform background and verify dispersion
relations in the two-fluid solver (and in the ideal-MHD closure when invoked).

\item \textbf{Whistler/Hall branch (optional extended-MHD check).}
If the generalized Ohm law \eqref{eq:ohm_general} is used with the Hall term retained,
verify the whistler dispersion relation. This provides a sensitive check of
\(\mathbf{J}\times\mathbf{B}\) coupling and current evaluation.
\end{enumerate}

\subsubsection*{(C) Standard 3D MHD code benchmarks}
These tests validate shock capturing, nonlinear coupling, and topology handling in
the controlled limit:
\begin{enumerate}[leftmargin=1.6em]
\item \textbf{Brio--Wu shock tube (1D MHD Riemann problem).}
Verify shock/rarefaction structures and convergence to a reference solution.

\item \textbf{Orszag--Tang vortex (2D).}
A standard benchmark for transition to MHD turbulence. Validate global invariants,
magnetic/kinetic energy evolution, and compare with reference MHD results.

\item \textbf{Field-loop advection (2D/3D).}
Advect a weak magnetic loop in a prescribed velocity field and verify minimal
numerical diffusion and preservation of \(\gradthree\cdot\B=0\).

\item \textbf{Kelvin--Helmholtz instability (2D/3D).}
Validate linear growth rates (early time) and nonlinear roll-up (late time), with
attention to divergence control and energy conservation.
\end{enumerate}

\subsubsection*{(D) Controlled reconnection tests (resistive and/or Hall MHD)}
To compare with known reconnection scalings in the controlled limit:
\begin{enumerate}[leftmargin=1.6em]
\item \textbf{Harris current sheet (2D).}
Run tearing-mode onset with either explicit resistivity \(\eta\) or Hall terms.
Measure reconnection rate versus \(\eta\) (Sweet--Parker scaling) and/or Hall length.

\item \textbf{GEM reconnection challenge configuration (optional).}
If kinetic or two-fluid physics is active, compare reconnected flux versus time and
outflow speeds to the canonical GEM setup as a reference behavior target.
\end{enumerate}

\paragraph{Acceptance criteria for controlled-limit benchmarks.}
A controlled-limit validation pass requires:
\begin{enumerate}[leftmargin=1.6em]
\item correct linear dispersion (waves) within resolution error;
\item correct shock structure (Riemann problems) and convergence with refinement;
\item \(\|\gradthree\cdot\B\|\) bounded and decreasing with refinement (or maintained
to machine tolerance with constrained transport);
\item agreement with standard MHD reference outputs to within discretization error.
\end{enumerate}

\subsection{Beyond-MHD tests (4D effects)}
\label{sec:validation_beyond}

The objective here is not to reproduce standard MHD results, but to validate that
the \emph{new} degrees of freedom behave as predicted, and that their brane-facing
signatures can be isolated cleanly while the bulk system remains conservative.
Unless otherwise stated, these tests are performed in the full \(4+1\) formulation
(direct \(w\)-grid or \(3\)D\(\times\)modes with \(N_w>1\)) and brane observables are
computed via \(\Proj{\cdot}\).

\subsubsection*{(A) Scalar-photon channel: free evolution and coupling}
\begin{enumerate}[leftmargin=1.6em]
\item \textbf{Vacuum \(A_w\) pulse propagation.}
Initialize \(A_w(t=0,\x,w)\) as a localized pulse (with all sources off) and verify:
(i) stable propagation under the localized Maxwell operator,
(ii) correct energy partition in the mixed sector via \(E_w=F_{w0}\) and
\(C_a=F_{aw}\), and
(iii) gauge-driver stability (no spurious growth of \(\partial\!\cdot\!A\)).

\item \textbf{Matter response to \(E_w\).}
Initialize a species density localized near the brane and apply a controlled \(A_w\)
configuration that produces a nonzero \(E_w\). Verify that a transverse drift
\(v_{s,w}\) develops in accordance with the bulk momentum equation
\eqref{eq:euler_species_w}, and that the resulting \(j_s^w\) generates a measurable
brane leakage \(\Sleak^{(s)}\) through \eqref{eq:Sleak_def}.
\end{enumerate}

\subsubsection*{(B) Leakage identity tests (open-system brane balance laws)}
These tests validate that the brane is an open subsystem and that the leakage terms
are quantitatively correct:
\begin{enumerate}[leftmargin=1.6em]
\item \textbf{Continuity leakage closure.}
Construct an initial condition with a known, imposed \(j_s^w(t,\x,w)\) (e.g.\ a
transverse pulse) and verify that the measured brane density obeys
\begin{equation}
\partial_t \rho_{s,\rm brane} + \partial_a j_{s,\rm brane}^a = \Sleak^{(s)}
\end{equation}
to numerical accuracy. This is an exact identity test (discrete verification of
\eqref{eq:projected_continuity_leakage}).

\item \textbf{Bulk conservation with nonzero brane leakage.}
Simultaneously verify that the bulk continuity equation remains satisfied while the
brane subsystem exhibits net exchange \(\mathcal{L}_s(t;\Omega)\neq 0\) for an
appropriate brane region \(\Omega\). This demonstrates that brane non-ideality is
consistent with conservative bulk evolution.
\end{enumerate}

\subsubsection*{(C) Finite-localization and mode-tower tests}
These tests validate the Hermite/KK tower structure and its physical consequences:
\begin{enumerate}[leftmargin=1.6em]
\item \textbf{Static point-source potential (Coulomb + Yukawa tower).}
For a brane-localized static charge, compute the brane potential \(A_0(r)\) and
verify agreement with the predicted Coulomb-plus-Yukawa form
(e.g.\ \eqref{eq:leading_yukawa_em}) in the regimes \(r\gg\lambda\) and \(r\sim\lambda\).
This directly validates mode masses \(m_n\) and coupling weights \(c_n\).

\item \textbf{Mode excitation spectrum under sharp gradients.}
Drive the system with a brane-localized current sheet whose thickness is varied
across \(\sim\lambda\). Verify that:
(i) the \(n=0\) mode dominates for broad structures,
(ii) higher modes are excited as gradients approach \(\lambda^{-1}\),
and (iii) the mode-energy spectrum converges as \(N_w\) increases.
\end{enumerate}

\subsubsection*{(D) Reconnection via topology export (key beyond-MHD demonstration)}
The flagship beyond-MHD validation is to show that reconnection-like topology change
in brane observables can occur through conservative \(4\)D transport channels:
\begin{enumerate}[leftmargin=1.6em]
\item \textbf{Current-sheet reconnection with suppressed resistivity.}
Initialize a brane current sheet in a configuration that would require explicit
resistivity (or numerical diffusion) to reconnect in ideal MHD. Run two variants:
\begin{itemize}[leftmargin=1.7em]
\item \emph{Controlled-limit run:} enforce \(A_w\approx 0\), \(J^w\approx 0\),
\(N_w=1\). Reconnection should be suppressed up to numerical diffusion.
\item \emph{Full 4D-interaction run:} allow \(A_w\), \(F_{\mu w}\), \(J^w\), and/or
\(N_w>1\). Measure enhanced reconnection signatures on the brane.
\end{itemize}
The validation target is not a single reconnection rate, but a \emph{ledger closure}:
the increase in brane reconnected flux must correlate with:
(i) nonzero transverse exchange (\(\Sleak\), \(J^wE_w\)),
(ii) growth of mixed-sector energy (via \(E_w\), \(F_{aw}\)),
and/or (iii) transfer of EM energy into higher transverse modes.

\item \textbf{Topology-stress correlation.}
Compute the projected unresolved-stress / ``topology stress tensor'' terms defined
in Section~\ref{sec:generalized_ohm_4d} and verify that their support localizes near
reconnection regions and correlates with the non-ideal term in the brane induction
equation inferred from \(\partial_t\B+\gradthree\times(\mathbf{v}\times\B)\).
\end{enumerate}

\paragraph{Acceptance criteria for beyond-MHD tests.}
A beyond-MHD validation pass requires:
\begin{enumerate}[leftmargin=1.6em]
\item exact projected identities (e.g.\ leakage balance) hold discretely to solver tolerance;
\item bulk global energy \(\mathcal{E}_{\rm tot}(t)\) remains conserved up to expected
numerical error when collisions are off;
\item mode spectra and leakage diagnostics converge as \(N_w\) and \(Q_w\) increase;
\item reconnection-like brane behavior is accompanied by quantifiable transport into
\(w\)/modes (not solely by numerical diffusion).
\end{enumerate}

\subsection{Convergence and truncation studies}
\label{sec:validation_convergence}

To support massive simulations, we require convergence studies across four axes:
\begin{enumerate}[leftmargin=1.6em]
\item \textbf{Brane spatial resolution:} refine \(\Delta x\) in \((x,y,z)\) and verify
expected convergence rates for smooth problems and correct shock capture for
discontinuous problems.
\item \textbf{Time step:} refine \(\Delta t\) and verify stability and convergence,
including sensitivity to gauge-driver parameters.
\item \textbf{Transverse resolution:} increase either \(W_{\max}\) and \(\Delta w\)
(direct solver) or \(N_w\) and quadrature order \(Q_w\) (mode solver) until brane
observables and diagnostics converge.
\item \textbf{Solver consistency:} demonstrate that the 3D\(\times\)modes solver
converges to the direct 4D-grid solver for shared test problems as \(N_w\) is
increased and/or as the direct \(w\)-grid is refined.
\end{enumerate}

\paragraph{Recommended convergence reporting.}
For each benchmark we report:
(i) \(\|\gradthree\cdot\B\|\) norms, (ii) bulk energy closure error,
(iii) brane leakage norm \(\|\Sleak\|\) and integrated exchange \(\mathcal{L}_s\),
(iv) mode-energy spectra \(E_{\rm EM}^{(n)}\), and (v) brane reconnection proxies
(when applicable). This provides a unified basis for distinguishing true physical
beyond-MHD behavior from truncation artifacts.

% ============================================================
\section{Discussion}
\label{sec:discussion}

This paper advances a ``4D-first'' viewpoint for plasma modeling: the fundamental
dynamics lives in a conservative bulk with four spatial dimensions \((x,y,z,w)\),
and the familiar \(3\)D plasma variables arise as \emph{brane-facing observables}
defined by a projection map \(\Proj{\cdot}\). In this setting, standard MHD is not
replaced; it is recovered as a controlled limit. The primary payoff is that the
mechanisms normally inserted into MHD as phenomenological non-idealities acquire a
geometric and dynamical origin: topology change on the brane can correspond to
conservative transport into transverse degrees of freedom in \(w\) and into the
transverse mode tower induced by localization.

\subsection{MHD as a controlled brane/zero-mode limit}
Section~\ref{sec:mhd_reduction} makes explicit that recovering MHD requires more than
taking a low-frequency limit. One must additionally suppress the \emph{degrees of freedom
that are absent in strictly \(3+1\) electromagnetism}:
\begin{equation}
\partial_w(\cdot)\approx 0,
\qquad
F_{\mu w}\approx 0,
\qquad
J^w\approx 0,
\label{eq:discussion_mhd_assumptions}
\end{equation}
together with the standard two-fluid \(\to\) MHD assumptions (quasi-neutrality and
electron inertia suppression). In this regime, the localized Maxwell sector reduces to
a \(3+1\) Maxwell system with effective coupling \(\mu_0^{\rm eff}=\mu_0/\!\int\Z\,\dd w\),
and the usual ideal induction equation follows. The controlled-limit derivation is
important because it cleanly identifies what must be ``turned back on'' to obtain
new plasma physics: transverse current \(J^w\), mixed field components \(F_{\mu w}\),
and excitation of higher transverse modes.

\subsection{Beyond MHD as conservative topology and energy transport}
A key conceptual distinction of this framework is the difference between:
\begin{enumerate}[leftmargin=1.6em]
\item \textbf{Dissipation to heat} (collisional resistivity, viscosity, Landau damping),
which changes entropy and reduces macroscopic free energy; and
\item \textbf{Export into additional degrees of freedom}, which can appear as
``effective non-ideality'' in brane-facing equations while leaving the bulk system
conservative.
\end{enumerate}
The latter is the characteristic novelty of the \(4\)D interaction model. When
\(A_w\), \(E_w\), \(F_{aw}\), and transverse modes are dynamically active, the brane
subsystem generically becomes an \emph{open} subsystem with explicit exchange terms.
From the brane perspective, those exchange terms can mimic the role usually played
by resistivity in enabling field-line slippage and reconnection, but without requiring
local conversion of magnetic energy directly into heat. Instead, energy may be routed
into mixed-sector fields and/or into higher transverse modes, and later returned.

This interpretation suggests a physically sharp diagnostic criterion for distinguishing
``numerical reconnection'' from ``4D-enabled reconnection'' in simulations: reconnection-like
brane signatures should correlate with measurable transport into transverse channels,
such as nonzero work \(J^wE_w\), growth of mixed-sector energy, and/or growth of mode
energies with \(n\ge 1\) in the Hermite/KK decomposition.

\subsection{Relationship to extended MHD and kinetic plasma theory}
On the brane, the standard extended-MHD hierarchy appears in the usual way:
Hall drift, electron pressure, electron inertia, and collisional resistivity arise
from controlled relaxations of the two-fluid assumptions in
Section~\ref{sec:two_fluid_algebra}. In this sense, the present model is compatible
with (and can reproduce) standard extended-MHD closures.

What is new is that additional correction terms arise even when conventional
extended-MHD closures are held fixed, simply because the parent theory contains
more degrees of freedom than \(3+1\) Maxwell:
\begin{itemize}[leftmargin=1.7em]
\item the scalar-photon channel \(A_w\) and its associated fields \(E_w\) and \(F_{aw}\),
\item explicit brane--bulk exchange via \(J^w\) (leakage sources in projected balances),
\item and finite-localization mode-tower corrections (massive modes with \(m_n^2\propto n\)).
\end{itemize}
In the language of generalized Ohm's law, these effects appear as additional terms
beyond Hall/pressure/inertia, including mixed-field forces and projected unresolved
stresses (Section~\ref{sec:generalized_ohm_4d}). This suggests a unifying view:
\emph{extended MHD} expands the matter closure at fixed \(3+1\) EM, while the
\emph{4D interaction model} expands the EM and topology sector itself.

The framework also provides a natural bridge to kinetic methods. In the kinetic
(Vlasov/PIC) realization, the ``massive simulation'' strategy becomes: evolve the
full \(3\)D phase-space dynamics while representing transverse structure in \(w\) by
a truncated mode basis (or a modest direct \(w\)-grid), thereby retaining the new
topology channels without requiring an intractable full \(4\)D spatial grid at
high resolution.

\subsection{Canonical vorticity constraint and the origin of brane non-ideality}
\label{sec:vorticity_identity}
A useful structural identity in the bulk is the London-type relation
\eqref{eq:london_identity_bulk}:
\begin{equation}
\partial_A v_{s,B}-\partial_B v_{s,A}
= -\frac{q_s}{m_s}F_{AB},
\qquad (A,B\in\{x,y,z,w\}),
\label{eq:discussion_london}
\end{equation}
valid away from phase singularities in the field-based realization (and interpreted
as a canonical-momentum constraint more generally). This identity can be read as a
geometric ``flux locking'' statement: canonical vorticity is tied to gauge flux in
the bulk. A brane observer, however, does not measure the full bulk tensor \(F_{AB}\);
it measures projected quantities \(\Proj{F_{\mu\nu}}\) and experiences explicit exchange
terms when \(J^w\neq 0\). Thus, what appears as a violation of ideal constraints on the
brane can correspond to perfectly consistent bulk evolution plus projection. This is
the mathematical origin of the ``topology export'' interpretation of reconnection in
this framework.

\subsection{The meaning of \texorpdfstring{$w$}{w}, \texorpdfstring{$\Z$}{Z}, and \texorpdfstring{$\W$}{W}}
The model separates two roles often conflated in dimensional reduction:
\begin{itemize}[leftmargin=1.7em]
\item \(\Z(w)\) controls how strongly electromagnetic dynamics is localized near the
operational brane. It sets an intrinsic transverse length scale \(\lambda\) (for
Gaussian localization) and thereby controls when the mode tower becomes important.
\item \(\W(w)\) defines what the brane observer measures. It can be matched to \(\Z/\!\int\Z\)
for a natural ``measure where the field lives'' choice, or chosen differently to
represent finite measurement resolution or a particular observational protocol.
\end{itemize}
In simulation practice, \(\lambda\) and the mode truncation \(N_w\) are therefore not
merely numerical parameters; they are physical modeling knobs that control the strength
of beyond-MHD corrections. A major objective for future work is to calibrate these knobs
against known plasma scales (e.g.\ current-sheet thickness, electron skin depth, or
other microphysical lengths) in specific regimes, thereby making the \(4\)D interaction
channels quantitatively predictive.

\subsection{Falsifiable signatures and near-term simulation targets}
The framework makes several concrete, testable predictions at the level of simulation
diagnostics:
\begin{enumerate}[leftmargin=1.6em]
\item \textbf{Mode-spectrum signature:} reconnection-like brane behavior should correlate with
energy transfer from the \(n=0\) Maxwell mode into higher transverse modes \(n\ge 1\) when
structures approach the localization scale \(r\sim\lambda\). In the Gaussian case, the
spectrum and coupling selection rules are fixed (odd modes decouple from strictly brane-localized
sources), providing a sharp test of correct transverse modeling.

\item \textbf{Scalar-photon signature:} activation of the \(A_w\) sector should be detectable as
nontrivial mixed-field energy (via \(E_w\) and \(F_{aw}\)) and nonzero work channel \(J^wE_w\),
particularly near strong-gradient regions.

\item \textbf{Leakage-topology correlation:} changes in brane helicity (or other topological
diagnostics) should correlate with explicit brane exchange measures such as \(\|\Sleak\|\) and
integrated leakage \(\mathcal{L}_s(t;\Omega)\), rather than correlating solely with collisional
heating.

\item \textbf{Topology-stress localization:} the projected unresolved-stress contribution (the
``topology stress tensor'' defined in Section~\ref{sec:generalized_ohm_4d}) should localize
near reconnection regions and provide a quantitative accounting of the non-ideal term in the
brane induction equation.
\end{enumerate}
These signatures are well suited to the validation/benchmark program in
Section~\ref{sec:validation} and provide a clear path to assessing whether \(4\)D
interaction effects are being captured rather than replaced by numerical diffusion.

\subsection{Limitations and future work}
The present work has several limitations that define a roadmap:
\begin{itemize}[leftmargin=1.7em]
\item \textbf{Matter-sector regime.} The baseline matter model is nonrelativistic. This is
appropriate for MHD-like regimes but will require upgrade (relativistic kinetic matter) for
ultra-relativistic plasmas or extreme-field environments.
\item \textbf{Closure choices.} The two-fluid reduction uses standard closures (scalar pressures,
optional collisional drag). Incorporating anisotropic pressure tensors, heat flux closures, and
collision operators consistent with the bulk conservation laws is a direct extension.
\item \textbf{Transverse calibration.} While the transverse spectral structure is fixed once \(\Z(w)\)
is chosen, mapping \(\lambda\) and \(N_w\) to physical plasma scales in specific regimes remains an
open calibration problem and a central target for future numerical experiments.
\item \textbf{Numerical realizations.} The 3D\(\times\)modes algorithm is designed for massive
simulations, but its robustness across strongly nonlinear regimes must be established by systematic
convergence studies and comparison to direct 4D-grid reference runs.
\item \textbf{Geometry and additional sectors.} The present paper treats \(\Z(w)\) as fixed.
Allowing back-reaction of plasma loads on localization/geometry degrees of freedom (when present in
the broader unified model) is a natural next step, as is coupling to gravitational sectors when
appropriate.
\end{itemize}

Despite these limitations, the model provides a coherent hierarchy: a conservative bulk theory,
an operational brane map, a controlled MHD limit, and a computable set of beyond-MHD correction
channels with clear simulation knobs. This hierarchy is designed to make reconnection and other
topology-changing plasma processes not a special exception requiring ad hoc terms, but an emergent,
diagnosable consequence of evolving the full \(4\)D interaction structure.

% ============================================================
\section{Conclusion}
\label{sec:conclusion}

We have presented a practical ``4D-first'' framework for plasma physics in which the
fundamental interaction dynamics is evolved in a conservative \(4+1\)D bulk and the
usual \(3\)D plasma variables arise as operational brane observables defined by a
projection map. The central objective is not to discard MHD, but to place it inside a
larger, computable hierarchy that (i) reproduces MHD in controlled limits and (ii)
predicts explicit, simulation-resolvable beyond-MHD correction channels when those
limits are relaxed.

\paragraph{What was constructed.}
The paper defined a coupled bulk system consisting of a localized electromagnetic
sector (Section~\ref{sec:em_sector}) and a multi-species plasma matter sector
(Section~\ref{sec:matter_sector}), collected as a single simulation target
(Section~\ref{sec:full_system}). The electromagnetic localization profile \(\Z(w)\)
controls how gauge dynamics is concentrated near the operational brane, while an
independent projection kernel \(\W(w)\) defines what the brane observer measures
(Section~\ref{sec:conventions}). This separation is essential: the bulk system can
remain conservative while the brane subsystem becomes an open system with explicit
exchange terms.

\paragraph{Controlled recovery of standard limits.}
We established a controlled reduction chain
\[
\text{(bulk 4+1 system)}
\;\to\;
\text{(effective 3+1 Maxwell + two-fluid)}
\;\to\;
\text{(MHD)},
\]
under explicit assumptions that suppress transverse structure:
zero-mode dominance, negligible mixed components \(F_{\mu w}\), negligible transverse
current \(J^w\), and quasi-neutral two-fluid ordering (Section~\ref{sec:mhd_reduction}).
In this limit, the localized Maxwell sector reduces to ordinary \(3+1\) Maxwell theory
with an effective coupling \(\mu_0^{\rm eff}=\mu_0/\!\int \Z(w)\dd w\), and the standard
ideal-MHD induction equation follows. The reduction is not merely conceptual: it is
designed to be algebraically checkable and compatible with code verification.

\paragraph{Beyond MHD as explicit 4D interaction channels.}
The principal scientific output is a clean organization of the leading extensions that
appear when the controlled assumptions are relaxed (Section~\ref{sec:beyond_mhd}):
\begin{itemize}[leftmargin=1.7em]
\item a dynamical scalar-photon sector \(A_w\) and mixed field components \(F_{\mu w}\),
introducing additional force channels and transverse electric structure;
\item explicit brane--bulk exchange through \(J^w\neq 0\), producing leakage/source terms in
projected brane balance laws even when bulk conservation is exact;
\item finite-localization transverse-mode corrections (Hermite/KK tower for Gaussian \(\Z\)),
which provide additional storage and transport channels for energy and topology and yield
fixed-coefficient Yukawa corrections in static limits;
\item projected unresolved-stress contributions (``topology stress'' terms) that quantify
how sub-\(w\) structure manifests as effective non-ideality in brane-facing equations.
\end{itemize}
Together, these effects provide a principled alternative to ad hoc resistivity for
topology-changing dynamics: brane reconnection-like behavior can be interpreted as
conservative transport into transverse degrees of freedom rather than local dissipation
into heat.

\paragraph{A concrete computational roadmap.}
We described two complementary simulation architectures (Section~\ref{sec:computation}):
direct \(4\)D-grid evolution in \((x,y,z,w)\) for reference and ground truth, and a
3D\(\times\)modes solver that represents transverse structure in a truncated Hermite basis.
The mode solver is the recommended route for massive runs because the linear transverse
operator becomes diagonal in mode index (fixed masses \(m_n^2\propto n\)), enabling strong
mode-parallel scaling, while nonlinear couplings are handled via pseudo-spectral
reconstruction/projection in \(w\). The diagnostic program emphasizes bulk energy closure,
mode-energy spectra, and explicit leakage measures \(\Sleak^{(s)}\) to distinguish physical
4D-enabled non-ideality from numerical diffusion.

\paragraph{Validation and near-term targets.}
We outlined a benchmark program (Section~\ref{sec:validation}) that (i) reproduces standard
Maxwell/two-fluid/MHD behavior in the controlled limit, and (ii) isolates beyond-MHD signals
through targeted activation of the \(A_w\) sector, transverse currents, and transverse modes.
The near-term implementation targets are therefore clear:
\begin{enumerate}[leftmargin=1.6em]
\item implement the 3D\(\times\)modes Maxwell update with generalized Lorenz gauge control and
verify the Coulomb-plus-Yukawa tower response;
\item couple to a two-fluid or GNLS plasma module with charge-conserving current construction;
\item reproduce controlled-limit MHD benchmarks with \(N_w=1\), then increase \(N_w\) and quantify
how reconnection proxies correlate with energy transfer into mixed fields and higher modes;
\item demonstrate convergence in both \((x,y,z)\) resolution and transverse truncation (\(N_w\),
quadrature order), with bulk energy ledger closure as a primary correctness gate.
\end{enumerate}

\paragraph{Outlook.}
The framework presented here provides a coherent hierarchy: a conservative bulk interaction
model, an operational brane-observer map, a controlled MHD limit, and a structured set of
beyond-MHD correction channels with explicit simulation knobs. Immediate extensions include:
(i) incorporating relativistic kinetic matter for extreme regimes, (ii) adding collision and
radiation models consistent with bulk conservation, and (iii) allowing back-reaction of plasma
loads on localization/geometry degrees of freedom when those sectors are activated. The key
promise is that reconnection and other topology-changing plasma processes become diagnosable
consequences of evolving a larger interaction system, rather than exceptions that require
phenomenological fixes at the outset.

\section*{Acknowledgments}
This work was carried out by an independent researcher. Generative-AI tools were used during
the development of the manuscript to help translate conceptual ideas into mathematical form,
assist with symbolic manipulations, prototype and revise companion Mathematica and SymPy
verification scripts, and improve the exposition. The conceptual direction, hypotheses, and
overall structure of the work were directed by the author. The resulting arguments and
calculations were checked for internal consistency using the accompanying symbolic scripts, but
the work should be regarded as exploratory and would benefit from review by subject-matter
experts. Any errors, omissions, or misunderstandings remain the author's.

% ============================================================
% Appendices
% ============================================================
\appendix

\section{Notation Quick Reference}
\label{app:notation}

This appendix lists the most frequently used symbols and conventions.

\subsection*{Indices and coordinates}

\begin{center}
\begin{tabular}{@{}ll@{}}
\toprule
Symbol & Meaning \\
\midrule
$x^M$ & Bulk spacetime coordinates: $x^M=(t,x^1,x^2,x^3,w)\equiv (t,\mathbf{x},w)$ \\
$\mathbf{x}$ & Brane spatial coordinates: $\mathbf{x}=(x,y,z)\in\mathbb{R}^3$ \\
$\mathbf{X}$ & Bulk spatial coordinates: $\mathbf{X}=(x,y,z,w)\in\mathbb{R}^4$ \\
$M,N,\ldots$ & Bulk spacetime indices in $\{0,1,2,3,4\}$ with $0\equiv t$, $4\equiv w$ \\
$\mu,\nu,\ldots$ & Brane spacetime indices in $\{0,1,2,3\}$ \\
$A,B,\ldots$ & Bulk spatial indices in $\{1,2,3,4\}$ \\
$a,b,c,\ldots$ & Brane spatial indices in $\{1,2,3\}$ \\
$\eta_{MN}$ & Flat bulk metric, $\eta_{MN}=\mathrm{diag}(-1,+1,+1,+1,+1)$ \\
$\epsilon_{abc}$ & Brane Levi--Civita symbol, $\epsilon_{xyz}=+1$ \\
\bottomrule
\end{tabular}
\end{center}

\subsection*{Differential operators and measures}

\begin{center}
\begin{tabular}{@{}ll@{}}
\toprule
Symbol & Meaning \\
\midrule
$\partial_M$ & Partial derivative $\partial/\partial x^M$ \\
$\nabla_3$ & Brane gradient $(\partial_x,\partial_y,\partial_z)$ \\
$\nabla_4$ & Bulk spatial gradient $(\nabla_3,\partial_w)$ \\
$\nabla_3^2$ & Brane Laplacian $\partial_x^2+\partial_y^2+\partial_z^2$ \\
$\nabla_4^2$ & Bulk spatial Laplacian $\nabla_3^2+\partial_w^2$ \\
$\square_4$ & Brane d'Alembertian $-\partial_t^2+\nabla_3^2$ \\
$\square_5$ & Bulk d'Alembertian $-\partial_t^2+\nabla_3^2+\partial_w^2$ \\
$\int \dd^3x$ & $\int_{\mathbb{R}^3}\dd x\,\dd y\,\dd z$ \\
$\int \dd w$ & $\int_{-\infty}^{+\infty}\dd w$ \\
$\int \dd^4X$ & $\int_{\mathbb{R}^3}\dd^3x\int_{\mathbb{R}}\dd w$ \\
\bottomrule
\end{tabular}
\end{center}

\subsection*{Localization, projection, and effective couplings}

\begin{center}
\begin{tabular}{@{}ll@{}}
\toprule
Symbol & Meaning \\
\midrule
$Z(w)$ & EM localization profile (weight in the bulk Maxwell action) \\
$W(w)$ & Brane projection / measurement kernel, normalized by $\int \dd w\,W(w)=1$ \\
$\overline{Q}$ & Brane observable (projection): $\overline{Q}(t,\mathbf{x})=\int \dd w\,W(w)\,Q(t,\mathbf{x},w)$ \\
$\mathcal{Z}_{\mathrm{int}}$ & Localization integral $\mathcal{Z}_{\mathrm{int}}=\int \dd w\,Z(w)$ \\
$\mu_0^{\mathrm{eff}}$ & Effective brane coupling $\mu_0^{\mathrm{eff}}=\mu_0/\mathcal{Z}_{\mathrm{int}}$ (controlled brane Maxwell limit) \\
$\epsilon_0^{\mathrm{eff}}$ & Effective permittivity satisfying $\epsilon_0^{\mathrm{eff}}\mu_0^{\mathrm{eff}}=c^{-2}$ \\
$\lambda$ & Localization width (Gaussian example: $Z(w)=e^{-w^2/\lambda^2}$) \\
\bottomrule
\end{tabular}
\end{center}

\subsection*{Electromagnetic sector (4+1D)}

\begin{center}
\begin{tabular}{@{}ll@{}}
\toprule
Symbol & Meaning \\
\midrule
$A_M$ & Bulk gauge potential $A_M=(A_0,A_a,A_w)$ \\
$F_{MN}$ & Field strength $F_{MN}=\partial_M A_N-\partial_N A_M$ \\
$J^M$ & Total bulk 5-current $J^M=(J^0,J^a,J^w)$ (sources Maxwell) \\
$\xi$ & Gauge-fixing parameter (generalized Lorenz gauge) \\
$E_A$ & Bulk electric components $E_A\equiv F_{A0}$, with $A\in\{x,y,z,w\}$ \\
$E_w$ & Transverse electric component $E_w\equiv F_{w0}=-\partial_t A_w-\partial_w A_0$ \\
$C_a$ & Mixed components $C_a\equiv F_{aw}=\partial_a A_w-\partial_w A_a$ \\
$\mathbf{E}^{(\mathrm{brane})}$ & Brane electric field: $E_a^{(\mathrm{brane})}\equiv \overline{F_{a0}}$ \\
$\mathbf{B}^{(\mathrm{brane})}$ & Brane magnetic field: $B_a^{(\mathrm{brane})}\equiv \tfrac12\epsilon_{abc}\,\overline{F_{bc}}$ \\
$\rho_q$ & Brane charge density (in the controlled $3+1$ Maxwell limit) \\
$\mathbf{J}$ & Brane current density (in the controlled $3+1$ Maxwell limit) \\
\bottomrule
\end{tabular}
\end{center}

\subsection*{Plasma matter sector (bulk and brane observables)}

\begin{center}
\begin{tabularx}{\linewidth}{@{}p{0.25\linewidth}X@{}}
\toprule
Symbol & Meaning \\
\midrule
$s$ & Species index (e.g.\ $s\in\{\mathrm{e},\mathrm{i}\}$ for electrons/ions) \\
$m_s,\,q_s$ & Mass and fixed species charge label of species $s$; for elementary defect branches $q_s$ may be viewed as the reduced brane charge inherited from the localized-charge ontology, while for composite ions it may be treated as the usual effective charge state \\
$\psi_s$ & Species field (GNLS option) \\
$f_s$ & Species distribution function (Vlasov/PIC option) \\
$\rho_s$ & Bulk species number density (GNLS: $\rho_s=|\psi_s|^2$; Vlasov: $\rho_s=\int f_s\,\dd^4v$) \\
$j_s^A$ & Bulk species number current ($A\in\{x,y,z,w\}$) \\
$v_s^A$ & Bulk species velocity, $v_s^A\equiv j_s^A/\rho_s$ (when $\rho_s>0$) \\
$J_s^M$ & Species charge 5-current: $J_s^0=q_s\rho_s$, $J_s^A=q_s j_s^A$ \\
$\overline{\rho_s}$ & Brane-observed species density $\overline{\rho_s}=\int \dd w\,W\,\rho_s$ \\
$\overline{j_s^a}$ & Brane-observed species current $\overline{j_s^a}=\int \dd w\,W\,j_s^a$ \\
$\mathbf{v}_{s,\mathrm{brane}}$ & Brane-observed species velocity $\mathbf{v}_{s,\mathrm{brane}}=\overline{\mathbf{j}_s}/\overline{\rho_s}$ \\
$S_{\mathrm{leak}}^{(s)}$ & Leakage/source in projected continuity: $\partial_t\overline{\rho_s}+\partial_a\overline{j_s^a}=S_{\mathrm{leak}}^{(s)}$ \\
$U_s(\rho_s),\,h_s(\rho_s)$ & EOS potential and enthalpy-like function: $h_s=\dd U_s/\dd\rho_s$ \\
$Q_s$ & Dispersive/quantum potential $Q_s=-(\hbar^2/2m_s)(\nabla_4^2\sqrt{\rho_s})/\sqrt{\rho_s}$ (GNLS option) \\
\bottomrule
\end{tabularx}
\end{center}

\subsection*{Two-fluid and MHD variables (controlled brane limit)}

\begin{center}
\begin{tabularx}{\linewidth}{@{}p{0.25\linewidth}X@{}}
\toprule
Symbol & Meaning \\
\midrule
$n_s$ & Brane species number density (controlled limit, with $s\in\{\mathrm{e},\mathrm{i}\}$) \\
$\mathbf{v}_s$ & Brane species velocity (3-vector) \\
$\rho$ & Single-fluid mass density $\rho=m_{\mathrm{i}}n_{\mathrm{i}}+m_{\mathrm{e}}n_{\mathrm{e}}$ \\
$\mathbf{v}$ & Center-of-mass (single-fluid) velocity on the brane (3-vector) \\
$\mathbf{J}$ & Brane current density, $\mathbf{J}=e n(\mathbf{v}_{\mathrm{i}}-\mathbf{v}_{\mathrm{e}})$ (quasi-neutral) \\
$p_s,\;p$ & Species pressures and total pressure $p=p_{\mathrm{i}}+p_{\mathrm{e}}$ \\
$\eta$ & (Optional) resistivity in resistive closures, $\mathbf{E}+\mathbf{v}\times\mathbf{B}=\eta\mathbf{J}+\cdots$ \\
\bottomrule
\end{tabularx}
\end{center}

\subsection*{Transverse-mode (Hermite/KK) notation}

\begin{center}
\begin{tabularx}{\linewidth}{@{}p{0.25\linewidth}X@{}}
\toprule
Symbol & Meaning \\
\midrule
$f_n(w),\,\phi_n(w)$ & Transverse basis functions (Gaussian case: Hermite family) \\
$a_\mu^{(n)}(t,\mathbf{x})$ & Brane mode amplitude of $A_\mu$ in the expansion $A_\mu=\sum_n a_\mu^{(n)} f_n(w)$ \\
$m_n$ & Mode mass, Gaussian case: $m_n^2=2n/\lambda^2$ \\
$c_n$ & Mode coupling weight to brane-localized sources (defined from $f_n(0)$ and $\|f_n\|_Z$) \\
$N_w$ & Number of retained transverse modes in a truncated expansion \\
$Q_w$ & Quadrature order / number of transverse nodes in pseudo-spectral projection \\
\bottomrule
\end{tabularx}
\end{center}

\section{Derivation Details: Maxwell Reduction and \texorpdfstring{$\mu_0^{\rm eff}$}{mu0eff}}
\label{app:maxwell_reduction}

This appendix provides the detailed derivation of the effective \(3+1\) Maxwell sector
and the appearance of the effective coupling
\(\mu_0^{\rm eff}=\mu_0/\!\int \dd w\,Z(w)\) used throughout the paper.
We give two equivalent derivations:
(i) reduction at the level of the action, and
(ii) reduction at the level of the field equations.
Both make clear the assumptions under which the reduction is valid and identify
the precise terms that spoil it when \(w\)-structure is excited.

\subsection{Localized Maxwell action in \texorpdfstring{\(4+1\)}{4+1}D}
\label{app:maxwell_action}

We take the bulk electromagnetic action (with an optional gauge-fixing term)
\begin{equation}
S_{\rm EM}[A;Z]
=
\int \dd t\,\dd^3x\,\dd w\;
\left[
-\frac{Z(w)}{4\mu_0}\,F_{MN}F^{MN}
-\frac{1}{2\xi\mu_0}(\partial\!\cdot\!A)^2
- A_M J^M
\right],
\label{eq:app_em_action}
\end{equation}
where \(M,N\in\{0,1,2,3,4\}\) with \(0\equiv t\) and \(4\equiv w\),
\(F_{MN}=\partial_M A_N-\partial_N A_M\), and \(Z(w)\ge 0\) is integrable:
\begin{equation}
Z_{\rm int}\equiv \int_{-\infty}^{+\infty}\dd w\;Z(w) < \infty.
\label{eq:app_Zint}
\end{equation}
(For the Gaussian example \(Z(w)=e^{-w^2/\lambda^2}\), one has \(Z_{\rm int}=\lambda\sqrt{\pi}\).)

\paragraph{Dimensional note.}
Because \(\int \dd w\,Z\) carries units of length if \(Z\) is dimensionless, the bulk
constant \(\mu_0\) in \eqref{eq:app_em_action} should be understood as the \(4+1\) coupling
whose units make the action consistent; after reduction, \(\mu_0^{\rm eff}\) has the usual
\(3+1\) units. (Equivalently, one may label the bulk constant \(\mu_{0,5}\) and define
\(\mu_0^{\rm eff}\equiv \mu_{0,5}/Z_{\rm int}\).)

\subsection{Reduction at the level of the action}
\label{app:action_reduction}

We now impose the controlled brane/zero-mode assumptions used in
Section~\ref{sec:mhd_reduction}, specialized here to the EM sector:

\begin{enumerate}[leftmargin=1.6em]
\item \textbf{Zero-mode dominance for brane components:}
\begin{equation}
A_\mu(t,\mathbf{x},w)\;\simeq\;a_\mu(t,\mathbf{x}),
\qquad
\partial_w A_\mu \simeq 0,
\qquad \mu\in\{0,1,2,3\}.
\label{eq:app_zero_mode}
\end{equation}

\item \textbf{Suppressed mixed sector in this limit:}
\begin{equation}
F_{\mu w}\simeq 0.
\label{eq:app_mixed_supp}
\end{equation}
A sufficient condition is axial gauge \(A_w=0\) together with \(\partial_w A_\mu\simeq 0\).
\end{enumerate}

Under \eqref{eq:app_zero_mode}--\eqref{eq:app_mixed_supp},
\begin{equation}
F_{\mu\nu}(t,\mathbf{x},w)\simeq f_{\mu\nu}(t,\mathbf{x})
\equiv \partial_\mu a_\nu-\partial_\nu a_\mu,
\qquad
F_{\mu w}\simeq 0,
\label{eq:app_F_reduction}
\end{equation}
and therefore
\begin{equation}
F_{MN}F^{MN}
\simeq
f_{\mu\nu}f^{\mu\nu}.
\label{eq:app_Fsq_reduction}
\end{equation}

Substituting \eqref{eq:app_Fsq_reduction} into \eqref{eq:app_em_action} yields an effective
\(3+1\) action:
\begin{align}
S_{\rm EM}
&\simeq
\int \dd t\,\dd^3x\;
\left[
-\frac{1}{4\mu_0}\left(\int \dd w\,Z(w)\right) f_{\mu\nu}f^{\mu\nu}
-\int \dd w\;a_\mu J^\mu(t,\mathbf{x},w)
\right]
\;+\;\text{(gauge-fixing)},
\label{eq:app_action_step1}
\\[4pt]
&=
\int \dd t\,\dd^3x\;
\left[
-\frac{1}{4\mu_0^{\rm eff}} f_{\mu\nu}f^{\mu\nu}
- a_\mu j_{\rm eff}^\mu(t,\mathbf{x})
\right]
\;+\;\text{(gauge-fixing)},
\label{eq:app_action_step2}
\end{align}
where we define the effective brane current and effective coupling by
\begin{equation}
j_{\rm eff}^\mu(t,\mathbf{x})
\equiv
\int_{-\infty}^{+\infty}\dd w\;J^\mu(t,\mathbf{x},w),
\qquad
\mu_0^{\rm eff}\equiv \frac{\mu_0}{Z_{\rm int}}.
\label{eq:app_mu0eff_def}
\end{equation}

\paragraph{Brane-localized sources as a special case.}
If the source is strictly brane-localized,\\
\(J^\mu(t,\mathbf{x},w)=j_b^\mu(t,\mathbf{x})\delta(w)\),
then \(j_{\rm eff}^\mu=j_b^\mu\), and the reduced action is exactly the standard Maxwell action
(with \(\mu_0\to\mu_0^{\rm eff}\)) coupled to the brane current.

\paragraph{Gauge-fixing term.}
If \((\partial\!\cdot\!A)\) is also \(w\)-independent in the controlled limit, the gauge-fixing term
reduces with the same prefactor \(Z_{\rm int}\) and is absorbed by the replacement
\(\mu_0\to\mu_0^{\rm eff}\) exactly as above.

\subsection{Reduction at the level of the field equations}
\label{app:eom_reduction}

Varying \eqref{eq:app_em_action} gives the localized Maxwell equations
\begin{equation}
\partial_M\!\left(Z(w)F^{MN}\right) + \frac{1}{\xi}\partial^N(\partial\!\cdot\!A)
=
\mu_0 J^N.
\label{eq:app_localized_maxwell}
\end{equation}

We focus on the brane components \(N=\nu\in\{0,1,2,3\}\) and integrate over \(w\):
\begin{equation}
\int \dd w\;\partial_M\!\left(ZF^{M\nu}\right)
+\frac{1}{\xi}\int \dd w\;\partial^\nu(\partial\!\cdot\!A)
=
\mu_0 \int \dd w\;J^\nu.
\label{eq:app_integrate_step1}
\end{equation}
Split the first term into \(\mu\) and \(w\) contributions and integrate by parts in \(w\):
\begin{equation}
\partial_\mu\!\left(\int \dd w\,Z F^{\mu\nu}\right)
+
\Big[Z F^{w\nu}\Big]_{-\infty}^{+\infty}
+
\frac{1}{\xi}\,\partial^\nu\!\left(\int \dd w\,(\partial\!\cdot\!A)\right)
=
\mu_0\,j_{\rm eff}^\nu.
\label{eq:app_integrated_maxwell}
\end{equation}
Assuming localized fields such that the boundary flux term vanishes,
\(\big[Z F^{w\nu}\big]_{-\infty}^{+\infty}=0\), and invoking the controlled-limit suppression
\(F^{w\nu}\simeq 0\) and \(w\)-independence \(F^{\mu\nu}\simeq f^{\mu\nu}\), we obtain
\begin{equation}
\partial_\mu\!\left(\int \dd w\,Z\right) f^{\mu\nu}
+\frac{1}{\xi}\,\partial^\nu(\partial_\mu a^\mu)\left(\int \dd w\,Z\right)
=
\mu_0\,j_{\rm eff}^\nu.
\label{eq:app_integrated_maxwell_zero_mode}
\end{equation}
Dividing by \(Z_{\rm int}\) yields the standard \(3+1\) Maxwell equation with effective coupling
\begin{equation}
\partial_\mu f^{\mu\nu} + \frac{1}{\xi}\partial^\nu(\partial_\mu a^\mu)
=
\mu_0^{\rm eff}\,j_{\rm eff}^\nu,
\qquad
\mu_0^{\rm eff}=\frac{\mu_0}{Z_{\rm int}}.
\label{eq:app_maxwell_eff_final}
\end{equation}

\paragraph{Where the reduction fails.}
Equation \eqref{eq:app_integrated_maxwell} makes the failure modes explicit:
\begin{itemize}[leftmargin=1.7em]
\item If \(F^{w\nu}\) is not negligible, then the boundary/flux term and the \(w\)-mixing in
\(ZF^{w\nu}\) re-enter and the brane sector is not closed.
\item If \(F^{\mu\nu}\) has significant \(w\)-dependence in the support of \(Z\), then
\(\int \dd w\,Z F^{\mu\nu}\neq Z_{\rm int} f^{\mu\nu}\) and higher transverse modes contribute.
\end{itemize}
These are precisely the channels that generate beyond-MHD corrections in the brane-facing plasma
equations.

\subsection{Alternative viewpoint: KK/Hermite normalization and the zero mode}
\label{app:mode_normalization_mu0eff}

For Gaussian localization (and more generally for any integrable \(Z\)), the transverse
Sturm--Liouville problem defines an orthogonal basis \(\{f_n(w)\}\) with weight \(Z\).
Writing a mode expansion for the brane components,
\begin{equation}
A_\mu(t,\mathbf{x},w)=\sum_{n\ge 0} a_\mu^{(n)}(t,\mathbf{x})\,f_n(w),
\label{eq:app_mode_expansion}
\end{equation}
and inserting into the quadratic action produces a sum of decoupled kinetic terms with
coefficients given by the weighted norms \(\|f_n\|_Z^2=\int \dd w\,Z f_n^2\).

The massless (Maxwell) mode is the constant mode \(f_0(w)=1\), whose norm is
\begin{equation}
\|f_0\|_Z^2 = \int \dd w\,Z(w) = Z_{\rm int}.
\label{eq:app_f0_norm}
\end{equation}
Thus the kinetic term for \(a_\mu^{(0)}\) is
\begin{equation}
-\frac{1}{4\mu_0}\|f_0\|_Z^2\,f_{\mu\nu}f^{\mu\nu}
=
-\frac{1}{4\mu_0^{\rm eff}}\,f_{\mu\nu}f^{\mu\nu},
\qquad
\mu_0^{\rm eff}=\frac{\mu_0}{\|f_0\|_Z^2}=\frac{\mu_0}{Z_{\rm int}},
\label{eq:app_mu0eff_from_modes}
\end{equation}
which reproduces the result from the action/equation reductions.
From this perspective, \(\mu_0^{\rm eff}\) is simply the bulk coupling rescaled by the
normalization of the zero mode.

\subsection{Connection to \texorpdfstring{$\epsilon_0^{\rm eff}$}{eps0eff} and the speed of light}
\label{app:eps0eff}

If one wishes to retain a fixed brane light speed \(c\) in SI form, then the effective
permittivity is defined by
\begin{equation}
\epsilon_0^{\rm eff}\,\mu_0^{\rm eff} = \frac{1}{c^2}.
\label{eq:app_eps0eff_def}
\end{equation}
In practice, once \(\mu_0^{\rm eff}\) is fixed by the localization integral \(Z_{\rm int}\),
this equation defines \(\epsilon_0^{\rm eff}\). In natural units \(c=1\), the relation is
\(\epsilon_0^{\rm eff}=1/\mu_0^{\rm eff}\).

\paragraph{Operational note.}
The reduction in this appendix concerns the \emph{effective} brane Maxwell sector sourced by the
\(w\)-integrated current \(j_{\rm eff}^\mu\). Separately, brane-measured fields in the main text
are defined by the operational projection kernel \(W(w)\) via \(\overline{Q}=\int \dd w\,WQ\).
These are distinct operations (reduction vs.\ measurement) and generally coincide only under the
controlled-limit assumptions where \(w\)-dependence is negligible.

\section{Derivation Details: Two-Fluid \texorpdfstring{$\to$}{->} MHD}
\label{app:twofluid_to_mhd}

This appendix records the standard algebraic reduction from a brane two-fluid plasma
(electrons + ions) to single-fluid MHD. The only modification relative to textbook
SI-MHD is that, in our framework, \(\mu_0\) may be replaced by \(\mu_0^{\rm eff}\)
in Amp\`ere's law if one is working in the controlled brane-Maxwell limit of
Appendix~\ref{app:maxwell_reduction}. Otherwise the reduction is unchanged.
Equivalently, one may package the same localization dependence into fixed reduced
brane source charges \(q_{\rm eff}=q_*/\sqrt{\Zint}\), but we keep
\(\mu_0^{\rm eff}\) as the plasma-facing notation in this appendix.

Throughout, all fields and fluid variables in this appendix are \emph{brane} \(3\)-vectors
and \(3\)-scalar densities; we use \(\gradthree\) as the brane gradient.

\subsection{Starting point: brane two-fluid + Maxwell}
\label{app:twofluid_start}

Consider two charged species: ions (\(\ion\)) and electrons (\(\ele\)):
\begin{equation}
q_\ion = +e,\qquad q_\ele=-e.
\end{equation}
Let \(n_s(t,\mathbf{x})\) be the number density, \(\mathbf{v}_s(t,\mathbf{x})\) the
species velocity, and \(p_s(t,\mathbf{x})\) the species pressure.
Here \(e>0\) denotes the fixed brane-observed elementary charge magnitude
(equivalently \(e_{\rm eff}\) in the localized-defect reading).

\paragraph{Continuity (each species).}
\begin{equation}
\partial_t n_s + \gradthree\cdot(n_s \mathbf{v}_s)=0,
\qquad s\in\{\ion,\ele\}.
\label{eq:app_species_cont}
\end{equation}

\paragraph{Momentum (each species).}
\begin{equation}
m_s n_s\Big(\partial_t + \mathbf{v}_s\cdot\gradthree\Big)\mathbf{v}_s
=
q_s n_s\Big(\mathbf{E} + \mathbf{v}_s\times\mathbf{B}\Big)
-\gradthree p_s
+\mathbf{R}_s,
\qquad s\in\{\ion,\ele\}.
\label{eq:app_species_mom}
\end{equation}
Here \(\mathbf{R}_s\) is the interspecies momentum-exchange (collisional) term, taken
to satisfy total momentum conservation
\begin{equation}
\mathbf{R}_\ion + \mathbf{R}_\ele = \mathbf{0}.
\label{eq:app_drag_sum_zero}
\end{equation}

\paragraph{Maxwell (brane).}
We use the standard brane Maxwell equations:
\begin{align}
\gradthree\cdot\mathbf{B} &= 0,
&
\gradthree\times\mathbf{E} &= -\partial_t\mathbf{B},
\label{eq:app_maxwell_hom}
\\
\gradthree\cdot\mathbf{E} &= \frac{\rho_q}{\epsilon_0^{\rm eff}},
&
\gradthree\times\mathbf{B} - \frac{1}{c^2}\partial_t\mathbf{E}
&= \mu_0^{\rm eff}\,\mathbf{J}.
\label{eq:app_maxwell_inhom}
\end{align}
If one is not in the controlled reduction regime, replace \(\epsilon_0^{\rm eff}\to\epsilon_0\)
and \(\mu_0^{\rm eff}\to\mu_0\).

The brane charge density and current are
\begin{equation}
\rho_q \equiv e(n_\ion - n_\ele),
\qquad
\mathbf{J} \equiv e(n_\ion\mathbf{v}_\ion - n_\ele\mathbf{v}_\ele).
\label{eq:app_rhoq_J_def}
\end{equation}

\subsection{Single-fluid variables and inversion identities}
\label{app:single_fluid_defs}

Define the total mass density \(\rho\), center-of-mass velocity \(\mathbf{v}\), and
total pressure \(p\) by
\begin{equation}
\rho \equiv m_\ion n_\ion + m_\ele n_\ele,
\qquad
\mathbf{v} \equiv \frac{m_\ion n_\ion \mathbf{v}_\ion + m_\ele n_\ele \mathbf{v}_\ele}{\rho},
\qquad
p\equiv p_\ion + p_\ele.
\label{eq:app_rho_v_p_def}
\end{equation}

\paragraph{Quasi-neutral approximation.}
For MHD-like regimes we adopt quasi-neutrality:
\begin{equation}
n_\ion \simeq n_\ele \equiv n,
\qquad \rho_q \simeq 0.
\label{eq:app_quasineutral}
\end{equation}
Then
\begin{equation}
\rho \simeq (m_\ion+m_\ele)n,
\qquad
\mathbf{J} = en(\mathbf{v}_\ion - \mathbf{v}_\ele).
\label{eq:app_quasineutral_rhoJ}
\end{equation}

\paragraph{Exact inversion for \(\mathbf{v}_\ion\) and \(\mathbf{v}_\ele\) (quasi-neutral).}
Solving \eqref{eq:app_rho_v_p_def} and \eqref{eq:app_quasineutral_rhoJ} for the species
velocities gives
\begin{equation}
\mathbf{v}_\ion
=
\mathbf{v} + \frac{m_\ele}{m_\ion+m_\ele}\frac{\mathbf{J}}{en},
\qquad
\mathbf{v}_\ele
=
\mathbf{v} - \frac{m_\ion}{m_\ion+m_\ele}\frac{\mathbf{J}}{en}.
\label{eq:app_vi_ve_exact}
\end{equation}
In the electron-inertia-suppressed limit \(m_\ele/m_\ion\ll 1\),
\begin{equation}
\mathbf{v}_\ion \approx \mathbf{v},
\qquad
\mathbf{v}_\ele \approx \mathbf{v} - \frac{\mathbf{J}}{en}.
\label{eq:app_vi_ve_approx}
\end{equation}

\subsection{Mass continuity}
\label{app:mass_continuity}

Multiply \eqref{eq:app_species_cont} by \(m_s\) and sum over species:
\begin{align}
\partial_t\rho + \gradthree\cdot\!\Big(m_\ion n_\ion \mathbf{v}_\ion + m_\ele n_\ele \mathbf{v}_\ele\Big)
&=0,
\nonumber\\
\Rightarrow\qquad
\partial_t\rho + \gradthree\cdot(\rho\,\mathbf{v})
&=0.
\label{eq:app_mass_continuity}
\end{align}
This is the standard single-fluid continuity equation.

\subsection{Single-fluid momentum equation}
\label{app:single_fluid_momentum}

Add the ion and electron momentum equations \eqref{eq:app_species_mom}. Using
\eqref{eq:app_drag_sum_zero}, the collisional exchange cancels. The electric forces also
cancel because \(q_\ion n_\ion + q_\ele n_\ele = \rho_q \simeq 0\) under quasi-neutrality.
The magnetic forces combine into \(\mathbf{J}\times\mathbf{B}\):
\begin{align}
q_\ion n_\ion(\mathbf{v}_\ion\times\mathbf{B}) + q_\ele n_\ele(\mathbf{v}_\ele\times\mathbf{B})
&=
e n_\ion(\mathbf{v}_\ion\times\mathbf{B}) - e n_\ele(\mathbf{v}_\ele\times\mathbf{B})
\nonumber\\
&= \big(e n_\ion\mathbf{v}_\ion - e n_\ele\mathbf{v}_\ele\big)\times\mathbf{B}
= \mathbf{J}\times\mathbf{B}.
\label{eq:app_JxB_identity}
\end{align}

Thus, the summed momentum equation becomes
\begin{equation}
m_\ion n_\ion\Big(\partial_t + \mathbf{v}_\ion\cdot\gradthree\Big)\mathbf{v}_\ion
+
m_\ele n_\ele\Big(\partial_t + \mathbf{v}_\ele\cdot\gradthree\Big)\mathbf{v}_\ele
=
\mathbf{J}\times\mathbf{B} - \gradthree p.
\label{eq:app_momentum_sum_exact}
\end{equation}

\paragraph{MHD closure on inertial terms.}
Expressing the left-hand side of \eqref{eq:app_momentum_sum_exact} exactly in terms of
\(\rho\) and \(\mathbf{v}\) introduces additional terms involving the relative drift
\(\mathbf{v}_\ion-\mathbf{v}_\ele \propto \mathbf{J}\). In standard MHD ordering these
``drift-stress'' corrections are of higher order and are neglected, yielding the usual
single-fluid momentum equation:
\begin{equation}
\rho\Big(\partial_t + \mathbf{v}\cdot\gradthree\Big)\mathbf{v}
=
\mathbf{J}\times\mathbf{B} - \gradthree p.
\label{eq:app_mhd_momentum}
\end{equation}
(If desired, the neglected terms can be retained and written as the divergence of a
current-related stress tensor; this is one pathway to extended-MHD corrections.)

\subsection{Generalized Ohm's law}
\label{app:ohm_law}

Start from the electron momentum equation \eqref{eq:app_species_mom}:
\begin{equation}
m_\ele n_\ele\Big(\partial_t + \mathbf{v}_\ele\cdot\gradthree\Big)\mathbf{v}_\ele
=
- e n_\ele\Big(\mathbf{E} + \mathbf{v}_\ele\times\mathbf{B}\Big)
-\gradthree p_\ele
+\mathbf{R}_\ele.
\label{eq:app_electron_mom_again}
\end{equation}
Solve for \(\mathbf{E}\):
\begin{equation}
\mathbf{E} + \mathbf{v}_\ele\times\mathbf{B}
=
-\frac{1}{e n_\ele}\gradthree p_\ele
+\frac{1}{e n_\ele}\mathbf{R}_\ele
-\frac{m_\ele}{e}\Big(\partial_t + \mathbf{v}_\ele\cdot\gradthree\Big)\mathbf{v}_\ele.
\label{eq:app_ohm_ve_form}
\end{equation}

\paragraph{Resistive closure (optional).}
A standard electron--ion drag model is
\begin{equation}
\mathbf{R}_\ele = m_\ele n \nu_{ei}(\mathbf{v}_\ion-\mathbf{v}_\ele),
\label{eq:app_Re_drag}
\end{equation}
which, under quasi-neutrality, becomes
\begin{equation}
\frac{1}{e n}\mathbf{R}_\ele
=
\frac{m_\ele \nu_{ei}}{e^2 n}\,\mathbf{J}
\equiv \eta\,\mathbf{J},
\qquad
\eta \equiv \frac{m_\ele \nu_{ei}}{e^2 n}.
\label{eq:app_eta_def}
\end{equation}

\paragraph{Transform from \(\mathbf{v}_\ele\) to \(\mathbf{v}\).}
Using \(\mathbf{v}_\ele \approx \mathbf{v} - \mathbf{J}/(en)\) from \eqref{eq:app_vi_ve_approx},
\begin{align}
\mathbf{E} + \mathbf{v}\times\mathbf{B}
&=
\Big(\mathbf{E} + \mathbf{v}_\ele\times\mathbf{B}\Big)
+ \Big(\mathbf{v}-\mathbf{v}_\ele\Big)\times\mathbf{B}
\nonumber\\
&=
\Big(\mathbf{E} + \mathbf{v}_\ele\times\mathbf{B}\Big)
+\frac{\mathbf{J}}{en}\times\mathbf{B}.
\label{eq:app_ohm_transform}
\end{align}
Substituting \eqref{eq:app_ohm_ve_form} and \eqref{eq:app_eta_def} yields the generalized
Ohm law in the standard \((\mathbf{E}+\mathbf{v}\times\mathbf{B})\) form:
\begin{equation}
\boxed{
\mathbf{E} + \mathbf{v}\times\mathbf{B}
=
\underbrace{\eta\,\mathbf{J}}_{\text{resistive}}
+
\underbrace{\frac{\mathbf{J}\times\mathbf{B}}{en}}_{\text{Hall}}
-
\underbrace{\frac{1}{en}\gradthree p_\ele}_{\text{electron pressure}}
-
\underbrace{\frac{m_\ele}{e}\Big(\partial_t + \mathbf{v}_\ele\cdot\gradthree\Big)\mathbf{v}_\ele}_{\text{electron inertia}}
}.
\label{eq:app_ohm_general_full}
\end{equation}
Dropping the electron-inertia term yields the Hall-resistive form used most commonly
in extended MHD.

\paragraph{Ideal MHD (Ohm).}
The ideal-MHD limit is the controlled sub-limit in which the right-hand side of
\eqref{eq:app_ohm_general_full} is negligible:
\begin{equation}
\mathbf{E} = -\mathbf{v}\times\mathbf{B}.
\label{eq:app_ohm_ideal}
\end{equation}

\subsection{Induction equation}
\label{app:induction}

Faraday's law \eqref{eq:app_maxwell_hom} implies
\begin{equation}
\partial_t\mathbf{B} = -\gradthree\times\mathbf{E}.
\label{eq:app_faraday}
\end{equation}
Insert the generalized Ohm law \eqref{eq:app_ohm_general_full} to obtain
\begin{align}
\partial_t\mathbf{B}
&=
\gradthree\times(\mathbf{v}\times\mathbf{B})
-\gradthree\times(\eta\,\mathbf{J})
-\gradthree\times\!\left(\frac{\mathbf{J}\times\mathbf{B}}{en}\right)
+\gradthree\times\!\left(\frac{1}{en}\gradthree p_\ele\right)
\nonumber\\
&\hspace{2.5em}
+\;\gradthree\times\!\left[
\frac{m_\ele}{e}\Big(\partial_t + \mathbf{v}_\ele\cdot\gradthree\Big)\mathbf{v}_\ele
\right].
\label{eq:app_induction_general}
\end{align}

\paragraph{Ideal induction equation.}
Using \eqref{eq:app_ohm_ideal}, \eqref{eq:app_faraday} becomes
\begin{equation}
\partial_t\mathbf{B} = \gradthree\times(\mathbf{v}\times\mathbf{B}),
\qquad
\gradthree\cdot\mathbf{B}=0,
\label{eq:app_induction_ideal}
\end{equation}
which is the standard ideal-MHD flux-freezing statement.

\subsection{Closing the MHD system}
\label{app:mhd_closure}

A minimal ideal-MHD closure consists of:
\begin{align}
\partial_t\rho + \gradthree\cdot(\rho\,\mathbf{v}) &= 0,
\label{eq:app_mhd_cont}
\\
\rho\Big(\partial_t + \mathbf{v}\cdot\gradthree\Big)\mathbf{v}
&= \mathbf{J}\times\mathbf{B} - \gradthree p,
\label{eq:app_mhd_mom}
\\
\partial_t\mathbf{B} &= \gradthree\times(\mathbf{v}\times\mathbf{B}),
\qquad
\gradthree\cdot\mathbf{B}=0,
\label{eq:app_mhd_induction}
\end{align}
supplemented by an equation of state (and, if needed, an energy equation). A common
additional closure in the MHD ordering is to neglect displacement current in Amp\`ere's
law \eqref{eq:app_maxwell_inhom}:
\begin{equation}
\mathbf{J} \simeq \frac{1}{\mu_0^{\rm eff}}\,\gradthree\times\mathbf{B},
\label{eq:app_ampere_no_disp}
\end{equation}
thereby eliminating \(\mathbf{J}\) in favor of \(\mathbf{B}\). In resistive MHD, one
keeps \(-\gradthree\times(\eta\mathbf{J})\) in \eqref{eq:app_induction_general} and uses
\eqref{eq:app_ampere_no_disp} to write it as a magnetic diffusion term.

\paragraph{Where the present framework modifies this appendix.}
The algebra above is the standard two-fluid \(\to\) MHD reduction on a \(3\)D domain.
In the full \(4+1\) theory, additional terms enter \(\mathbf{E}\) and the brane balances
when \(J^w\), \(F_{\mu w}\), and transverse-mode excitations are not negligible. Those
effects appear as explicit beyond-MHD corrections in the main text and do not alter the
internal consistency of the classical reduction recorded here.

\section{Mode Decomposition in \texorpdfstring{$w$}{w} (Hermite / Soft-Wall)}
\label{sec:kk_modes}

This appendix collects the transverse (\(w\)) mode machinery induced by the localized
Maxwell kinetic term. The practical purpose is twofold:
(i) to make the \emph{finite-localization corrections} computable and interpretable
(Coulomb \(+\) Yukawa tower), and
(ii) to provide the mathematical foundation for the \(\text{3D}\times\text{modes}\)
HPC algorithm described in Section~\ref{sec:computation}.

\subsection{Sturm--Liouville problem from localized Maxwell}
\label{app:sl_from_localized_maxwell}

Consider the localized Maxwell sector in the bulk,
\begin{equation}
\partial_M\!\left(\Z(w)F^{MN}\right) + \frac{1}{\xi}\partial^N(\partial\!\cdot\!A)
= \mu_0\,J^N,
\qquad
F_{MN}=\partial_M A_N-\partial_N A_M,
\qquad
M,N\in\{0,1,2,3,4\},
\end{equation}
with \(\Z(w)\ge 0\) integrable and time-independent. In a source-free region and with
a hyperbolic gauge choice (e.g.\ generalized Lorenz gauge), the brane components
\(A_\mu\) (\(\mu\in\{0,1,2,3\}\)) separate into a brane wave operator and a transverse
operator acting on \(w\):
\begin{equation}
\square_4 A_\mu + \frac{1}{\Z(w)}\frac{\dd}{\dd w}\!\left(\Z(w)\frac{\dd A_\mu}{\dd w}\right)
= 0,
\qquad
\square_4\equiv -\partial_t^2+\nabla_3^2.
\label{eq:w_operator_separation}
\end{equation}

This motivates a mode expansion of the form
\begin{equation}
A_\mu(x,w)=\sum_{n\ge 0} a_\mu^{(n)}(x)\,f_n(w),
\qquad x\equiv (t,\mathbf{x}),
\label{eq:kk_expansion_app}
\end{equation}
where the transverse profiles \(f_n(w)\) are eigenfunctions of the weighted
Sturm--Liouville operator:
\begin{equation}
-\frac{\dd}{\dd w}\!\left(\Z(w)\frac{\dd f_n}{\dd w}\right)=m_n^2\,\Z(w)\,f_n(w).
\label{eq:sl_problem_app}
\end{equation}

The natural inner product and norm are the ones induced by the action:
\begin{equation}
\langle f,g\rangle_{\Z}\equiv \int_{-\infty}^{+\infty}\dd w\;\Z(w)\,f(w)g(w),
\qquad
\|f_n\|_{\Z}^2 \equiv \langle f_n,f_n\rangle_{\Z}.
\label{eq:Z_inner_product_app}
\end{equation}
With appropriate boundary/localization conditions (e.g.\ sufficiently fast decay so that
boundary terms vanish), the operator in \eqref{eq:sl_problem_app} is self-adjoint with
respect to \(\langle\cdot,\cdot\rangle_{\Z}\), implying orthogonality of distinct modes:
\(\langle f_n,f_m\rangle_{\Z}\propto \delta_{nm}\).

\subsection{Gaussian soft-wall and the Hermite spectrum}
\label{app:hermite_soft_wall}

A particularly convenient ``soft-wall'' choice is a Gaussian localization profile,
\begin{equation}
\Z(w)=\exp\!\left(-\frac{w^2}{\lambda^2}\right),
\label{eq:Z_gaussian_app}
\end{equation}
where \(\lambda\) is the localization width.

Writing \(y\equiv w/\lambda\), equation \eqref{eq:sl_problem_app} becomes
\begin{equation}
\frac{\dd^2 f_n}{\dd y^2} - 2y\,\frac{\dd f_n}{\dd y} + (m_n^2\lambda^2)\,f_n = 0,
\label{eq:hermite_ode_app}
\end{equation}
which is the Hermite differential equation. A complete set of polynomial solutions is
given by the physicists' Hermite polynomials:
\begin{equation}
f_n(w)=H_n\!\left(\frac{w}{\lambda}\right),
\qquad
m_n^2=\frac{2n}{\lambda^2},
\qquad n=0,1,2,\ldots
\label{eq:hermite_spectrum_app}
\end{equation}
The \(n=0\) mode is massless and corresponds to the standard \(3+1\) Maxwell field on
the brane. The \(n\ge 1\) modes behave as a tower of massive \(3+1\) Proca-type modes,
with masses set by the localization scale \(\lambda\).

\paragraph{Normalized basis (optional).}
It is often convenient to work with \(\Z\)-orthonormal functions
\(\{\phi_n\}\) rather than raw polynomials:
\begin{equation}
\phi_n(w)\equiv \frac{1}{\sqrt{\|f_n\|_{\Z}^2}}\,f_n(w)
=
\frac{H_n(w/\lambda)}{\sqrt{\lambda\sqrt{\pi}\,2^n n!}},
\qquad
\langle \phi_n,\phi_m\rangle_{\Z}=\delta_{nm},
\label{eq:hermite_normalized_app}
\end{equation}
using the norm computed in \eqref{eq:hermite_norm_app} below.

\subsection{Mode normalization and brane coupling weights}
\label{app:mode_norm_and_coupling}

For the Gaussian soft-wall \eqref{eq:Z_gaussian_app}, the Hermite orthogonality integral gives
\begin{equation}
\|f_n\|_{\Z}^2
=
\int_{-\infty}^{+\infty}\dd w\;
\exp\!\left(-\frac{w^2}{\lambda^2}\right)
H_n\!\left(\frac{w}{\lambda}\right)^2
=
\lambda\sqrt{\pi}\,2^n n!.
\label{eq:hermite_norm_app}
\end{equation}

If the bulk current is brane-localized, \(J^\mu(x,w)=J_b^\mu(x)\delta(w)\), then the
\(n\)-th mode is sourced in proportion to its value at \(w=0\). A convenient coupling
weight is
\begin{equation}
c_n \equiv \frac{f_n(0)^2}{\|f_n\|_{\Z}^2}.
\label{eq:cn_def_app}
\end{equation}
This weight appears directly in the effective brane-to-brane propagator and in static
potentials.

\paragraph{Parity selection rule.}
For Hermite polynomials, \(H_n(0)=0\) for odd \(n\). Therefore
\begin{equation}
c_{2m+1}=0,
\qquad m=0,1,2,\ldots,
\label{eq:odd_modes_decouple_app}
\end{equation}
so odd transverse modes do \emph{not} couple to strictly brane-localized sources at
the symmetry point \(w=0\). Physically this follows from:
(i) \(\Z(w)\) is even,
(ii) the brane sits at the symmetry point, and
(iii) odd transverse profiles integrate to zero against \(\delta(w)\).

For even \(n=2m\), the closed-form value \(H_{2m}(0)=(-1)^m(2m)!/m!\) yields
\begin{equation}
c_{2m}
= \frac{H_{2m}(0)^2}{\lambda\sqrt{\pi}\,2^{2m}(2m)!}
= \frac{1}{\lambda\sqrt{\pi}}\,
\frac{1}{4^m}\binom{2m}{m}.
\label{eq:cn_even_app}
\end{equation}
In particular,
\begin{equation}
c_0=\frac{1}{\lambda\sqrt{\pi}}
=\frac{1}{\Zint},
\qquad
\Zint\equiv \int_{-\infty}^{+\infty}\dd w\;\Z(w)=\lambda\sqrt{\pi},
\label{eq:c0_and_Zint_app}
\end{equation}
which is the same normalization that produces \(\mu_0^{\rm eff}=\mu_0/\Zint\) in the
controlled Maxwell reduction (Appendix~\ref{app:maxwell_reduction}).

\subsection{Effective \texorpdfstring{$3+1$}{3+1} field equations in mode form}
\label{app:mode_field_equations}

Insert the expansion \eqref{eq:kk_expansion_app} into the brane component equation
with sources and project onto mode \(n\) by multiplying by \(\Z(w)f_n(w)\) and integrating
over \(w\). Using orthogonality and assuming boundary terms vanish, one obtains
\begin{equation}
\left(\square_4 + m_n^2\right)a_\mu^{(n)}(x)
=
\mu_0\,J_{\mu}^{(n)}(x)
\;+\;\text{(gauge-driver terms)},
\label{eq:mode_wave_equation_app}
\end{equation}
where the projected current is
\begin{equation}
J_{\mu}^{(n)}(x)\equiv
\frac{1}{\|f_n\|_{\Z}^2}
\int_{-\infty}^{+\infty}\dd w\;\Z(w)\,f_n(w)\,J_\mu(x,w).
\label{eq:mode_projected_current_app}
\end{equation}

For a brane-localized current \(J_\mu(x,w)=J_{b,\mu}(x)\delta(w)\),
\begin{equation}
J_{\mu}^{(n)}(x)=\frac{f_n(0)}{\|f_n\|_{\Z}^2}J_{b,\mu}(x),
\label{eq:mode_source_brane_local_app}
\end{equation}
so the brane-to-brane response involves \(f_n(0)^2/\|f_n\|_{\Z}^2=c_n\).

\paragraph{Interpretation.}
Equation \eqref{eq:mode_wave_equation_app} shows that the localized \(4+1\) Maxwell system
induces a \emph{tower} of \(3+1\) fields:
the \(n=0\) field is ordinary Maxwell;
each \(n\ge 1\) mode behaves like a massive vector field with mass \(m_n\).
In simulations, one truncates to \(n=0,\ldots,N_w-1\) and interprets \(N_w\) as the
amount of transverse structure retained.

\subsection{Coulomb law and Yukawa corrections}
\label{app:yukawa_corrections}

The mode decomposition implies an effective brane-to-brane propagator
\begin{equation}
D_{\rm eff}(k^2)
=
\sum_{n\ge 0}\frac{c_n}{k^2-m_n^2+\mathrm{i}\epsilon},
\qquad
m_n^2=\frac{2n}{\lambda^2},
\label{eq:Deff_app}
\end{equation}
where \(k^2\) is the brane Lorentz scalar built from \((\omega,\mathbf{k})\).

In the static limit (\(\omega=0\)), each massive mode contributes a Yukawa potential.
For a fixed point source \(q_{\rm src}\) on the brane, with
\(J_b^0=q_{\rm src}\,\delta^{(3)}(\mathbf{x})\),
the scalar potential on the brane takes the form
\begin{equation}
A_0(r)
=
\frac{\mu_0 q_{\rm src}}{4\pi r}\sum_{n\ge 0}c_n\,e^{-m_n r},
\qquad r\equiv|\mathbf{x}|.
\label{eq:A0_KK_sum_app}
\end{equation}
Separating the zero mode and using \(\mu_0^{\rm eff}=\mu_0 c_0\), one can write
\begin{equation}
A_0(r)
=
\frac{\mu_0^{\rm eff} q_{\rm src}}{4\pi r}
\left[
1+\sum_{m\ge 1}\frac{c_{2m}}{c_0}\,e^{-m_{2m}r}
\right],
\qquad
m_{2m}=\frac{2\sqrt{m}}{\lambda},
\label{eq:A0_mu_eff_form_app}
\end{equation}
with coefficient ratios
\begin{equation}
\frac{c_{2m}}{c_0}
=
\frac{1}{4^m}\binom{2m}{m}
= 1,\;\frac12,\;\frac{3}{8},\;\frac{5}{16},\;\ldots
\qquad (m=0,1,2,3,\ldots).
\label{eq:cn_ratio_app}
\end{equation}
Therefore the leading deviation from Coulomb at large \(r\) is
\begin{equation}
A_0(r)
=
\frac{\mu_0^{\rm eff} q_{\rm src}}{4\pi r}
\left[
1+\frac12\,e^{-2r/\lambda}
+\mathcal{O}\!\left(e^{-2\sqrt{2}\,r/\lambda}\right)
\right].
\label{eq:leading_yukawa_em_app}
\end{equation}
Here \(q_{\rm src}\) is a fixed source label; for an elementary defect branch in the
canonically normalized brane packaging one may identify
\(q_{\rm src}=q_{\rm eff}=\eta_Q e_{\rm eff}\).

\paragraph{Brane interpretation.}
For \(r\gg \lambda\), brane observers see ordinary Coulomb behavior with coupling
\(\mu_0^{\rm eff}\). For \(r\sim \lambda\), Yukawa corrections become significant and
encode finite-localization physics. In plasma applications, this identifies \(\lambda\)
as a new physical scale controlling when transverse degrees of freedom can influence
brane-facing electromagnetic interactions.

\subsection{How the mode tower is used in large simulations}
\label{app:mode_tower_in_computation}

The \(\text{3D}\times\text{modes}\) algorithm (Section~\ref{sec:computation}) follows directly
from \eqref{eq:mode_wave_equation_app}:

\begin{enumerate}[leftmargin=1.6em]
\item \textbf{Evolve mode amplitudes.}
For each retained \(n\in\{0,\ldots,N_w-1\}\), evolve \(a_M^{(n)}(t,\mathbf{x})\) using a
\(3\)D hyperbolic update (wave stencil) with a diagonal ``mass'' term \(m_n^2 a_M^{(n)}\).
The \(n\)-dependence is local, enabling mode-parallel scaling.

\item \textbf{Reconstruct fields at transverse quadrature nodes (pseudo-spectral).}
Choose quadrature nodes \(\{w_q\}\) adapted to the Gaussian weight and reconstruct
\(A_M(t,\mathbf{x},w_q)=\sum_{n=0}^{N_w-1}a_M^{(n)}(t,\mathbf{x})\,\phi_n(w_q)\).

\item \textbf{Compute nonlinear plasma sources in physical \(w\)-space.}
Evaluate currents \(J^M(t,\mathbf{x},w_q)\) from the matter model (two-fluid, GNLS, or PIC).

\item \textbf{Project sources back to modes.}
Compute \(J^{M,(n)}(t,\mathbf{x})\) by quadrature approximations to
\eqref{eq:mode_projected_current_app}. This yields the mode-forcing terms for the next
Maxwell update.

\end{enumerate}

\paragraph{What is gained by the soft-wall/Hermite choice.}
Unlike a hard-wall truncation (Fourier box), the Gaussian soft-wall:
(i) handles \(w\in\mathbb{R}\) without artificial reflecting boundaries,
(ii) yields a simple spectrum \(m_n^2\propto n\),
(iii) provides an exact parity selection rule \eqref{eq:odd_modes_decouple_app} for brane sources,
and (iv) makes linear evolution mode-diagonal, which is critical for massive simulations.

\paragraph{Extension to the scalar-photon sector \(A_w\).}
Nothing in the Sturm--Liouville construction restricts it to \(A_\mu\) alone:
the same basis may be used to represent \(A_w\) and mixed components \(F_{\mu w}\).
In plasma runs, \(A_w\) is retained as a dynamical degree of freedom unless a controlled
reduction explicitly suppresses it (Section~\ref{sec:mhd_reduction}); its excitation is
diagnosed by mixed-sector energy and by the work channel \(J^wE_w\).

\section{Projected Energy and Helicity Budgets with Leakage}
\label{app:proj_energy_helicity}

This appendix derives the brane-facing (projected) energy and magnetic-helicity budgets,
making explicit how the projection \(\overline{(\cdot)}\) turns the brane subsystem into an
\emph{open} system even when the bulk evolution is conservative. The key message is:

\begin{quote}
\emph{Bulk conservation laws close on the bulk fields. Brane-facing conservation laws close
only after adding explicit leakage terms (transport in \(w\)) and ``projection-stress''
covariances (unresolved transverse structure).}
\end{quote}

\subsection{Projection identities and Reynolds decomposition in \texorpdfstring{\(w\)}{w}}
\label{app:proj_identities}

Recall the operational brane observable map
\begin{equation}
\overline{Q}(t,\mathbf{x})
\;\equiv\;
\int_{-\infty}^{+\infty}\dd w\;W(w)\,Q(t,\mathbf{x},w),
\qquad
\int\dd w\,W(w)=1,
\label{eq:proj_def_app}
\end{equation}
with fixed (time-independent) kernel \(W(w)\). For any sufficiently smooth \(Q\),
\begin{equation}
\overline{\partial_t Q}=\partial_t \overline{Q},
\qquad
\overline{\partial_a Q}=\partial_a \overline{Q},
\label{eq:proj_commute_tx_app}
\end{equation}
while for \(w\)-derivatives one has the integration-by-parts identity
\begin{equation}
\overline{\partial_w Q}
=
\Big[W(w)\,Q\Big]_{-\infty}^{+\infty}
-
\int_{-\infty}^{+\infty}\dd w\;W'(w)\,Q
\;\equiv\;
\Big[WQ\Big]_{-\infty}^{+\infty}
-
\overline{Q}^{\,(\partial W)}.
\label{eq:proj_parts_app}
\end{equation}
If \(W(w)\) decays sufficiently fast (or is compactly supported with \(W=0\) at its
endpoints), the boundary term vanishes and \(\overline{\partial_w Q}=-\overline{Q}^{(\partial W)}\).

\paragraph{Resolved + subscale decomposition.}
For any quantity \(Q(t,\mathbf{x},w)\) define the fluctuation
\begin{equation}
Q'(t,\mathbf{x},w)\equiv Q-\overline{Q},
\qquad \overline{Q'}=0.
\label{eq:reynolds_decomp_app}
\end{equation}
Then quadratic products satisfy
\begin{equation}
\overline{QR}
=
\overline{Q}\,\overline{R}
+
\overline{Q'R'}.
\label{eq:proj_product_app}
\end{equation}
The covariance \(\overline{Q'R'}\) is the precise mathematical object that encodes
``unresolved transverse physics'' (higher modes in the Hermite/KK expansion, or fine
structure in a direct \(w\)-grid).

\subsection{Bulk electromagnetic energy balance (localized Maxwell)}
\label{app:bulk_poynting}

For the localized Maxwell system (no explicit time dependence in \(Z(w)\)), define
bulk electric and magnetic-type components
\begin{equation}
E_A\equiv F_{A0},
\qquad
F_{AB}=\partial_A A_B-\partial_B A_A,
\qquad A,B\in\{x,y,z,w\}.
\end{equation}
A convenient gauge-invariant electromagnetic energy density is
\begin{equation}
u_{\rm EM}
\equiv
\frac{Z(w)}{2\mu_0}
\left(
E_AE_A+\frac12 F_{AB}F_{AB}
\right)
=
\frac{Z(w)}{2\mu_0}
\left(
|\mathbf{E}|^2 + E_w^2 + |\mathbf{B}|^2 + |\mathbf{C}|^2
\right),
\label{eq:uem_decomp_app}
\end{equation}
where \(\mathbf{E}\) and \(\mathbf{B}\) denote the brane \(3\)-vectors
\(E_a=F_{a0}\), \(B_a=\tfrac12\epsilon_{abc}F_{bc}\), and the mixed field
\(\mathbf{C}\) is \(C_a\equiv F_{aw}\).

A natural \(4\)-spatial-dimensional Poynting flux is
\begin{equation}
S_{\rm EM}^A \equiv \frac{Z(w)}{\mu_0}\,E_B F_{BA}.
\label{eq:poynting_def_app}
\end{equation}
In components this splits into the brane flux and the transverse (into-\(w\)) flux:
\begin{align}
S_{\rm EM}^a
&=
\frac{Z}{\mu_0}\Big(E_b F_{ba}-E_w C_a\Big)
=
\frac{Z}{\mu_0}\Big[(\mathbf{E}\times\mathbf{B})_a - E_w C_a\Big],
\label{eq:poynting_brane_app}
\\
S_{\rm EM}^w
&=
\frac{Z}{\mu_0}E_a F_{aw}
=
\frac{Z}{\mu_0}\,\mathbf{E}\cdot\mathbf{C}.
\label{eq:poynting_w_app}
\end{align}
Thus, nonzero mixed structure \(\mathbf{C}\) generically enables energy transport along
the extra dimension.

For conservative sources (and ignoring gauge-fixing terms which do not affect gauge-invariant
budgets), the localized Maxwell equations imply the bulk energy identity
\begin{equation}
\partial_t u_{\rm EM}
+
\partial_A S_{\rm EM}^A
=
-\,J^A E_A
=
-\,(J^aE_a + J^wE_w).
\label{eq:bulk_poynting_app}
\end{equation}
The channel \(J^wE_w\) is the direct ``scalar-photon work'' term: when \(A_w\) is dynamical,
the bulk plasma can exchange energy with the transverse electric field \(E_w\).

\subsection{Projected electromagnetic energy budget and explicit \texorpdfstring{\(w\)}{w}-leakage}
\label{app:proj_energy_budget}

Multiply \eqref{eq:bulk_poynting_app} by \(W(w)\) and integrate over \(w\). Using
\eqref{eq:proj_commute_tx_app} and separating brane and \(w\) divergences yields
\begin{equation}
\partial_t \overline{u_{\rm EM}}
+
\partial_a \overline{S_{\rm EM}^a}
=
-\overline{J^A E_A}
\;+\;
\mathcal{L}_{\rm EM}^{(w)},
\label{eq:proj_energy_local_app}
\end{equation}
where the \emph{transverse leakage power density} is
\begin{align}
\mathcal{L}_{\rm EM}^{(w)}
&\equiv
-\int\dd w\;W(w)\,\partial_w S_{\rm EM}^w
\nonumber\\
&=
-\Big[W S_{\rm EM}^w\Big]_{-\infty}^{+\infty}
+
\int\dd w\;W'(w)\,S_{\rm EM}^w.
\label{eq:energy_leak_term_app}
\end{align}
When the boundary term vanishes, \(\mathcal{L}_{\rm EM}^{(w)}=\int \dd w\,W' S_{\rm EM}^w\).
This is the precise analog of the species-leakage term in the projected continuity equation:
even if the bulk system is conservative, the brane-measured EM energy can increase or decrease
because Poynting flux crosses the measurement profile in \(w\).

\paragraph{Integral form over a brane region.}
For a brane region \(\Omega\subset\mathbb{R}^3\),
\begin{equation}
\frac{\dd}{\dd t}\int_{\Omega}\dd^3x\;\overline{u_{\rm EM}}
=
-\int_{\partial\Omega}\dd S\;n_a\,\overline{S_{\rm EM}^a}
-\int_{\Omega}\dd^3x\;\overline{J^A E_A}
+\int_{\Omega}\dd^3x\;\mathcal{L}_{\rm EM}^{(w)}.
\label{eq:energy_integral_brane_app}
\end{equation}
The last term is the explicit \(w\)-exchange contribution to the brane EM energy ledger.

\subsection{Resolved versus projected EM energy}
\label{app:resolved_vs_projected_energy}

The quantity \(\overline{u_{\rm EM}}=\overline{u_{\rm EM}(E_A,F_{AB})}\) is the
projection of the \emph{bulk} quadratic energy density. A brane observer, however,
often constructs an energy density from \emph{projected fields}, e.g.
\begin{equation}
u_{\rm EM}^{\rm (res)}(\mathbf{x})
\equiv
\frac{1}{2\mu_0^{\rm eff}}
\Big(
|\overline{\mathbf{E}}|^2 + |\overline{\mathbf{B}}|^2
\Big),
\label{eq:resolved_energy_def_app}
\end{equation}
in the controlled limit where \(\mu_0^{\rm eff}\) is meaningful. These two notions differ
once \(w\)-structure is present.

Using the Reynolds decomposition \eqref{eq:reynolds_decomp_app}, the projected quadratic energy
splits into a resolved part plus a covariance (subscale) part:
\begin{equation}
\overline{|\mathbf{E}|^2}
=
|\overline{\mathbf{E}}|^2 + \overline{|\mathbf{E}'|^2},
\qquad
\overline{|\mathbf{B}|^2}
=
|\overline{\mathbf{B}}|^2 + \overline{|\mathbf{B}'|^2},
\qquad\text{etc.}
\label{eq:energy_reynolds_split_app}
\end{equation}
Thus, even in a regime where one focuses only on brane fields, the correct energy ledger
includes an explicitly measurable ``transverse-mode energy'' reservoir given by the covariance
terms. In the Hermite/KK formulation these covariance terms correspond to the energy stored in
the \(n\ge 1\) modes and in the mixed sector \((E_w,\mathbf{C})\).

\paragraph{Interpretation for reconnection diagnostics.}
In resistive MHD, apparent magnetic-energy loss on the brane is often attributed directly to
\(\eta J^2\) heating. Here, brane magnetic-energy loss can instead correspond to:
(i) transport into \(w\) (the \(\mathcal{L}_{\rm EM}^{(w)}\) term), and/or
(ii) transfer into unresolved transverse structure (covariance energy),
without requiring collisional dissipation. These mechanisms are conservative in the bulk.

\subsection{Magnetic helicity identity at fixed \texorpdfstring{\(w\)}{w}}
\label{app:helicity_identity}

Fix \(w\) and define the brane vector potential \(\mathbf{A}(t,\mathbf{x},w)\) and scalar
potential \(\Phi(t,\mathbf{x},w)\equiv A_0(t,\mathbf{x},w)\). Define the brane magnetic field
\(\mathbf{B}\equiv \nabla_3\times\mathbf{A}\) and brane electric field
\(\mathbf{E}\equiv -\nabla_3\Phi-\partial_t\mathbf{A}\). Then the standard local helicity identity
holds pointwise in \(w\):
\begin{equation}
\partial_t\!\big(\mathbf{A}\cdot\mathbf{B}\big)
+
\nabla_3\cdot\Big(\Phi\,\mathbf{B}+\mathbf{E}\times\mathbf{A}\Big)
=
-2\,\mathbf{E}\cdot\mathbf{B}.
\label{eq:helicity_local_app}
\end{equation}
This identity is purely \(3+1\) (with \(w\) acting as a parameter), and is independent of the
details of the \(4+1\) dynamics.

\paragraph{Gauge note.}
Magnetic helicity is gauge invariant on a periodic domain, or for fields satisfying standard
boundary conditions (e.g.\ \(\mathbf{B}\cdot \mathbf{n}=0\) and fixed tangential \(\mathbf{A}\)
on \(\partial\Omega\)), or when formulated as \emph{relative helicity} with respect to a
reference field. In what follows, the boundary flux term is kept explicit.

\subsection{Projected helicity budget and ``helicity leakage'' into transverse structure}
\label{app:proj_helicity_budget}

Project \eqref{eq:helicity_local_app} with \(W(w)\) and integrate over \(w\). Define the
\emph{projected helicity density} and projected helicity flux:
\begin{equation}
\overline{h}(\mathbf{x})
\equiv
\overline{\mathbf{A}\cdot\mathbf{B}}
=
\int \dd w\,W\,(\mathbf{A}\cdot\mathbf{B}),
\qquad
\overline{\mathbf{F}}_h(\mathbf{x})
\equiv
\overline{\Phi\,\mathbf{B}+\mathbf{E}\times\mathbf{A}}.
\label{eq:proj_helicity_defs_app}
\end{equation}
Then
\begin{equation}
\partial_t \overline{h}
+
\nabla_3\cdot \overline{\mathbf{F}}_h
=
-2\,\overline{\mathbf{E}\cdot\mathbf{B}}.
\label{eq:proj_helicity_local_app}
\end{equation}
Integrating over a brane region \(\Omega\) gives
\begin{equation}
\frac{\dd}{\dd t}\int_{\Omega}\dd^3x\;\overline{h}
=
-2\int_{\Omega}\dd^3x\;\overline{\mathbf{E}\cdot\mathbf{B}}
-\int_{\partial\Omega}\dd S\; \mathbf{n}\cdot \overline{\mathbf{F}}_h.
\label{eq:proj_helicity_integral_app}
\end{equation}

\paragraph{Resolved helicity and subscale (transverse) helicity.}
A brane-only description typically uses projected fields
\(\overline{\mathbf{A}},\overline{\mathbf{B}},\overline{\Phi},\overline{\mathbf{E}}\)
and defines the \emph{resolved} helicity density
\begin{equation}
h_{\rm res}\equiv \overline{\mathbf{A}}\cdot\overline{\mathbf{B}},
\qquad
\mathbf{F}_{h,{\rm res}} \equiv
\overline{\Phi}\,\overline{\mathbf{B}}
+
\overline{\mathbf{E}}\times\overline{\mathbf{A}}.
\label{eq:resolved_helicity_defs_app}
\end{equation}
Because \(\overline{\mathbf{B}}=\nabla_3\times\overline{\mathbf{A}}\) and
\(\overline{\mathbf{E}}=-\nabla_3\overline{\Phi}-\partial_t\overline{\mathbf{A}}\),
the resolved fields satisfy the same identity:
\begin{equation}
\partial_t h_{\rm res}
+
\nabla_3\cdot \mathbf{F}_{h,{\rm res}}
=
-2\,\overline{\mathbf{E}}\cdot\overline{\mathbf{B}}.
\label{eq:resolved_helicity_budget_app}
\end{equation}

However, the projected helicity \(\overline{h}\) is not equal to \(h_{\rm res}\) when
\(w\)-structure is present. Using \eqref{eq:proj_product_app}, define the \emph{subscale}
(transverse-structure) helicity density
\begin{equation}
h_{\rm sub}
\equiv
\overline{h}-h_{\rm res}
=
\overline{\mathbf{A}'\cdot\mathbf{B}'},
\label{eq:subscale_helicity_def_app}
\end{equation}
and the corresponding subscale helicity flux
\begin{equation}
\mathbf{F}_{h,{\rm sub}}
\equiv
\overline{\mathbf{F}}_h - \mathbf{F}_{h,{\rm res}}
=
\overline{\Phi\,\mathbf{B}+\mathbf{E}\times\mathbf{A}}
-
\overline{\Phi}\,\overline{\mathbf{B}}
-
\overline{\mathbf{E}}\times\overline{\mathbf{A}}.
\label{eq:subscale_helicity_flux_app}
\end{equation}
Subtracting \eqref{eq:resolved_helicity_budget_app} from \eqref{eq:proj_helicity_local_app} yields
the exact subscale helicity transfer equation:
\begin{equation}
\partial_t h_{\rm sub}
+
\nabla_3\cdot \mathbf{F}_{h,{\rm sub}}
=
-2\Big(
\overline{\mathbf{E}\cdot\mathbf{B}}
-
\overline{\mathbf{E}}\cdot\overline{\mathbf{B}}
\Big)
=
-2\,\overline{\mathbf{E}'\cdot\mathbf{B}'}.
\label{eq:subscale_helicity_budget_app}
\end{equation}

\paragraph{Interpretation as ``leakage'' into transverse modes.}
Equation \eqref{eq:subscale_helicity_budget_app} isolates a mathematically precise mechanism by
which brane-facing (resolved) helicity can change without invoking collisional dissipation:
helicity can be transferred into transverse structure (nonzero \(h_{\rm sub}\)) through the
covariance source \(\overline{\mathbf{E}'\cdot\mathbf{B}'}\). In the Hermite/KK formulation,
this covariance is exactly the contribution of excited higher transverse modes (and mixed-sector
structure) to helicity budgets.

\paragraph{Connection to reconnection diagnostics.}
In ideal MHD one often expects \(\mathbf{E}\cdot\mathbf{B}=0\) (or equivalently conservation of
magnetic helicity up to boundary fluxes). In the present framework, even if the \emph{resolved}
fields satisfy \(\overline{\mathbf{E}}\cdot\overline{\mathbf{B}}\approx 0\), the projected
quadratic \(\overline{\mathbf{E}\cdot\mathbf{B}}\) can be nonzero due to transverse structure,
leading to genuine evolution of \(\overline{h}\) as seen by a finite-resolution brane
measurement protocol. This provides a conservative, fully diagnostic pathway for topology change:
the helicity ledger closes when one includes the subscale reservoir \(h_{\rm sub}\) and the
explicit \(w\)-exchange channels in the energy ledger (Section~\ref{sec:full_system} and
Appendix~\ref{app:maxwell_reduction}).

\subsection{Summary: what must be recorded in simulations}
\label{app:budgets_summary}

For verification and for interpreting beyond-MHD behavior, simulations should record:
\begin{itemize}[leftmargin=1.7em]
\item the projected EM energy budget \eqref{eq:proj_energy_local_app} including the explicit
\(w\)-leakage term \(\mathcal{L}_{\rm EM}^{(w)}\) in \eqref{eq:energy_leak_term_app};
\item the partition of projected quadratic energies into resolved \(+\) covariance parts as in
\eqref{eq:energy_reynolds_split_app} (mode energies in the \(\text{3D}\times\text{modes}\) scheme);
\item the projected helicity budget \eqref{eq:proj_helicity_local_app} and its resolved/subscale
decomposition \eqref{eq:subscale_helicity_budget_app}, which directly measures transfer of
topological content into transverse structure.
\end{itemize}
These ledgers are the bookkeeping backbone needed to distinguish true dissipation from conservative
export into extra-dimensional degrees of freedom.

\section{Derivation Details: Projected Generalized Ohm's Law and the Topology EMF}
\label{app:ohm_projection}

This appendix provides a short derivation of the projected generalized Ohm law used in
Section~\ref{sec:generalized_ohm_4d}. The logic is:

\begin{enumerate}[leftmargin=1.6em]
\item write the bulk (pointwise in \(w\)) electron momentum equation in \(4+1\)D,
\item solve it for the brane electric field \(\mathbf{E}(t,\mathbf{x},w)\),
\item apply the operational projection \(\overline{(\cdot)}\),
\item add and subtract \(\overline{\mathbf{v}}\times\overline{\mathbf{B}}\) to expose standard
two-fluid/MHD terms plus the new \(4\)D/projection corrections.
\end{enumerate}

\subsection{Bulk two-fluid momentum equation (electron) in \texorpdfstring{\(4+1\)}{4+1}D}
\label{app:ohm_bulk}

Let \(A,B\in\{x,y,z,w\}\) denote bulk spatial indices and \(a,b,c\in\{x,y,z\}\) denote
brane spatial indices. For species \(s\), the bulk number density and bulk velocity are
\(n_s(t,\mathbf{x},w)\) and \(v_s^A(t,\mathbf{x},w)\). The bulk electromagnetic tensor is
\(F_{MN}\), with
\begin{equation}
E_A \equiv F_{A0},\qquad
B_a \equiv \frac12\epsilon_{abc}F_{bc},\qquad
C_a \equiv F_{aw}=\partial_a A_w-\partial_w A_a.
\label{eq:app_EBC_defs}
\end{equation}
Write the full \(4\)D convective derivative (for electrons) as
\begin{equation}
D_t^{(4)} \equiv \partial_t + v_e^a\partial_a + v_e^w\partial_w.
\label{eq:app_Dt4}
\end{equation}

A generic (nonrelativistic) bulk electron momentum balance has the form
\begin{equation}
m_e n_e\,D_t^{(4)} v_{e,A}
=
q_e n_e\Big(E_A + F_{AB}v_e^B\Big)
-\partial_A p_e
+R_{e,A}
+\mathcal{F}^{(\mathrm{extra})}_{e,A},
\qquad q_e=-e,
\label{eq:app_euler_e_bulk}
\end{equation}
where \(p_e\) is the electron pressure (or the equivalent enthalpy-gradient term in the GNLS
realization), \(R_{e,A}\) is electron--ion momentum exchange (optional), and
\(\mathcal{F}^{(\mathrm{extra})}_{e,A}\) collects any additional modeled forces (e.g.\ dispersive
terms if retained).

\subsection{Mixed-sector decomposition of the Lorentz force}
\label{app:ohm_mixed_decomp}

Restrict \eqref{eq:app_euler_e_bulk} to brane components \(A=a\). Using antisymmetry of \(F_{AB}\),
\begin{equation}
F_{ab}v_e^b = (\mathbf{v}_e\times\mathbf{B})_a,
\qquad
F_{aw}v_e^w = v_e^w C_a,
\label{eq:app_F_split}
\end{equation}
so the brane-component Lorentz term becomes
\begin{equation}
E_a + F_{ab}v_e^b + F_{aw}v_e^w
=
E_a + (\mathbf{v}_e\times\mathbf{B})_a + v_e^w C_a.
\label{eq:app_lorentz_brane_split}
\end{equation}
This identifies the key beyond-\(3+1\) channel: \(v_e^w C_a\), which vanishes if either
\(v_e^w=0\) (no transverse drift) or \(C_a=F_{aw}=0\) (no mixed \(aw\) structure).

\subsection{Pointwise (in \texorpdfstring{\(w\)}{w}) bulk generalized Ohm law}
\label{app:ohm_pointwise}

Solve \eqref{eq:app_euler_e_bulk} for the brane electric field \(\mathbf{E}=(E_x,E_y,E_z)\) using
\eqref{eq:app_lorentz_brane_split}. In vector notation (all quantities still depend on \(w\)):
\begin{equation}
\mathbf{E}
=
-\mathbf{v}_e\times\mathbf{B}
-\;v_e^w\,\mathbf{C}
-\frac{1}{e n_e}\nabla_3 p_e
+\frac{1}{e n_e}\mathbf{R}_e
-\frac{m_e}{e}\,D_t^{(4)}\mathbf{v}_e
+\frac{1}{e n_e}\boldsymbol{\mathcal{F}}^{(\mathrm{extra})}_e,
\label{eq:app_ohm_pointwise}
\end{equation}
where \(\nabla_3\) differentiates only brane coordinates \(\mathbf{x}=(x,y,z)\) and
\(\mathbf{C}=(C_x,C_y,C_z)\).

\subsection{Operational projection and exact brane-facing identity}
\label{app:ohm_project_exact}

Recall the operational brane projection
\begin{equation}
\overline{Q}(t,\mathbf{x}) \equiv \int_{-\infty}^{+\infty}\dd w\,W(w)\,Q(t,\mathbf{x},w),
\qquad
\int \dd w\,W(w)=1.
\label{eq:app_proj_def_ohm}
\end{equation}
Project \eqref{eq:app_ohm_pointwise} to obtain the \emph{exact} identity
\begin{equation}
\overline{\mathbf{E}}
=
-\overline{\mathbf{v}_e\times\mathbf{B}}
-\overline{v_e^w\,\mathbf{C}}
-\overline{\frac{1}{e n_e}\nabla_3 p_e}
+\overline{\frac{1}{e n_e}\mathbf{R}_e}
-\frac{m_e}{e}\,\overline{D_t^{(4)}\mathbf{v}_e}
+\overline{\frac{1}{e n_e}\boldsymbol{\mathcal{F}}^{(\mathrm{extra})}_e}.
\label{eq:app_ohm_projected_exact}
\end{equation}
At this stage, no MHD closure has been assumed; the only operation performed is the
measurement/projection map.

\subsection{Add--subtract \texorpdfstring{$\overline{\mathbf{v}}\times\overline{\mathbf{B}}$}{vbar x Bbar} and expose standard terms}
\label{app:ohm_add_subtract}

Introduce a brane single-fluid velocity \(\overline{\mathbf{v}}\) and brane current
\(\overline{\mathbf{J}}\) in the usual quasi-neutral ordering (Appendix~\ref{app:twofluid_to_mhd}).
Add and subtract \(\overline{\mathbf{v}}\times\overline{\mathbf{B}}\) to \eqref{eq:app_ohm_projected_exact}:
\begin{equation}
\overline{\mathbf{E}}+\overline{\mathbf{v}}\times\overline{\mathbf{B}}
=
\underbrace{\Big(\overline{\mathbf{v}}-\overline{\mathbf{v}_e}\Big)\times\overline{\mathbf{B}}}_{\text{Hall rewrite}}
\;+\;
\underbrace{\Big(\overline{\mathbf{v}_e}\times\overline{\mathbf{B}}-\overline{\mathbf{v}_e\times\mathbf{B}}\Big)}_{\text{projection covariance}}
\;+\;
\Big[\text{all remaining projected terms from \eqref{eq:app_ohm_projected_exact}}\Big].
\label{eq:app_ohm_split_step}
\end{equation}

\paragraph{Hall term (standard).}
Under quasi-neutrality with \(\overline{n}\equiv \overline{n_e}\simeq \overline{n_i}\) and the
standard ordering \(m_e\ll m_i\), one uses
\begin{equation}
\overline{\mathbf{v}}-\overline{\mathbf{v}_e}
\;\approx\;
\frac{\overline{\mathbf{J}}}{e\,\overline{n}},
\label{eq:app_hall_approx}
\end{equation}
which gives the usual Hall contribution
\begin{equation}
\Big(\overline{\mathbf{v}}-\overline{\mathbf{v}_e}\Big)\times\overline{\mathbf{B}}
\;\approx\;
\frac{\overline{\mathbf{J}}\times\overline{\mathbf{B}}}{e\,\overline{n}}.
\label{eq:app_hall_term}
\end{equation}

\paragraph{Covariance EMF (new, exact).}
Define the cross-covariance EMF
\begin{equation}
\boldsymbol{\mathcal{E}}_{\rm cov}
\equiv
\overline{\mathbf{v}_e}\times\overline{\mathbf{B}}-\overline{\mathbf{v}_e\times\mathbf{B}}
=
-\Big(\overline{\mathbf{v}_e\times\mathbf{B}}-\overline{\mathbf{v}_e}\times\overline{\mathbf{B}}\Big).
\label{eq:app_Ecov_def}
\end{equation}
This term vanishes if the fields have negligible \(w\)-dependence (controlled reduction), but is
otherwise generically nonzero even in laminar flows.

\subsection{Mixed-sector term and projection residuals}
\label{app:ohm_mixed_and_residuals}

The remaining terms in \eqref{eq:app_ohm_projected_exact} are grouped as follows.

\paragraph{Mixed-sector transverse EMF (new).}
The \(4+1\)D mixed coupling produces the brane-facing contribution
\begin{equation}
\boldsymbol{\mathcal{E}}_{w}
\equiv
-\overline{v_e^w\,\mathbf{C}},
\qquad \mathbf{C}=(F_{xw},F_{yw},F_{zw}).
\label{eq:app_Ew_def}
\end{equation}
This is a direct, computable pathway for brane non-ideality sourced by transverse drift and mixed fields.

\paragraph{Projection residuals (exact bookkeeping).}
Projection does not commute with ratios or nonlinear closures. For compactness, define
\begin{align}
\boldsymbol{\mathcal{E}}_{p}
&\equiv
-\overline{\frac{1}{e n_e}\nabla_3 p_e}
+\frac{1}{e\,\overline{n}}\nabla_3\overline{p_e},
\label{eq:app_Ep_def}
\\
\boldsymbol{\mathcal{E}}_{R}
&\equiv
\overline{\frac{1}{e n_e}\mathbf{R}_e}-\eta\,\overline{\mathbf{J}},
\label{eq:app_ER_def}
\\
\boldsymbol{\mathcal{E}}_{I}
&\equiv
-\frac{m_e}{e}\Big(\overline{D_t^{(4)}\mathbf{v}_e}-\mathcal{I}_e\Big),
\label{eq:app_EI_def}
\end{align}
where \(\eta\,\overline{\mathbf{J}}\) and \(\mathcal{I}_e\) are the \emph{chosen} brane closures for resistivity
and electron inertia (if retained). In ideal-MHD-like limits these are often dropped; in Hall/extended MHD
they are kept in standard forms (Appendix~\ref{app:twofluid_to_mhd}).

\subsection{Final projected generalized Ohm law}
\label{app:ohm_final}

Collecting \eqref{eq:app_hall_term}--\eqref{eq:app_EI_def} yields the brane-facing generalized Ohm law:
\begin{equation}
\boxed{
\overline{\mathbf{E}}+\overline{\mathbf{v}}\times\overline{\mathbf{B}}
=
\frac{\overline{\mathbf{J}}\times\overline{\mathbf{B}}}{e\,\overline{n}}
-\frac{1}{e\,\overline{n}}\nabla_3\overline{p_e}
+\eta\,\overline{\mathbf{J}}
-\frac{m_e}{e}\mathcal{I}_e
+\boldsymbol{\mathcal{E}}_{\rm topo}
+\overline{\frac{1}{e n_e}\boldsymbol{\mathcal{F}}^{(\mathrm{extra})}_e},
}
\label{eq:app_ohm_final}
\end{equation}
where the topology EMF is the sum of explicitly \(4\)D/projection-driven contributions:
\begin{equation}
\boxed{
\boldsymbol{\mathcal{E}}_{\rm topo}
\equiv
\boldsymbol{\mathcal{E}}_{\rm cov}
+\boldsymbol{\mathcal{E}}_{w}
+\boldsymbol{\mathcal{E}}_{p}
+\boldsymbol{\mathcal{E}}_{R}
+\boldsymbol{\mathcal{E}}_{I}.
}
\label{eq:app_Etopo_def}
\end{equation}
Equations \eqref{eq:app_ohm_final}--\eqref{eq:app_Etopo_def} are equivalent to the statement in
Section~\ref{sec:generalized_ohm_4d}, but with the intermediate bookkeeping made explicit.

\paragraph{Use in the induction equation.}
Projected Faraday gives \(\partial_t\overline{\mathbf{B}}=-\nabla_3\times\overline{\mathbf{E}}\), so any component of
\(\boldsymbol{\mathcal{E}}_{\rm topo}\) with nonzero curl acts as a brane-facing non-ideal source. In particular,
\(\boldsymbol{\mathcal{E}}_{w}=-\overline{v_e^w\mathbf{C}}\) provides a conservative, mixed-sector pathway for
reconnection-like behavior without inserting ad hoc resistivity, while the covariance term \(\boldsymbol{\mathcal{E}}_{\rm cov}\)
quantifies the direct impact of transverse structure (mode tower / unresolved \(w\)-variation) on brane topology.

\section{Geometry DOFs and Coupling to Plasma Loads}
\label{app:geometry_dofs_plasma}

This appendix describes how to endow the model with a small set of \emph{collective}
geometry degrees of freedom (DOFs) and how those DOFs are driven by \emph{plasma loads}
(currents, fields, and pressure/enthalpy). The intent is computational: rather than
solving a full geometric PDE, we treat the transverse localization and throat geometry
as a low-dimensional subsystem coupled (in a ledger-consistent way) to the bulk plasma--EM
fields. This provides a principled mechanism for (i) time-dependent localization width,
(ii) dynamic tube/throat breathing, and (iii) conservative export of energy/topology into
transverse/geometry channels.

\subsection{Collective geometry variables}
\label{app:geom_vars}

We introduce a vector of geometry variables \(\Theta^I(t)\). The minimal set used in the
\(4+1\) model papers is
\begin{equation}
\Theta^I(t) \in \{a(t),\,L(t)\},
\end{equation}
where \(a(t)\) is an effective transverse radius (in the \((x,y,z)\) directions) and
\(L(t)\) is an effective extent in the transverse direction \(w\).

For plasma applications, a second (often more directly useful) parameterization is to
treat the electromagnetic localization width \(\lambda(t)\) as a geometry DOF through
\begin{equation}
Z(w;\lambda) = \exp\!\left(-\frac{w^2}{\lambda^2}\right),
\qquad \Theta^I(t)\in\{\lambda(t)\},
\end{equation}
or, more generally, to allow \(Z\) to depend on \(\Theta^I\):
\begin{equation}
Z(w) \;\longrightarrow\; Z\bigl(w;\Theta(t)\bigr).
\end{equation}
In the \(\text{3D}\times\text{modes}\) formulation (Appendix~\ref{sec:kk_modes}), \(\lambda\)
controls the KK/Hermite spectrum \(m_n^2\propto 1/\lambda^2\), so a time-dependent \(\lambda(t)\)
is a direct, low-dimensional knob for transverse-mode accessibility.

\paragraph{Operational note.}
We treat the projection kernel \(W(w)\) as part of the \emph{measurement definition} and keep it
fixed, even when \(Z(w;\Theta)\) changes. (One can also parametrize \(W\), but that changes what the
brane observer means, rather than changing the underlying interaction geometry.)

\subsection{Where geometry enters the coupled system}
\label{app:geom_enters_action}

At the level of the action/energy ledger, geometry variables enter only through explicit
parameter dependence of (i) the EM localization profile \(Z(w;\Theta)\), (ii) any confinement or
core-regularization potential \(V_{\rm conf}(\mathbf{X};\Theta)\) in the matter sector, and (iii)
a standalone geometry energy \(E_{\rm geom}(\Theta)\) representing wall tension, volume work, etc.
Schematically:
\begin{equation}
S[A,\{\text{matter}\},\Theta]
=
S_{\rm EM}[A;Z(\cdot;\Theta)]
+
S_{\rm mat}[\{\text{matter}\};V_{\rm conf}(\cdot;\Theta)]
+
S_{\rm geom}[\Theta].
\end{equation}

A referee-clean closure for the geometry subsystem is an effective mechanical Lagrangian
\begin{equation}
\boxed{
\mathcal{L}_{\rm geom}
=
\frac12\,M_{IJ}\,\dot\Theta^I\dot\Theta^J
-
E_{\rm geom}(\Theta)
-
E_{\rm ratio}(\Theta)
},
\label{eq:app_Lgeom_general}
\end{equation}
with positive-definite inertia matrix \(M_{IJ}\). The optional penalty \(E_{\rm ratio}\) enforces
a preferred manifold (e.g.\ \(L\simeq \alpha a\)) when desired:
\begin{equation}
E_{\rm ratio}(a,L)=\frac{\kappa}{2}\,\bigl(L-\alpha a\bigr)^2,
\qquad (\kappa\ge 0).
\label{eq:app_Eratio}
\end{equation}

\subsection{Generalized forces from the total Hamiltonian}
\label{app:geom_forces}

Define the conservative total energy
\begin{equation}
H_{\rm tot}
=
H_{\rm EM}[A;\Theta]
+
H_{\rm mat}[\{\text{matter}\};\Theta]
+
E_{\rm geom}(\Theta)
+
E_{\rm ratio}(\Theta),
\label{eq:app_Htot_geom}
\end{equation}
where \(H_{\rm EM}\) is the bulk electromagnetic Hamiltonian (including \(Z(w;\Theta)\)),
and \(H_{\rm mat}\) is the bulk plasma Hamiltonian appropriate to the chosen matter model
(two-fluid, GNLS, Vlasov/PIC, etc.).

The \emph{ledger definition} of the generalized force on \(\Theta^I\) is
\begin{equation}
\boxed{
F_I \;\equiv\; -\frac{\partial H_{\rm tot}}{\partial \Theta^I},
}
\label{eq:app_generalized_force_def}
\end{equation}
where the derivative is taken at fixed instantaneous fields (Hellmann--Feynman / virtual work).

With \eqref{eq:app_Htot_geom}, the force splits into contributions:
\begin{equation}
F_I = F_I^{({\rm EM})}+F_I^{({\rm mat})}+F_I^{({\rm geom})}+F_I^{({\rm ratio})}.
\end{equation}

\paragraph{(i) Geometry force from localized Maxwell.}
Using the localized Maxwell energy density (Appendix~\ref{app:proj_energy_helicity}),
\begin{equation}
H_{\rm EM}
=
\int \dd^3x\,\dd w\;
\frac{Z(w;\Theta)}{2\mu_0}
\left(
E_AE_A+\frac12 F_{AB}F_{AB}
\right),
\qquad A,B\in\{x,y,z,w\},
\end{equation}
the explicit parametric force contribution is
\begin{equation}
\boxed{
F_I^{({\rm EM})}
=
-\frac{\partial H_{\rm EM}}{\partial \Theta^I}
=
-\int \dd^3x\,\dd w\;
\frac{\partial_{\Theta^I}Z(w;\Theta)}{2\mu_0}
\left(
E_AE_A+\frac12 F_{AB}F_{AB}
\right)
\;+\;F_{I,\rm extra}^{({\rm EM})}.
}
\label{eq:app_FI_EM}
\end{equation}
The term \(F_{I,\rm extra}^{({\rm EM})}\) is present only if additional EM coefficients
(or external-source profiles) are made geometry-dependent. If \(Z\) is the only explicit
geometry dependence in the EM sector, then \(F_{I,\rm extra}^{({\rm EM})}=0\).

\paragraph{(ii) Geometry force from plasma/matter loading.}
The matter-sector contribution depends on the matter model, but the geometry dependence is
typically through a confinement/core potential \(V_{\rm conf}(\mathbf{X};\Theta)\).
Two common cases:

\emph{GNLS / field-theory matter:} if the Hamiltonian contains \(\sum_s \int \rho_s V_{\rm conf}\),
then
\begin{equation}
\boxed{
F_I^{({\rm mat})}
=
-\frac{\partial H_{\rm mat}}{\partial \Theta^I}
=
-\sum_s \int \dd^3x\,\dd w\;
\rho_s(t,\mathbf{x},w)\,\partial_{\Theta^I}V_{\rm conf}(\mathbf{x},w;\Theta)
\;+\;\cdots
}
\label{eq:app_FI_mat_GNLS}
\end{equation}
where \(\cdots\) denotes any additional explicit geometry dependence of EOS or dispersive terms
(if present).

\emph{Vlasov / PIC matter:} if particles experience a geometry-dependent potential \(V_{\rm conf}\),
then in continuum form
\begin{equation}
H_{\rm mat}\supset \sum_s \int f_s(t,\mathbf{x},w,\mathbf{v},v_w)\,V_{\rm conf}(\mathbf{x},w;\Theta)\,
\dd^3x\,\dd w\,\dd^3v\,\dd v_w,
\end{equation}
and hence
\begin{equation}
\boxed{
F_I^{({\rm mat})}
=
-\sum_s \int f_s\;\partial_{\Theta^I}V_{\rm conf}\;
\dd^3x\,\dd w\,\dd^3v\,\dd v_w
=
-\sum_s \int \rho_s(t,\mathbf{x},w)\,\partial_{\Theta^I}V_{\rm conf}(\mathbf{x},w;\Theta)\;
\dd^3x\,\dd w,
}
\label{eq:app_FI_mat_PIC}
\end{equation}
which is the same loading form, written either in phase space or using \(\rho_s=\int f_s\,\dd^3v\,\dd v_w\).

\paragraph{(iii) Ratio-penalty forces.}
For \eqref{eq:app_Eratio}, the explicit forces are
\begin{equation}
\boxed{
F_a^{({\rm ratio})} = -\frac{\partial E_{\rm ratio}}{\partial a}
= +\kappa\alpha\,(L-\alpha a),
\qquad
F_L^{({\rm ratio})} = -\frac{\partial E_{\rm ratio}}{\partial L}
= -\kappa\,(L-\alpha a).
}
\label{eq:app_ratio_forces}
\end{equation}

\subsection{Dynamics: damped collective-coordinate evolution and quasi-static limits}
\label{app:geom_dynamics}

A baseline (simulation-friendly) geometry evolution law is a damped second-order system:
\begin{equation}
\boxed{
M_{IJ}\,\ddot\Theta^J + \Gamma_{IJ}\,\dot\Theta^J = F_I,
}
\label{eq:app_geom_eom_general}
\end{equation}
with positive semidefinite damping matrix \(\Gamma_{IJ}\). In diagonal form for \(\Theta^I=(a,L)\),
this becomes
\begin{equation}
M_a \ddot a + \Gamma_a \dot a = F_a,
\qquad
M_L \ddot L + \Gamma_L \dot L = F_L.
\end{equation}

Two controlled limits are frequently useful:
\begin{itemize}[leftmargin=1.7em]
\item \textbf{Overdamped relaxation:} \(M_{IJ}\to 0\), yielding \(\Gamma_{IJ}\dot\Theta^J = F_I\).
\item \textbf{Adiabatic elimination / instantaneous minimization:} when geometry relaxation is much faster
than plasma timescales, one sets
\begin{equation}
\frac{\partial H_{\rm tot}}{\partial \Theta^I}\approx 0,
\qquad\text{i.e.}\qquad F_I\approx 0,
\end{equation}
and solves an algebraic equilibrium problem for \(\Theta(t)\) given instantaneous plasma loads.
\end{itemize}

\paragraph{Optional: geometry fields on the brane.}
If one wants \emph{spatially varying} localization/throat response, promote
\(\Theta^I(t)\to \Theta^I(t,\mathbf{x})\) and use a gradient-regularized relaxation PDE,
e.g.
\begin{equation}
\tau_I\,\partial_t \Theta^I
=
-\frac{\delta H_{\rm tot}}{\delta \Theta^I}
+ D_I\,\nabla_3^2 \Theta^I,
\label{eq:app_geom_field_relax}
\end{equation}
where \(\delta/\delta\Theta^I\) is a functional derivative and the diffusion term enforces
smoothness (prevents grid-scale localization-width noise). This is optional but can be valuable
for modeling localized thinning/thickening that correlates with reconnection sites.

\subsection{Worked example: dynamic Gaussian width \texorpdfstring{$\lambda(t)$}{lambda(t)}}
\label{app:geom_lambda_example}

Take
\begin{equation}
Z(w;\lambda)=\exp\!\left(-\frac{w^2}{\lambda^2}\right),
\qquad
\partial_\lambda Z(w;\lambda)=\frac{2w^2}{\lambda^3}\,Z(w;\lambda).
\label{eq:app_dZ_dlambda}
\end{equation}
Then the EM-driven generalized force from \eqref{eq:app_FI_EM} becomes
\begin{equation}
\boxed{
F_\lambda^{({\rm EM})}
=
-\int \dd^3x\,\dd w\;
\frac{1}{2\mu_0}
\left(\frac{2w^2}{\lambda^3}Z(w;\lambda)\right)
\left(
E_AE_A+\frac12 F_{AB}F_{AB}
\right).
}
\label{eq:app_Flambda_EM}
\end{equation}
This force is always nonpositive for nontrivial fields (it tends to reduce \(\lambda\))
unless balanced by \(E_{\rm geom}(\lambda)\) or matter contributions that favor expansion.
A stable \(\lambda_\star\) requires competing effects in \(H_{\rm tot}\).

\paragraph{Impact on reduced brane parameters.}
When \(\lambda\) changes, both the effective Maxwell coupling and KK masses change:
\begin{equation}
\mu_0^{\rm eff}(\lambda) = \frac{\mu_0}{\int \dd w\,Z(w;\lambda)}
= \frac{\mu_0}{\lambda\sqrt{\pi}},
\qquad
m_n^2(\lambda)=\frac{2n}{\lambda^2}.
\end{equation}
Equivalently, in a canonically normalized brane packaging one may write
\(q_{\rm eff}(\lambda)=q_*/\sqrt{\Zint(\lambda)}\) for a fixed microscopic branch
charge \(q_*=\eta_Q e_*\). Thus changing \(\lambda\) repackages the reduced brane
coupling and the KK spectrum, but it does not make the underlying branch sign or
the fixed microscopic charge \(q_*\) itself vary.
In a \(\text{3D}\times\text{modes}\) code, this means the diagonal mode operator
\((\square_4+m_n^2)\) acquires a time-dependent mass term. If \(\lambda(t)\) varies slowly
relative to the EM wave timestep, this is straightforward to implement as a time-dependent
coefficient.

\subsection{Worked example: tube DOFs \texorpdfstring{$(a(t),L(t))$}{(a(t),L(t))} and a minimal geometry energy}
\label{app:geom_aL_example}

To model breathing/ringing of a localized tube/throat, one may use \(\Theta=(a,L)\) with a minimal
geometry energy (volume work plus surface tension):
\begin{equation}
E_{\rm geom}(a,L) = P_{\rm vac}\,V(a,L) + \sigma\,A(a,L),
\label{eq:app_Egeom_aL}
\end{equation}
where a natural \(4\)D tube measure for a corridor \(B^3(a)\times[0,L]\subset\mathbb{R}^4\) is
\begin{equation}
V(a,L)=\frac{4\pi}{3}a^3\,L,
\qquad
A(a,L)=4\pi a^2\,L + 2\cdot\frac{4\pi}{3}a^3.
\label{eq:app_VA_aL}
\end{equation}
The resulting purely geometric restoring forces are
\begin{equation}
F_a^{({\rm geom})}=-\frac{\partial E_{\rm geom}}{\partial a},
\qquad
F_L^{({\rm geom})}=-\frac{\partial E_{\rm geom}}{\partial L},
\end{equation}
with explicit derivatives obtained directly from \eqref{eq:app_Egeom_aL}--\eqref{eq:app_VA_aL}.

The plasma loads enter through \(F_a^{({\rm mat})}\) and \(F_L^{({\rm mat})}\) via
\(\partial_{a,L}V_{\rm conf}\), and (optionally) through \(F_{a,L}^{({\rm EM})}\) if the EM sector is
given explicit geometry dependence (e.g.\ geometry-dependent \(Z\), geometry-dependent imposed drive).

\subsection{How this is used in massive plasma simulations}
\label{app:geom_in_hpc}

In large runs, geometry DOFs are intended to be cheap and ledger-consistent. A practical operator-split
update per timestep \(\Delta t\) is:

\begin{enumerate}[leftmargin=1.7em]
\item \textbf{Advance plasma and EM fields} from \(t^n\to t^{n+1}\) using the current geometry \(\Theta^n\)
(i.e.\ current \(Z(w;\Theta^n)\), \(V_{\rm conf}(\cdot;\Theta^n)\), and mode masses if applicable).

\item \textbf{Compute instantaneous geometry forces} \(F_I^n\) using \eqref{eq:app_generalized_force_def} with
numerical quadrature in \(w\). In a \(\text{3D}\times\text{modes}\) solver, the EM force contribution can be
computed from mode energies because \(H_{\rm EM}\) is a sum of diagonal mode contributions plus any explicit
\(\Theta\)-dependence.

\item \textbf{Update geometry} \(\Theta^n\to \Theta^{n+1}\) using \eqref{eq:app_geom_eom_general} with a stable ODE
integrator (explicit with damping is often sufficient; implicit is optional if stiffness is large).

\item \textbf{Update derived coefficients} (e.g.\ \(\mu_0^{\rm eff}(\Theta)\), \(m_n(\Theta)\), coupling weights \(c_n(\Theta)\))
for the next field step.
\end{enumerate}

\paragraph{Energy accounting.}
If \(\Gamma_{IJ}=0\) and the time integration is done conservatively, the combined field+geometry energy
\eqref{eq:app_Htot_geom} is conserved up to numerical error. With \(\Gamma_{IJ}\neq 0\), geometry damping removes
energy from the resolved system; this can be interpreted as coupling to unresolved internal wall modes and should
be tracked explicitly in the simulation energy ledger.

\paragraph{Why this matters for beyond-MHD physics.}
Allowing \(\Theta\) (and hence \(Z\), confinement, and transverse spectra) to respond to plasma loads provides a
physically grounded mechanism for apparent non-ideality on the brane: magnetic energy and helicity can be exported
into transverse and geometry reservoirs (Appendix~\ref{app:proj_energy_helicity}) rather than being represented
solely by an \emph{ad hoc} resistivity. This is precisely the kind of channel needed to model reconnection-like
phenomena while keeping the bulk dynamics conservative and diagnosable.

\end{document}
